% La Philosophie africaine — Philosophie 1re Tech — Cours signé OnBuch+
% Programme officiel MINESEC. Compiler avec : tectonic N01-philosophie-africaine.tex
\newcommand{\DOCMATIERE}{Philosophie}
\newcommand{\DOCNIVEAU}{1re Tech}
\newcommand{\DOCMODULE}{Philosophie – classe de Première technique}
\newcommand{\DOCLECON}{N01}
\newcommand{\DOCTITRE}{La Philosophie africaine}
% OnBuch+ — préambule « cours signés » ENRICHI (compilé avec tectonic / XeLaTeX)
% Système visuel « Pop » d'OnBuch+ : fond crème (jamais blanc), bordures encre
% épaisses, ombres « dures » décalées, titres Archivo Black en capitales.
% Version MATHÉMATIQUES : graphes (pgfplots/tikz), boîtes pédagogiques
% (définition, propriété, méthode, attention, le savais-tu, exercice résolu,
% exercice, TP, bilan), page de garde et signature OnBuch+ ancrée en bas.
%
% Macros attendues dans main.tex AVANT \input{preamble} :
%   \DOCMATIERE  \DOCNIVEAU  \DOCMODULE  \DOCLECON (numéro)  \DOCTITRE
\documentclass[11pt]{article}

\usepackage{fontspec}
\usepackage[french]{babel}
\usepackage[a4paper,margin=2cm,top=3.4cm,bottom=2.8cm,headheight=1.4cm,footskip=1.3cm]{geometry}
\usepackage{amsmath,amssymb,amsfonts}
\usepackage{mathtools} % \xmapsto, \coloneqq... souvent utilisés par les modèles
\usepackage{cancel} % simplifications barrées (bilans de forces)
% \abs{z} (module, valeur absolue) et \norm{u} (norme), utilisés par les modèles
\providecommand{\abs}{}\let\abs\relax\DeclarePairedDelimiter{\abs}{\lvert}{\rvert}
\providecommand{\norm}{}\let\norm\relax\DeclarePairedDelimiter{\norm}{\lVert}{\rVert}
\usepackage[table]{xcolor} % [table] : fournit aussi \rowcolors (alternance de lignes)
\usepackage{pagecolor}
\usepackage{tikz}
\usetikzlibrary{arrows.meta,positioning,calc,shapes.geometric,shapes.misc,shapes.arrows,decorations.pathmorphing,decorations.pathreplacing,decorations.markings,patterns,shadows,mindmap,trees,fit,backgrounds,matrix,intersections,angles,quotes}
\usepackage{pgfplots}
\usepackage[version=4]{mhchem} % équations nucléaires \ce{^{235}_{92}U}
\usepackage[european,siunitx]{circuitikz} % schémas électriques
\pgfplotsset{compat=1.18}
\usepgfplotslibrary{fillbetween,groupplots} % aires entre courbes (calcul intégral)
\usepackage{enumitem}
\usepackage{fancyhdr}
% [most] charge toutes les bibliothèques tcolorbox (listings, minted, raster,
% documentation...) et gonfle inutilement la mémoire TeX sur un document long
% (~300 boîtes) : on ne charge que ce qui est réellement utilisé.
\usepackage[skins,breakable]{tcolorbox}
\usepackage{graphicx}
\usepackage{colortbl}
\usepackage{booktabs}
\usepackage{tabularx}
\usepackage{multirow}
\usepackage{diagbox} % case d'en-tête diagonale des tableaux à double entrée
% Colonnes extensibles centrée / gauche / droite (C, L, R), souvent utilisées par les modèles
\newcolumntype{C}{>{\centering\arraybackslash}X}
\newcolumntype{L}{>{\raggedright\arraybackslash}X}
\newcolumntype{R}{>{\raggedleft\arraybackslash}X}
\usepackage{array}
\usepackage{multicol}
\usepackage{siunitx}
% parse-numbers=false : accepte aussi \SI{2,5 \times 10^{-3}}{...}
\sisetup{locale=FR,output-decimal-marker={,},group-minimum-digits=4,parse-numbers=false}
\DeclareSIUnit\ohm{\ensuremath{\symup{\Omega}}} % U+2126 absent de la police
\DeclareSIUnit\division{div} % réglages d'oscilloscope (ms/div, V/div)
\DeclareSIUnit\tr{tr} % tours (vitesse de rotation, ex. \SI{1500}{\tr\per\minute})
\DeclareSIUnit\molar{M} % concentration molaire (ex. \SI{3}{\molar}, \SI{1}{\micro\molar})
\DeclareSIUnit\an{an} % année (le modèle confond parfois avec \year, un compteur TeX) — ex. \SI{5}{\kilo\meter\per\an}
\DeclareSIUnit\FCFA{FCFA} % franc CFA (ex. \SI{500}{\FCFA\per\kilo\gram})
\DeclareSIUnit\degreeBrix{\textdegree Brix} % teneur en sucre d'un sirop/jus (réfractométrie)
\DeclareSIUnit\month{mois} % le modèle utilise parfois \month, un compteur TeX, comme unité de durée
\DeclareSIUnit\microliter{\micro\liter} % le modèle écrit parfois \microliter en un seul token
\DeclareSIUnit\million{millions} % dénombrement (ex. \SI{15}{\million\per\mL} spermatozoïdes)
\usepackage{qrcode}
% \degree : commande fréquemment utilisée par les modèles pour un angle ou une
% température en dehors de \SI{}{} (non fournie par les paquets chargés).
\providecommand{\degree}{^{\circ}}
% \qed (fin de démonstration), utilisé par les modèles sans amsthm
\providecommand{\qed}{\hfill$\square$}
% \barre : double barre des valeurs interdites dans un tableau de variations (tkz-tab)
\providecommand{\barre}{\ensuremath{\|}}
% environnement proof (amsthm non chargé) utilisé par les modèles
\ifdefined\proof\else\newenvironment{proof}[1][Démonstration]{\par\noindent\textit{#1.}\ }{\qed\par}\fi
\usepackage{lastpage}
\usepackage{hyperref}
\hypersetup{hidelinks}

% ============================================================
% Palette « Pop » OnBuch+ — mêmes hex que la classe P (plus_ui.dart)
% ============================================================
\definecolor{poporange}{HTML}{F59321}
\definecolor{popdark}{HTML}{E07A0C}
\definecolor{popcream}{HTML}{FFF6E8}
\definecolor{popink}{HTML}{1C1714}
\definecolor{poppurple}{HTML}{7A5AE0}
\definecolor{popgreen}{HTML}{2E9E6A}
\definecolor{poppink}{HTML}{E0457B}
\definecolor{popblue}{HTML}{2F7DE1}
\definecolor{popgold}{HTML}{C98A12}
\definecolor{popmuted}{HTML}{8A7F76}
\definecolor{popline}{HTML}{EADFD2}
\definecolor{popgray}{HTML}{8A7F76}
\colorlet{popgrey}{popgray}
% Alias tolérants (le modèle nomme parfois les couleurs différemment)
\colorlet{popwhite}{white}
\colorlet{popblack}{popink}
\colorlet{popred}{poppink}
\colorlet{popyellow}{popgold}
% Teintes claires (fonds de tableaux/encadrés) — le modèle nomme parfois
% les couleurs avec un suffixe L (ex. popblueL) sans les avoir définies.
\colorlet{poporangeL}{poporange!20}
\colorlet{popdarkL}{popdark!20}
\colorlet{poppurpleL}{poppurple!20}
\colorlet{popgreenL}{popgreen!20}
\colorlet{poppinkL}{poppink!20}
\colorlet{popblueL}{popblue!20}
\colorlet{popgoldL}{popgold!20}
\colorlet{popredL}{popred!20}
\colorlet{popgrayL}{popgray!20}
% Teintes claires pour fonds de tableaux / zones
\colorlet{popblueL}{popblue!10!white}
\colorlet{poporangeL}{poporange!14!white}
\colorlet{popgreenL}{popgreen!12!white}
\colorlet{poppurpleL}{poppurple!12!white}
\colorlet{poppinkL}{poppink!12!white}
\colorlet{popgoldL}{popgold!14!white}

\pagecolor{popcream}

% ============================================================
% Typographie
% ============================================================
\defaultfontfeatures{Ligatures=TeX}
\setmainfont{PlusJakartaSans-Medium.ttf}[Path=../../../fonts/,BoldFont=PlusJakartaSans-Bold.ttf]
% Maths en sans-serif (Fira Math) avec chiffres et lettres droites de Plus
% Jakarta, pour que les formules se fondent dans le texte.
\usepackage[math-style=ISO,bold-style=ISO]{unicode-math}
\setmathfont{FiraMath-Regular.otf}[Path=../../../fonts/]
\setmathfont{PlusJakartaSans-Medium.ttf}[Path=../../../fonts/,range={up/num,up/{Latin,latin},\mathup}]
\usepackage{icomma} % virgule décimale sans espace en mode math
% Caractères mathématiques Unicode absents des polices de texte (Plus Jakarta,
% Archivo Black) : rendus par leur équivalent mathématique, en texte comme en
% formule — sinon ils s'affichent en carré vide (ex. « Arithmétique dans ℤ »).
\usepackage{newunicodechar}
\newunicodechar{ℝ}{\ensuremath{\mathbb{R}}}
\newunicodechar{ℤ}{\ensuremath{\mathbb{Z}}}
\newunicodechar{ℂ}{\ensuremath{\mathbb{C}}}
\newunicodechar{ℕ}{\ensuremath{\mathbb{N}}}
\newunicodechar{ℚ}{\ensuremath{\mathbb{Q}}}
\newunicodechar{𝔻}{\ensuremath{\mathbb{D}}}
\newunicodechar{α}{\ensuremath{\alpha}}
\newunicodechar{β}{\ensuremath{\beta}}
\newunicodechar{γ}{\ensuremath{\gamma}}
\newunicodechar{δ}{\ensuremath{\delta}}
\newunicodechar{ε}{\ensuremath{\varepsilon}}
\newunicodechar{θ}{\ensuremath{\theta}}
\newunicodechar{λ}{\ensuremath{\lambda}}
\newunicodechar{μ}{\ensuremath{\mu}}
\newunicodechar{σ}{\ensuremath{\sigma}}
\newunicodechar{τ}{\ensuremath{\tau}}
\newunicodechar{φ}{\ensuremath{\varphi}}
\newunicodechar{ψ}{\ensuremath{\psi}}
\newunicodechar{ω}{\ensuremath{\omega}}
\newunicodechar{Σ}{\ensuremath{\Sigma}}
% majuscules produites par \MakeUppercase dans les titres (α-aminés -> Α)
\newunicodechar{Α}{\ensuremath{\alpha}}
\newunicodechar{Β}{\ensuremath{\beta}}
\newunicodechar{Γ}{\ensuremath{\Gamma}}
\newunicodechar{Δ}{\ensuremath{\Delta}}
\newunicodechar{Ε}{\ensuremath{\varepsilon}}
\newunicodechar{Θ}{\ensuremath{\Theta}}
\newunicodechar{Λ}{\ensuremath{\Lambda}}
\newunicodechar{Μ}{\ensuremath{\mu}}
\newunicodechar{Τ}{\ensuremath{\tau}}
\newunicodechar{Φ}{\ensuremath{\Phi}}
\newunicodechar{Ψ}{\ensuremath{\Psi}}
\newunicodechar{Ω}{\ensuremath{\Omega}}
\newunicodechar{↦}{\ensuremath{\mapsto}}
\newunicodechar{∈}{\ensuremath{\in}}
\newunicodechar{∉}{\ensuremath{\notin}}
\newunicodechar{∧}{\ensuremath{\wedge}}
\newunicodechar{⇔}{\ensuremath{\Leftrightarrow}}
\newunicodechar{⟺}{\ensuremath{\Longleftrightarrow}}
\newunicodechar{⇒}{\ensuremath{\Rightarrow}}
\newunicodechar{⟹}{\ensuremath{\Longrightarrow}}
\newunicodechar{∘}{\ensuremath{\circ}}
\newunicodechar{∪}{\ensuremath{\cup}}
\newunicodechar{∩}{\ensuremath{\cap}}
\newunicodechar{≡}{\ensuremath{\equiv}}
\newunicodechar{⊂}{\ensuremath{\subset}}
\newunicodechar{⊥}{\ensuremath{\perp}}
\newunicodechar{∃}{\ensuremath{\exists}}
\newunicodechar{∀}{\ensuremath{\forall}}
\newunicodechar{∛}{\ensuremath{\sqrt[3]{\,}}}
\newunicodechar{∜}{\ensuremath{\sqrt[4]{\,}}}
\newunicodechar{⃗}{\ensuremath{{}^{\to}}}
\newunicodechar{ⁿ}{\ifmmode^{n}\else\textsuperscript{\ensuremath{n}}\fi}
\newunicodechar{ˣ}{\ifmmode^{x}\else\textsuperscript{\ensuremath{x}}\fi}
\newunicodechar{ᵇ}{\ifmmode^{b}\else\textsuperscript{\ensuremath{b}}\fi}
\newunicodechar{ʳ}{\ifmmode^{r}\else\textsuperscript{\ensuremath{r}}\fi}
\newunicodechar{ᵘ}{\ifmmode^{u}\else\textsuperscript{\ensuremath{u}}\fi}
\newunicodechar{ʸ}{\ifmmode^{y}\else\textsuperscript{\ensuremath{y}}\fi}
\newunicodechar{ᵃ}{\ifmmode^{a}\else\textsuperscript{\ensuremath{a}}\fi}
\newunicodechar{ᵉ}{\ifmmode^{e}\else\textsuperscript{\ensuremath{e}}\fi}
\newunicodechar{ᵏ}{\ifmmode^{k}\else\textsuperscript{\ensuremath{k}}\fi}
\newunicodechar{ᵖ}{\ifmmode^{p}\else\textsuperscript{\ensuremath{p}}\fi}
\newunicodechar{ᴺ}{\ifmmode^{N}\else\textsuperscript{\ensuremath{N}}\fi}
\newunicodechar{ᶜ}{\ifmmode^{c}\else\textsuperscript{\ensuremath{c}}\fi}
\newunicodechar{ʰ}{\ifmmode^{h}\else\textsuperscript{\ensuremath{h}}\fi}
\newunicodechar{ᵠ}{\ifmmode^{q}\else\textsuperscript{\ensuremath{q}}\fi}
\newunicodechar{ₙ}{\ifmmode_{n}\else\textsubscript{\ensuremath{n}}\fi}
\newunicodechar{ₐ}{\ifmmode_{a}\else\textsubscript{\ensuremath{a}}\fi}
\newunicodechar{ᵢ}{\ifmmode_{i}\else\textsubscript{\ensuremath{i}}\fi}
\newunicodechar{ᵣ}{\ifmmode_{r}\else\textsubscript{\ensuremath{r}}\fi}
\newunicodechar{ₖ}{\ifmmode_{k}\else\textsubscript{\ensuremath{k}}\fi}
\newunicodechar{ᵦ}{\ifmmode_{\beta}\else\textsubscript{\ensuremath{\beta}}\fi}
\newunicodechar{ₓ}{\ifmmode_{x}\else\textsubscript{\ensuremath{x}}\fi}
\newfontfamily\popdisplay{ArchivoBlack-Regular.ttf}[Path=../../../fonts/]
\newfontfamily\poplogo{SpaceGrotesk-Bold.ttf}[Path=../../../fonts/]
\newfontfamily\popsemi{PlusJakartaSans-SemiBold.ttf}[Path=../../../fonts/]
\newfontfamily\popxbold{PlusJakartaSans-ExtraBold.ttf}[Path=../../../fonts/]
\newfontfamily\popxboldls{PlusJakartaSans-ExtraBold.ttf}[Path=../../../fonts/,LetterSpace=3.0]

\setlist[enumerate]{leftmargin=1.7em,itemsep=5pt,topsep=4pt}
\setlist[itemize]{leftmargin=1.7em,itemsep=4pt,topsep=4pt,label={\color{poporange}\textbullet}}
\setlength{\parindent}{0pt}
\setlength{\parskip}{5pt}
\renewcommand{\arraystretch}{1.25}

% pgfplots : style Pop par défaut
\pgfplotsset{
  every axis/.append style={axis line style={popink,line width=1pt},tick style={popink},
    grid style={popline,line width=0.5pt},label style={font=\small},tick label style={font=\footnotesize}},
  popcurve/.style={poporange,line width=1.8pt,no markers,smooth},
}
% Même style disponible dans un \draw TikZ simple (hors pgfplots)
\tikzset{popcurve/.style={poporange,line width=1.8pt}}
\tikzset{popgrid/.style={popline,very thin}}
\colorlet{popcurve}{poporange} % le nom du style est parfois utilisé comme couleur

% ============================================================
% Pill / squiggle
% ============================================================
\newcommand{\poppill}[3]{%
  \tikz[baseline=-0.55ex]\node[draw=popink,line width=1.2pt,rounded corners=2.6pt,%
    fill=#2,inner xsep=6.5pt,inner ysep=3pt]{%
    \popxboldls\fontsize{7}{7}\selectfont\color{#3}\MakeUppercase{#1}};%
}
\newcommand{\popsquiggle}[1]{%
  \tikz[baseline=-0.4ex]{
    \draw[#1,line width=2.6pt,line cap=round]
      plot [smooth] coordinates {(0,0.09) (0.34,0.01) (0.68,0.21) (1.02,0.01) (1.36,0.10)};
  }%
}
% Mot-clé surligné (marqueur orange sous le texte)
\newcommand{\cle}[1]{{\popxbold\color{popdark}#1}}

% ============================================================
% En-tête / pied
% ============================================================
\pagestyle{fancy}
\fancyhf{}
\renewcommand{\headrulewidth}{0pt}
\renewcommand{\footrulewidth}{0pt}
\lhead{\raisebox{-0.15ex}{\poplogo\fontsize{13}{13}\selectfont\color{popink}OnBuch\color{poporange}+}}
\rhead{\poppill{\DOCMATIERE\ \textbullet\ \DOCNIVEAU\ \textbullet\ Leçon \DOCLECON}{white}{popink}}
\lfoot{\popxbold\fontsize{7.6}{7.6}\selectfont\color{popmuted}\MakeUppercase{Cours signé OnBuch+}\ \textbullet\ {\popsemi\DOCTITRE}}
\rfoot{\tikz[baseline=-0.6ex]\node[rounded corners=8.5pt,fill=popink,minimum height=17pt,minimum width=17pt,inner xsep=5pt,inner sep=0]{\popxbold\fontsize{8}{8}\selectfont\color{popcream}\thepage\,/\,\pageref*{LastPage}};}

% ============================================================
% Titres
% ============================================================
\newcommand{\chaptitre}[2]{%
  \begin{tcolorbox}[enhanced,colback=poporange,colframe=popink,boxrule=2.6pt,arc=11pt,
      drop shadow=popink,left=15pt,right=15pt,top=12pt,bottom=11pt,before skip=3pt,after skip=18pt]
    \poppill{#2}{popink}{popcream}\par\vspace{9pt}
    {\raggedright\hyphenpenalty=10000\exhyphenpenalty=10000\popdisplay\fontsize{22}{24}\selectfont\color{popcream}\MakeUppercase{#1}\par}
  \end{tcolorbox}
}
% Section numérotée : pastille orange avec le numéro + titre Archivo
\newcounter{popsec}
\newcommand{\coursec}[1]{%
  \stepcounter{popsec}\setcounter{popsub}{0}%
  \par\vspace{16pt}\noindent
  \begin{minipage}{\linewidth}
  \tikz[baseline=-0.7ex]\node[draw=popink,line width=1.6pt,fill=poporange,rounded corners=4pt,minimum size=22pt,inner sep=0]{\popdisplay\fontsize{12}{12}\selectfont\color{popcream}\thepopsec};%
  \hspace{8pt}{\popdisplay\fontsize{15}{17}\selectfont\color{popink}\MakeUppercase{#1}}\par
  \vspace{4pt}\noindent{\color{poporange}\rule{36pt}{3.6pt}}
  \end{minipage}\par\nopagebreak\vspace{8pt}
}
\newcounter{popsub}
\newcommand{\courssub}[1]{%
  \stepcounter{popsub}%
  \par\vspace{9pt}\noindent{\popxbold\fontsize{11.5}{13}\selectfont\color{popink}{\color{poporange}\thepopsec.\thepopsub}\hspace{6pt}#1}\par\nopagebreak\vspace{4pt}
}

% ============================================================
% Boîtes pédagogiques — même recette : fond blanc, bordure épaisse colorée,
% ombre dure encre, badge plein. Titre optionnel [#1].
% ============================================================
\tcbset{popbase/.style={breakable,enhanced,colback=white,boxrule=2.3pt,arc=9pt,
  shadow={3pt}{-3pt}{0pt}{popink},left=14pt,right=14pt,top=11pt,bottom=11pt,before skip=11pt,after skip=15pt,
  fontupper=\popsemi\fontsize{10.6}{14.5}\selectfont\color{popink},
  fontlower=\popsemi\fontsize{10.6}{14.5}\selectfont\color{popink}}}
% Coche dessinée (le glyphe ✓ n'existe pas dans les polices du document)
\newcommand{\popcheck}[1]{\tikz[baseline=-0.5ex]\draw[#1,line width=1.6pt,line cap=round,line join=round](-0.13,0.0)--(-0.04,-0.09)--(0.13,0.1);}
\providecommand{\coche}{\popcheck{popgreen}}
\providecommand{\resultat}[1]{\par\smallskip\noindent\fcolorbox{popgreen}{popgreenL}{\parbox{\dimexpr\linewidth-2\fboxsep-2\fboxrule\relax}{\textbf{Résultat.} #1}}\par\smallskip}
\newcommand{\popboxhead}[3]{\poppill{#1}{#2}{white}\ifx\relax#3\relax\else\hspace{6pt}{\popxbold\fontsize{10.6}{12}\selectfont\color{popink}#3}\fi\par\vspace{7pt}}

\newtcolorbox{aretenir}[1][]{popbase,colframe=popgreen,before upper={\popboxhead{À retenir}{popgreen}{#1}}}
\newtcolorbox{exemplebox}[1][]{popbase,colframe=poporange,before upper={\popboxhead{Exemple}{poporange}{#1}}}
\newtcolorbox{definition}[1][]{popbase,colframe=popblue,before upper={\popboxhead{Définition}{popblue}{#1}}}
\newtcolorbox{propriete}[1][]{popbase,colframe=poppurple,before upper={\popboxhead{Propriété}{poppurple}{#1}}}
\newtcolorbox{methode}[1][]{popbase,colframe=popgold,before upper={\popboxhead{Méthode}{popgold}{#1}}}
\newtcolorbox{attention}[1][]{popbase,colframe=poppink,before upper={\popboxhead{Attention}{poppink}{#1}}}
\newtcolorbox{experience}[1][]{popbase,colframe=popblue,colback=popblueL,before upper={\popboxhead{Expérience}{popblue}{#1}}}
% « Le savais-tu ? » : carte violette pleine (contexte, culture, Cameroun)
\newtcolorbox{savaistu}[1][]{popbase,colback=poppurple,colframe=popink,
  fontupper=\popsemi\fontsize{10.4}{14.2}\selectfont\color{white},
  before upper={\poppill{Le savais-tu\,?}{white}{poppurple}\ifx\relax#1\relax\else\hspace{6pt}{\popxbold\color{white}#1}\fi\par\vspace{7pt}}}
% Exercice résolu : énoncé en haut, \tcblower puis la solution sur fond crème
\newcounter{popexo}\newcounter{popexr}
\newtcolorbox{exoresolu}[1][]{popbase,colframe=poporange,colbacklower=popcream,
  segmentation style={popink,line width=1.2pt,dashed},
  before upper={\stepcounter{popexr}\popboxhead{Exercice résolu \thepopexr}{poporange}{#1}},
  before lower={\poppill{Solution}{popink}{popcream}\par\vspace{6pt}}}
% Exercice (non résolu en place) — la correction est dans la partie Corrigés
\newtcolorbox{exercice}[1][]{popbase,colframe=popink,
  before upper={\stepcounter{popexo}\popboxhead{Exercice \thepopexo}{popink}{#1}}}
% Corrigé
\newtcolorbox{corrige}[1][]{popbase,colframe=popgreen,colback=popgreenL,
  before upper={\popboxhead{Corrigé}{popgreen}{#1}}}
% Objectifs / prérequis / bilan
\newtcolorbox{objectifs}{popbase,colframe=popink,colback=poporangeL,
  before upper={\popboxhead{Objectifs de la leçon}{popink}{}}}
\newtcolorbox{prerequis}{popbase,colframe=popink,colback=popblueL,
  before upper={\popboxhead{Prérequis}{popblue}{}}}
\newtcolorbox{bilan}{popbase,colframe=popink,colback=white,boxrule=2.8pt,
  before upper={\popboxhead{Fiche bilan}{poporange}{L'essentiel à connaître pour l'examen}}}
% Figure encadrée avec légende
\newtcolorbox{popfigure}[1][]{enhanced,breakable=false,colback=white,colframe=popline,boxrule=1.4pt,arc=8pt,
  left=10pt,right=10pt,top=10pt,bottom=8pt,before skip=10pt,after skip=12pt,halign=center,
  lower separated=false,halign lower=center,
  fontlower=\popsemi\fontsize{9}{11.5}\selectfont\color{popmuted},
  #1}
\newcommand{\legende}[1]{\par\vspace{6pt}{\popsemi\fontsize{9}{11.5}\selectfont\color{popmuted}#1}}

% ============================================================
% Signature OnBuch+
% ============================================================
\newcommand{\poplogoinline}[1]{{\poplogo\fontsize{#1}{#1}\selectfont\color{popink}OnBuch\color{poporange}+}}
\newcommand{\signatureonbuch}{%
  \begin{tcolorbox}[enhanced,colback=popink,colframe=popink,boxrule=2.2pt,arc=10pt,
      drop shadow=poporange,left=14pt,right=14pt,top=10pt,bottom=10pt,before skip=0pt,after skip=0pt]
    \tikz[baseline=-0.7ex]\node[circle,fill=popgreen,draw=popcream,line width=1.3pt,minimum size=22pt,inner sep=0]{\popcheck{white}};%
    \hspace{9pt}\begin{minipage}[c]{0.62\linewidth}
      {\popxbold\fontsize{12}{13}\selectfont\color{popcream}Signé\ \textbullet\ {\poplogo OnBuch\color{poporange}+}}\par\vspace{2pt}
      {\raggedright\popsemi\fontsize{8}{10.5}\selectfont\color{popline}\MakeUppercase{Équipe pédagogique OnBuch+}\\\MakeUppercase{Conforme au programme officiel MINESEC}\par}
    \end{minipage}\hfill
    {\poplogo\fontsize{22}{22}\selectfont\color{popcream}OnBuch\color{poporange}+}
  \end{tcolorbox}%
}

% ============================================================
% Page de garde
% #1 titre · #2 matière · #3 niveau · #4 module · #5 n° leçon · #6 durée ·
% #7 description / objectifs (texte ou itemize)
% ============================================================
\newcommand{\pagegarde}[7]{%
  \thispagestyle{empty}
  \newgeometry{a4paper,margin=1.7cm,top=1.6cm,bottom=1.4cm}%
  \begin{tikzpicture}[remember picture,overlay]
    \fill[poporange!18] ([xshift=-3.2cm,yshift=-3.6cm]current page.north east) circle (4.6cm);
    \fill[poppurple!14] ([xshift=2.4cm,yshift=7.6cm]current page.south west) circle (3.8cm);
  \end{tikzpicture}%
  {\poplogo\fontsize{30}{30}\selectfont\color{popink}OnBuch\color{poporange}+}\hfill
  \raisebox{0.6ex}{\poppill{Cours signé \textbullet\ Premium}{popink}{popcream}}\par
  \vspace{1pt}\popsquiggle{poppurple}\par
  \vspace{8pt}
  \begin{tcolorbox}[enhanced,colback=poporange,colframe=popink,boxrule=2.8pt,arc=13pt,
      drop shadow=popink,left=18pt,right=18pt,top=12pt,bottom=13pt,after skip=10pt]
    \poppill{#2 \textbullet\ #3}{popink}{popcream}\hspace{5pt}\poppill{#4}{white}{popink}\par\vspace{8pt}
    {\popdisplay\fontsize{14}{14}\selectfont\color{popink}LEÇON #5}\par\vspace{4pt}
    {\raggedright\hyphenpenalty=10000\exhyphenpenalty=10000\popdisplay\fontsize{25}{27}\selectfont\color{popcream}\MakeUppercase{#1}\par}
    \vspace{8pt}
    {\popxbold\fontsize{9.5}{9.5}\selectfont\color{popink}\MakeUppercase{Durée conseillée : #6}}
  \end{tcolorbox}
  \begin{tcolorbox}[enhanced,colback=white,colframe=popink,boxrule=2.2pt,arc=10pt,
      drop shadow=popink,left=16pt,right=16pt,top=10pt,bottom=10pt,after skip=8pt]
    \poppill{Dans cette leçon}{poppurple}{white}\par\vspace{5pt}
    \popsemi\fontsize{9.3}{12.6}\selectfont\color{popink}#7
  \end{tcolorbox}
  \noindent\begin{minipage}[t]{0.487\linewidth}
    \begin{tcolorbox}[enhanced,colback=popgreen,colframe=popink,boxrule=2.2pt,arc=10pt,
        drop shadow=popink,left=11pt,right=11pt,top=10pt,bottom=10pt,height=3.1cm,valign=center]
      \tcbox[colback=white,colframe=popink,boxrule=1.3pt,arc=5pt,boxsep=0pt,nobeforeafter,
        left=3pt,right=3pt,top=3pt,bottom=3pt,valign=center]{\qrcode[height=1.9cm]{https://play.google.com/store/apps/details?id=com.luvvix.onbuch&hl=fr}}%
      \hspace{8pt}\begin{minipage}[c]{0.5\linewidth}
        {\popdisplay\fontsize{11}{12.5}\selectfont\color{white}\MakeUppercase{Télécharge OnBuch}}\par\vspace{4pt}
        {\raggedright\popsemi\fontsize{8.4}{11}\selectfont\color{white}Gratuit \textbullet\ Google Play\\Tuteur IA, annales, concours, résultats\par}
      \end{minipage}
    \end{tcolorbox}
  \end{minipage}\hfill
  \begin{minipage}[t]{0.487\linewidth}
    \begin{tcolorbox}[enhanced,colback=poppurple,colframe=popink,boxrule=2.2pt,arc=10pt,
        drop shadow=popink,left=11pt,right=11pt,top=10pt,bottom=10pt,height=3.1cm,valign=center]
      \tikz[baseline=-0.7ex]\node[circle,fill=white,draw=popink,line width=1.3pt,minimum size=1.6cm,inner sep=0]{\popdisplay\fontsize{24}{24}\selectfont\color{poppurple}+};%
      \hspace{8pt}\begin{minipage}[c]{0.6\linewidth}
        {\popdisplay\fontsize{11}{12.5}\selectfont\color{white}\MakeUppercase{Passe à OnBuch+}}\par\vspace{4pt}
        {\raggedright\popsemi\fontsize{8.4}{11}\selectfont\color{white}Tous les cours signés, Tuteur IA illimité, fiches PDF — onglet OnBuch+ de l'app\par}
      \end{minipage}
    \end{tcolorbox}
  \end{minipage}
  \vspace{8pt}
  \popcontenu
  \vfill
  \signatureonbuch
  \restoregeometry
  \newpage
}

% Rangée « Ce cours contient » (page de garde)
\newcommand{\poptile}[3]{% #1 couleur #2 gros texte #3 légende
  \begin{tcolorbox}[enhanced,colback=#1,colframe=popink,boxrule=2pt,arc=9pt,shadow={2.5pt}{-2.5pt}{0pt}{popink},
      left=6pt,right=6pt,top=9pt,bottom=9pt,halign=center,valign=center,height=1.75cm,width=\linewidth,nobeforeafter]
    {\popdisplay\fontsize{12.5}{13}\selectfont\color{white}\MakeUppercase{#2}}\par\vspace{3pt}
    {\centering\popsemi\fontsize{7.8}{9.5}\selectfont\color{white}#3\par}
  \end{tcolorbox}}
\newcommand{\popcontenu}{%
  \par\noindent{\popxboldls\fontsize{7.5}{7.5}\selectfont\color{popmuted}CE COURS SIGNÉ CONTIENT}\par\vspace{7pt}
  \noindent\begin{minipage}[t]{0.235\linewidth}\poptile{popblue}{Cours}{complet, illustré, pas à pas}\end{minipage}\hfill
  \begin{minipage}[t]{0.235\linewidth}\poptile{poporange}{Méthodes}{et exercices résolus}\end{minipage}\hfill
  \begin{minipage}[t]{0.235\linewidth}\poptile{poppink}{Exercices}{type examen + corrigés}\end{minipage}\hfill
  \begin{minipage}[t]{0.235\linewidth}\poptile{popgreen}{Fiche bilan}{l'essentiel à retenir}\end{minipage}\par}

% Sommaire
\newtcolorbox{sommaire}{enhanced,colback=white,colframe=popink,boxrule=2.2pt,arc=10pt,
  drop shadow=popink,left=18pt,right=18pt,top=13pt,bottom=13pt,before skip=6pt,after skip=20pt,
  before upper={\poppill{Sommaire}{poppurple}{white}\par\vspace{9pt}\popsemi\fontsize{10.6}{14.5}\selectfont\color{popink}}}

% Signature de fin de cours — poussée tout en bas de la dernière page
\newcommand{\findecours}{%
  \par\vfill\nopagebreak
  \begin{center}\popsquiggle{poporange}\end{center}\vspace{4pt}
  \signatureonbuch
}
\AtBeginDocument{\def\llbracket{\mathopen{[\mkern-3.5mu[}}\def\rrbracket{\mathclose{]\mkern-3.5mu]}}} % intervalles d'entiers (glyphe ⟦ absent de la police)
% Glyphes absents de Fira Math : redéfinis à partir de glyphes disponibles
\AtBeginDocument{%
  \def\setminus{\mathbin{\backslash}}\let\smallsetminus\setminus
  \def\vdots{\mathord{\vbox{\baselineskip3.2pt\lineskiplimit0pt\kern4pt\hbox{.}\hbox{.}\hbox{.}}}}%
  \def\ddots{\mathinner{\mkern1mu\raise6.5pt\hbox{.}\mkern2mu\raise3.6pt\hbox{.}\mkern2mu\raise0.7pt\hbox{.}\mkern1mu}}%
  \def\triangle{\mathord{\scalebox{0.9}{$\symup{\Delta}$}}}%
  \def\nabla{\mathord{\raisebox{\depth}{\scalebox{1}[-1]{$\symup{\Delta}$}}}}%
  \def\checkmark{\popcheck{popdark}}%
  \def\star{\mathbin{\ast}}\def\bigstar{\mathord{\ast}}%
  \def\diamond{\mathbin{\scalebox{0.6}{$\Diamond$}}}%
  \def\bigcup{\mathop{\mathchoice{\raisebox{-0.3em}{\scalebox{1.7}{$\cup$}}}{\scalebox{1.25}{$\cup$}}{\cup}{\cup}}\displaylimits}%
  \def\bigcap{\mathop{\mathchoice{\raisebox{-0.3em}{\scalebox{1.7}{$\cap$}}}{\scalebox{1.25}{$\cap$}}{\cap}{\cap}}\displaylimits}%
  \def\lesseqqgtr{\mathrel{\lessgtr}}\def\bigoplus{\mathop{\oplus}\displaylimits}%
}
\providecommand{\ssi}{si et seulement si} % abréviation fréquente
\AtBeginDocument{\providecommand{\wideparen}{\overparen}} % arc de cercle

%% FIGFIX-BEGIN v5 — correctifs globaux des figures (docs/onbuch-generation/tools/figfix)
\usepackage{adjustbox}\usepackage{etoolbox}\tcbuselibrary{raster}
% toute figure TikZ/pgfplots plus large que la ligne est réduite à la largeur disponible
\BeforeBeginEnvironment{tikzpicture}{\begin{adjustbox}{max width=\linewidth}}
\AfterEndEnvironment{tikzpicture}{\end{adjustbox}}
% légendes pgfplots lisibles par-dessus les courbes
\pgfplotsset{every axis/.append style={legend style={fill=white,fill opacity=0.92,text opacity=1,draw=black!15,inner sep=3pt}}}
% étiquettes d'arêtes (syntaxe edge["..."]) avec fond blanc
\tikzset{every edge quotes/.append style={fill=white,inner sep=1.2pt}}
%% FIGFIX-END
\begin{document}

\pagegarde{La Philosophie africaine}{Philosophie}{1re Tech}{Philosophie – classe de Première technique}{N01}{3 séances (fiche de progression harmonisée)}{Cette leçon initie l'élève de Première Technique à la problématique de la philosophie africaine : de la dénonciation des thèses ethnocentristes coloniales à l'émergence de courants contemporains (ethnophilosophie, philosophie sagace, nationale, critique), pour aboutir à la réorientation actuelle du débat vers la pertinence sociale et l'ancrage dans les défis du développement au Cameroun.
\vspace{4pt}
\begin{multicols}{2}\small
\begin{itemize}
\item Identifier et déconstruire les arguments des thèses ethnocentristes occidentales (Lévy-Bruhl, Hegel, etc.) sur l'absence de philosophie en Afrique.
\item Différencier les quatre grandes tendances de la philosophie africaine (ethnophilosophie, philosophie sagace, philosophies nationales-idéologiques, philosophie critique/professionnelle) et citer un représentant majeur pour chacune.
\item Analyser un extrait de texte (ex: Paulin Hountondji, Kwasi Wiredu, Fabien Eboussi Boulaga) pour en extraire la thèse centrale sur la nature de la philosophie africaine.
\item Construire un tableau comparatif (croquis heuristique) opposant « philosophie comme folklore » et « philosophie comme discours rationnel, argumenté et critique ".
\item Rédiger une dissertation structurée ou une explication de texte justifiant la réorientation actuelle du débat : « Pour une philosophie africaine de la libération et du développement ".
\item Mobiliser les concepts de « Raison », « Tradition », « Modernité », « Universalisme » et « Particularisme » dans une argumentation personnelle sur l'identité culturelle camerounaise.
\end{itemize}\end{multicols}}

\begin{sommaire}
\begin{multicols}{2}
\begin{enumerate}
\item 1. Le procès de la raison africaine : thèses ethnocentristes et coloniales
\item 2. La naissance de la philosophie africaine moderne : les quatre tendances fondatrices
\item 3. La réorientation actuelle : décolonialité, développement et pertinence sociale
\item 4. Atelier pratique : Analyse de documents et construction d'outils heuristiques (TP)
\item 5. Intégration : Éthique technique, gouvernance et philosophie africaine aujourd'hui
\item Activité d'intégration
\item Méthodes et exercices résolus
\item Exercices
\item Corrigés des exercices
\item Fiche bilan
\end{enumerate}
\end{multicols}
\end{sommaire}

\begin{objectifs}
\begin{itemize}
\item Identifier et déconstruire les arguments des thèses ethnocentristes occidentales (Lévy-Bruhl, Hegel, etc.) sur l'absence de philosophie en Afrique.
\item Différencier les quatre grandes tendances de la philosophie africaine (ethnophilosophie, philosophie sagace, philosophies nationales-idéologiques, philosophie critique/professionnelle) et citer un représentant majeur pour chacune.
\item Analyser un extrait de texte (ex: Paulin Hountondji, Kwasi Wiredu, Fabien Eboussi Boulaga) pour en extraire la thèse centrale sur la nature de la philosophie africaine.
\item Construire un tableau comparatif (croquis heuristique) opposant « philosophie comme folklore » et « philosophie comme discours rationnel, argumenté et critique ".
\item Rédiger une dissertation structurée ou une explication de texte justifiant la réorientation actuelle du débat : « Pour une philosophie africaine de la libération et du développement ".
\item Mobiliser les concepts de « Raison », « Tradition », « Modernité », « Universalisme » et « Particularisme » dans une argumentation personnelle sur l'identité culturelle camerounaise.
\item Évaluer la pertinence de la pensée africaine (ex: Ubuntu, Palabre, Njangi) pour résoudre des conflits concrets dans l'établissement ou le quartier (justice réparatrice, gouvernance locale).
\item Produire une carte mentale synthétisant l'évolution historiographique du débat philosophique africain des années 1940 à nos jours.
\end{itemize}
\end{objectifs}

\begin{prerequis}
\begin{itemize}
\item Définition philosophique de la raison, du logos, du mythe et de la tradition (notions de 4ème / Seconde).
\item Contexte historique de la colonisation et de la mission civilisatrice en Afrique (Histoire, classe de 3ème/2nde).
\item Méthode de l'explication de texte philosophique (identification thèse/arguments) et de la dissertation (problématiser, argumenter, conclure).
\item Notion de subjectivité, d'objectivité et de scientificité (Epistémologie, début 1ère).
\item Connaissance basique des grandes ethnies et chefferies traditionnelles du Cameroun (Géographie/Histoire locale).
\end{itemize}
\end{prerequis}

\newpage

\coursec{Le procès de la raison africaine : thèses ethnocentristes et coloniales}

Dans un lycée de Yaoundé, un élève affirme en cours que « la philosophie est une invention grecque, l'Afrique n'a que des contes et des légendes ». Un camarade lui répond en citant l'\cle{ubuntu}, la sagesse des anciens Bamiléké sur le vivre-ensemble et les écrits de Fabien Eboussi Boulaga. Le professeur demande : \textbf{comment dépasser ce choc des représentations pour construire une pensée philosophique africaine rigoureuse, ancrée dans nos réalités camerounaises et ouverte à l'universel ?}

Cette question nous plonge au cœur du « procès de la raison africaine ». Avant d'être une discipline académique, la philosophie africaine moderne est née d'une nécessité vitale : répondre à un déni systématique de l'humanité et de la rationalité des peuples africains, orchestré par la pensée occidentale des XVIII\textsuperscript{e} et XIX\textsuperscript{e} siècles. Comprendre ce procès, c'est déjà faire de la philosophie : c'est analyser les mécanismes d'un pouvoir épistémique qui a justifié la colonisation.

\courssub{La thèse hégélienne : l'Afrique hors de l'Histoire et de la Philosophie}

Georg Wilhelm Friedrich Hegel, dans son cours \emph{La Raison dans l'histoire} (1837), livre l'argumentaire le plus influent de l'exclusion de l'Afrique. Pour Hegel, l'Histoire est le progrès de la conscience de la liberté. Il divise le monde en trois parties : l'Orient (un seul est libre), le monde gréco-romain (quelques-uns sont libres), le monde germanique (tous sont libres). L'Afrique, elle, est purement et simplement exclue de ce mouvement dialectique.

\begin{exemplebox}[Extrait de \emph{La Raison dans l'histoire}]
« L'Afrique proprement dite, en deçà du Sahara, n'est pas une partie historique du monde ; elle n'a ni mouvement ni développement à exhiber [...] Ce que nous entendons par l'Afrique, c'est l'Afrique isolée, telle qu'elle est en elle-même, sans rapport avec le reste du monde [...] L'esprit n'y a pas encore atteint la conscience de sa liberté ; il s'y enfonce dans le sommeil de la nature, dans la pure subjectivité qui ne s'est pas encore objectivée en droit, en moralité, en État. »
\end{exemplebox}

\noindent La structure de l'argument hégélien repose sur trois piliers :
\begin{enumerate}
    \item \textbf{L'absence d'écriture} : L'histoire suppose des archives, donc l'écriture. L'oralité africaine est disqualifiée comme source historique.
    \item \textbf{L'absence d'État} : La politique hégélienne exige l'État comme réalisation de l'éthique substantielle. Les sociétés africaines sans État centralisé (sociétés à lignages, chefferies) sont reléguées dans la « nature ».
    \item \textbf{L'enfermement dans la nature} : L'Africain n'est pas un sujet historique qui se fait lui-même ; il est un objet naturel, passif, « endormi ».
\end{enumerate}

\begin{definition}[Ethnocentrisme]
Attitude consistant à prendre sa propre culture (ses valeurs, ses catégories de pensée, ses institutions) comme modèle universel de référence et à juger les autres cultures comme inférieures, arriérées ou déviantes. L'ethnocentrisme philosophique nie la pluralité des formes de rationalité.
\end{definition}

\noindent Cette thèse n'est pas une simple erreur d'appréciation : elle fonde une \textbf{ontologie coloniale}. Si l'Africain n'a pas d'histoire, il n'a pas de destin à construire ; il n'est qu'une matière brute à modeler. C'est la justification intellectuelle de la « mission civilisatrice ».

\begin{popfigure}
\begin{tikzpicture}[
    grow'=right,
    level distance=3.5cm,
    sibling distance=1.8cm,
    edge from parent/.style={very thick, draw=popdark, -{Stealth[length=3mm]}},
    every node/.style={font=\small, align=center, inner sep=4pt},
    root/.style={rectangle, rounded corners, fill=popblue, text=white, font=\bfseries, minimum width=3cm},
    branch/.style={rectangle, rounded corners, fill=popblueL, draw=popblue, text=popdark, minimum width=3.5cm},
    leaf/.style={ellipse, fill=poporangeL, draw=poporange, text=popdark, minimum width=3cm},
    fruit/.style={rectangle, rounded corners, fill=popred, text=white, font=\bfseries, minimum width=4cm}
]
\node[root, align=center]{Préjugé ontologique\\ \emph{(L'Africain = Nature)}}
    child { node[branch, align=center]{Pas d'Histoire\\ \emph{(Hegel)}}
        child { node[leaf] {Pas d'écriture} }
        child { node[leaf] {Pas d'État} }
        child { node[leaf] {Immuabilité} }
    }
    child { node[branch, align=center]{Mentalité pré-logique\\ \emph{(Lévy-Bruhl)}}
        child { node[leaf] {Participation mystique} }
        child { node[leaf] {Refus contradiction} }
        child { node[leaf] {Absence concept} }
    }
    child { node[branch, align=center]{Oralité = Irrationnel\\ \emph{(Ethnologie coloniale)}}
        child { node[leaf] {Mythe vs Logos} }
        child { node[leaf] {Tradition figée} }
        child { node[leaf] {Pas de critique} }
    }
    child { node[fruit, yshift=-1.5cm, align=center]{Infériorisation $\rightarrow$ Colonisation\\ \emph{Mission civilisatrice / Assimilation} } };
\end{tikzpicture}
\legende{Arbre des arguments ethnocentristes : du préjugé ontologique (racine) aux justifications coloniales (fruit) en passant par les thèses historico-philosophiques (branches) et anthropologiques (feuilles).}
\end{popfigure}

\courssub{L'ethnologie coloniale et la « mentalité pré-logique » de Lévy-Bruhl}

Alors que Hegel philosophe depuis son fauteuil berlinois, Lucien Lévy-Bruhl, professeur à la Sorbonne, prétend fonder scientifiquement l'infériorité intellectuelle dans \emph{Les Fonctions mentales dans les sociétés inférieures} (1910). Il oppose deux types de mentalité : la « mentalité logique » (occidentale, moderne) et la « mentalité pré-logique » (primitive, africaine).

\begin{definition}[Mentalité pré-logique (Lévy-Bruhl)]
Caractère attribué aux sociétés « primitives » : pensée mystique régie par la \cle{loi de participation} (l'être peut être à la fois lui-même et un autre, le symbole \emph{est} la chose), indifférence à la contradiction (violation du principe de non-contradiction), absence de concept abstrait et de raisonnement déductif formel.
\end{definition}

\noindent Le mécanisme de la \cle{participation} est central : « Les choses, les êtres, peuvent être à la fois eux-mêmes et autre chose qu'eux-mêmes. » Pour Lévy-Bruhl, l'Africain ne \emph{pense} pas le monde, il le \emph{vit} dans une confusion mystique. Le masque \emph{est} l'ancêtre ; le totem \emph{est} l'animal. Il n'y a pas de distance critique, pas de sujet s'objectivant face à l'objet.

\begin{attention}[Piège épistémologique]
Lévy-Bruhl confond \emph{logique formelle} (structure du raisonnement) et \emph{contenu culturel} (croyances, symboles). Un proverbe bamiléké comme « \emph{La main qui donne est au-dessus de celle qui reçoit} » utilise une structure logique (comparaison, inférence morale) pour exprimer une vérité sociale. Le nier, c'est confondre le \emph{quoi} (le mythe) et le \emph{comment} (la raison).
\end{attention}

\noindent Cette ethnologie n'est pas neutre. La \textbf{Mission Dakar-Djibouti (1931-1933)}, dirigée par Marcel Griaule, incarne l'ethnologie d'État. Son but : collecter des objets, fixer des « coutumes » immuables, photographier des « types raciaux ». On fige l'Africain dans un « traditionnel » intemporel pour nier sa modernité, sa capacité d'abstraction et son historicité. L'ethnologue devient le gardien du « vrai » Africain, celui d'avant la colonisation, confortant l'idée que la modernité (l'école, l'écriture, la raison) lui est étrangère.

\begin{aretenir}[Conséquences épistémologiques]
\begin{itemize}
    \item \textbf{Déni d'humanité} : La raison étant l'attribut définitoire de l'homme (Aristote : \emph{zoon logon echon}), nier la raison à l'Africain, c'est nier son humanité pleine.
    \item \textbf{Justification juridique} : Le Code de l'indigénat, le travail forcé, l'assimilation culturelle (écoles missionnaires interdisant les langues maternelles) trouvent leur fondement « scientifique » dans cette anthropologie négative.
    \item \textbf{Epistémicide} : Destruction des systèmes de savoirs locaux (médecine, agronomie, astronomie, droit coutumier) au profit du savoir unique occidental.
\end{itemize}
\end{aretenir}

\courssub{La riposte identitaire : la Négritude et ses limites}

Face à ce déni, la première riposte majeure vient de la diaspora et des intellectuels formés en France : la \cle{Négritude} (Aimé Césaire, Léopold Sédar Senghor, Léon-Gontran Damas). Dans \emph{Cahier d'un retour au pays natal} (1939) et \emph{Liberté I : Négritude et Humanisme} (1964), Senghor retourne l'argument de Lévy-Bruhl : « L'émotion est nègre, la raison est hellène. »

\begin{exemplebox}[Analyse de la thèse sénghorienne]
Senghor \textbf{accepte le découpage} de Lévy-Bruhl (Émotion vs Raison, Dionysiaque vs Apollinien) mais \textbf{inverse la valeur} : l'émotion, l'intuition, la participation, le rythme sont supérieurs à la raison froide, analytique, destructrice de l'Occident. L'Africain n'est pas « pré-logique », il est « supra-logique » ou « analogique ».
\end{exemplebox}

\noindent C'est une riposte \textbf{identitaire} salutaire : elle redonne fierté et dignité. Mais elle reste prisonnière du cadre colonial :
\begin{itemize}
    \item \textbf{Essentialisme} : Elle essentialise « l'Africain » (un bloc monolithique, immuable, naturellement rythmé/émotif).
    \item \textbf{Hétéro-référence} : Elle se définit encore par rapport à l'Occident (en négatif ou en miroir inversé).
    \item \textbf{Dépolitisation} : Elle gomme les différences de classes, de genres, de contextes historiques concrets au profit d'une « âme noire » abstraite.
\end{itemize}

\begin{attention}[Erreur fréquente]
Confondre « \cle{Négritude} » (mouvement littéraire et politique des années 1930-1960, essentialiste, poétique) et « \cle{Philosophie africaine} » (démarche rationnelle, critique, académique, argumentée, née dans les années 1970 avec Hountondji, Wiredu, Eboussi Boulaga, Mudimbe). La Négritude est une \emph{condition de possibilité} historique, mais pas la philosophie elle-même.
\end{attention}

\begin{popfigure}
\begin{tikzpicture}[
    mindmap,
    concept color=popblue,
    level 1 concept/.append style={level distance=4cm, sibling angle=120, concept color=popblueL},
    level 2 concept/.append style={level distance=3cm, sibling angle=60, concept color=popgreenL},
    level 3 concept/.append style={level distance=2.5cm, sibling angle=45, concept color=poporangeL},
    every node/.style={font=\small},
    text=popdark
]
\node[concept, align=center]{Négritude\\ (Senghor)}
    child[concept color=poppurple] { node[concept] {Essentialisme}
        child { node[concept] {Nature unique} }
        child { node[concept] {Intuition/Émotion} }
        child { node[concept] {Civilisation de l'Universel} }
    }
    child[concept color=popgreen] { node[concept] {Universalisme}
        child { node[concept] {Don au monde} }
        child { node[concept] {Métissage culturel} }
        child { node[concept] {Humanisme intégral} }
    }
    child[concept color=poporange] { node[concept] {Limites critiques}
        child { node[concept] {Figement identitaire} }
        child { node[concept] {Récupération néocoloniale} }
        child { node[concept] {Absence critique interne} }
    };
\end{tikzpicture}
\legende{Carte mentale : La Négritude de Senghor oscille entre essentialisme (identité close) et universalisme (don au monde), tout en portant les germes de sa propre critique par la philosophie africaine professionnelle.}
\end{popfigure}

\courssub{Le contexte camerounais : l'école coloniale et la formation d'auxiliaires}

Au Cameroun, le « procès de la raison » s'incarne dans l'institution scolaire. L'enseignement officiel colonial (allemand puis français) n'a pas vocation à former des philosophes.

\begin{savaistu}[L'école des fils de chefs]
Créée en 1920 à Yaoundé par l'administration française, l'« École des fils de chefs et de notables » (devenue plus tard l'École primaire supérieure) vise explicitement à former des \emph{auxiliaires de l'administration} : interprètes, commis, chefs de poste. Le programme : français, arithmétique élémentaire, hygiène, morale civique française, coutumes locales (codifiées par l'administration). \textbf{La philosophie, le latin, l'histoire générale en sont absents.} L'objectif est la docilité, non l'esprit critique.
\end{savaistu}

\noindent Cette stratégie produit une double coupure :
\begin{enumerate}
    \item \textbf{Coupure épistémique} : Le savoir légitime est le savoir français. Les savoirs endogènes (initiations, sociétés secrètes, pharmacopée, droit coutumier bamiléké ou bassa) sont relégués au rang de « coutumes » ou de « superstitions ».
    \item \textbf{Coupure sociale} : Une petite élite « évoluée » (les \emph{évolués}) se coupe de sa base culturelle, méprisant ses propres racines pour imiter le colonisateur, tout en étant barrée aux plus hautes responsabilités (plafond de verre).
\end{enumerate}

\noindent C'est dans ce contexte que des figures comme \textbf{Fabien Eboussi Boulaga} (1934-2018), jésuite, philosophe, rompent radicalement.

\begin{savaistu}[Fabien Eboussi Boulaga : « Philosophe en Afrique »]
Dans \emph{Christianisme sans fétiche} (1981) et \emph{La Crise du Muntu} (1984), Eboussi Boulaga refuse l'étiquette de « philosophe africain » qui l'assignerait à une identité essentialisée. Il revendique le titre de « \textbf{philosophe en Afrique} » : un sujet universel de la raison, situé dans un contexte africain, qui pense \emph{à partir de} l'Afrique (ses langues, ses problèmes concrets : dictature, pauvreté, aliénation religieuse) pour l'universel. Il applique la méthode phénoménologique et marxiste aux réalités camerounaises.
\end{savaistu}

\courssub{Méthode : Déconstruire un préjugé ethnocentriste (Démarche d'investigation)}

Pour répondre à la problématique de notre classe de Yaoundé, il ne suffit pas de polémiquer. Il faut une méthode rigoureuse de déconstruction.

\begin{methode}[Déconstruire un préjugé ethnocentriste en 4 étapes]
\begin{enumerate}
    \item \textbf{Citer l'auteur et l'extrait exact} : Ne pas caricaturer. Reproduire le texte incriminé (ex: Hegel § 173 ; Lévy-Bruhl Ch. III).
    \item \textbf{Identifier le présupposé implicite} : Quel critère universel non dit est utilisé comme mètre étalon ? (Ex: « L'histoire = Écriture » ; « Raison = Logique formelle aristotélicienne » ; « Philosophie = Tradition grecque »).
    \item \textbf{Analyser la fonction idéologique} : À qui profite ce déni ? Quel pouvoir politique, économique, épistémique légitime-t-il ? (Mission civilisatrice, extraction des ressources, domination symbolique).
    \item \textbf{Apporter le contre-exemple factuel ou logique} : Mobiliser l'ethnophilosophie critique, l'histoire des sciences africaines (mathématiques égyptiennes, métallurgie du fer au Cameroun, architecture soudanais), la logique formelle dans les proverbes, l'existence de concepts abstraits dans les langues africaines (ex: \emph{Muntu}, \cle{Nommo}, \cle{Ubuntu}).
\end{enumerate}
\end{methode}

\begin{exoresolu}[Application : Déconstruire « L'Afrique n'a pas de philosophie »]
\textbf{Énoncé :} Un manuel scolaire de 1950 affirme : « Le nègre ne philosophe pas ; il n'a pas de concepts abstraits, il vit dans le concret. »

\textbf{Solution détaillée :}
\begin{enumerate}
    \item \textbf{Citation} : Extrait du manuel, p. 12. Auteur anonyme (voix de l'institution coloniale).
    \item \textbf{Présupposé} : La philosophie = concepts abstraits \emph{séparés} de l'expérience vécue (idéalisme platonicien/kantien) + Langue philosophique = Grec/Allemand/Français. L'oralité ne peut pas porter de concept.
    \item \textbf{Fonction} : Justifier l'enseignement du français seul, l'interdiction des langues nationales à l'école, la suppression des initiations traditionnelles (lieux de transmission du savoir philosophique).
    \item \textbf{Contre-exemples} :
        \begin{itemize}
            \item \textbf{Logique formelle} : Le syllogisme chez les Dogon (recherche de Griaule/Dieterlen, mais réinterprété) ; les proverbes bamiléké structurés en \emph{modus ponens} (\emph{Si tu respectes l'aîné, tu gagnes sa bénédiction ; tu respectes l'aîné ; donc tu gagnes sa bénédiction}).
            \item \textbf{Concepts abstraits} : \emph{Muntu} (force vitale/être) chez les Bantu ; \emph{Nommo} (parole créatrice) chez les Dogon ; \cle{Ubuntu} (\emph{Umuntu ngumuntu ngabantu} : « Je suis parce que nous sommes ») — ontologie relationnelle rigoureuse.
            \item \textbf{Histoire des sciences et techniques} : L'écriture Shümom (système syllabique inventé par le sultan Njoya vers 1896), les mathématiques de l'Égypte antique (Papyrus de Rhind), la métallurgie du fer à Djoba (Adamaoua) dès 500 av. J.-C.
        \end{itemize}
    \item \textbf{Conclusion} : L'affirmation est fausse factuellement et logiquement circulaire (elle définit la philosophie pour exclure l'Afrique, puis constate l'exclusion pour valider la définition).
\end{enumerate}
\end{exoresolu}

\begin{aretenir}[Synthèse de la section]
Le « procès de la raison africaine » n'est pas un débat d'historiens : c'est un \textbf{dispositif de pouvoir}. Hegel (philosophie de l'histoire), Lévy-Bruhl (anthropologie scientifique), l'école coloniale (institution) forment un système cohérent qui dénie à l'Africain le statut de sujet pensant pour le réduire à objet de domination. La Négritude a brisé le silence, mais en essentialisant la réponse. La tâche de la philosophie africaine moderne — et de ce cours — est de passer de l'\emph{identité} (essentialisme) à la \emph{critique} (universel concret), en s'appuyant sur nos réalités camerounaises (langues, histoires, problèmes techniques, éthiques, politiques) pour produire une pensée rigoureuse.

\medskip\noindent\textbf{Transition :} Maintenant que le procès est instruit et la défense entendue, comment la philosophie africaine s'est-elle constituée comme discipline académique autonome dans les années 1970 ? C'est l'objet de la section 2 : les quatre tendances fondatrices (Ethnophilosophie, Nationaliste-idéologique, Professionnelle, Sage).
\end{aretenir}


\coursec{La naissance de la philosophie africaine moderne : les quatre tendances fondatrices}

Nous avons vu dans la section précédente comment la pensée occidentale, de Hegel à Lévy-Bruhl, a dressé un véritable \cle{procès de la raison africaine} : l'Africain y était présenté comme un être « prélogique », incapable d'abstraction et de pensée conceptuelle. Face à cette humiliation intellectuelle, les penseurs africains et certains occidentaux bienveillants ont réagi. Mais ils ne se sont pas accordés sur la réponse à apporter. Faut-il montrer que les cultures africaines contiennent déjà une philosophie ? Faut-il interroger les sages du village ? Faut-il forger des doctrines politiques pour les jeunes États indépendants ? Ou faut-il exiger les mêmes critères de rigueur que partout ailleurs dans le monde ? De ces réponses divergentes sont nées \textbf{quatre grandes tendances} de la philosophie africaine moderne, que nous allons étudier tour à tour : l'\cle{ethnophilosophie}, la \cle{philosophie sagace}, les \cle{philosophies nationales-idéologiques} et la \cle{philosophie critique ou professionnelle}.

\courssub{L'ethnophilosophie : la philosophie comme sagesse collective implicite}

\textbf{Intuition.} Quand un ancien du village explique, au moment d'un deuil, que « le mort n'est pas mort, il continue de veiller sur les siens », il exprime une vision du monde : une conception de la vie, de la mort, de la personne. L'ethnophilosophie part de cette observation : les cultures africaines traditionnelles contiennent une pensée profonde, mais celle-ci n'est pas écrite dans des livres ; elle est \emph{implicite}, cachée dans les mythes, les proverbes, les rites, la langue et les institutions.

\begin{definition}[L'ethnophilosophie]
L'\cle{ethnophilosophie} est l'étude systématique des visions du monde implicites dans les cultures africaines traditionnelles (mythes, rites, langue, proverbes, institutions), présentées comme une philosophie \textbf{collective et anonyme}, partagée par tous les membres du groupe.
\end{definition}

\textbf{Le texte fondateur.} L'ouvrage qui inaugure cette tendance est \emph{La Philosophie bantoue} (1945) du missionnaire belge Placide \cle{Tempels}. Vivant au Congo, Tempels observe attentivement les comportements, les croyances et le droit coutumier des Bantous. Il en tire une thèse audacieuse pour l'époque : les Bantous ont une ontologie, c'est-à-dire une théorie de l'être. Pour eux, l'être n'est pas défini par la substance (comme chez Aristote) mais par la \cle{force vitale} : « être, c'est posséder une force », et toute la réalité — Dieu, les ancêtres, les vivants, les animaux, les choses — forme une hiérarchie de forces qui s'influencent mutuellement. Le bien est ce qui augmente la force vitale ; le mal est ce qui la diminue.

\textbf{Les continuateurs.} Le prêtre rwandais Alexis \cle{Kagamé} systématise cette démarche dans \emph{La Philosophie bantu-rwandaise de l'être} (1956), en analysant la langue kinyarwanda pour y dégager des catégories philosophiques comparables à celles d'Aristote. Le théologien kényan John \cle{Mbiti}, dans \emph{African Religions and Philosophy} (1969), montre que la conception africaine du temps et de la communauté — résumée par sa formule célèbre « Je suis parce que nous sommes, et parce que nous sommes, je suis » — constitue une authentique pensée philosophique.

\begin{exemplebox}[La force vitale en situation concrète]
Chez de nombreux peuples du Cameroun, lorsqu'un enfant tombe malade, on interroge d'abord l'état des relations familiales : un conflit non réglé avec un ancêtre ou un parent aurait « affaibli » la force de l'enfant. La guérison passe alors par la réconciliation et le rite, autant que par le remède. Cette logique illustre exactement l'ontologie de la force vitale décrite par Tempels : la santé est un équilibre de forces relationnelles, pas seulement un état biologique.
\end{exemplebox}

\textbf{Portée et limite.} Le mérite de l'ethnophilosophie est immense : elle a \emph{réhabilité} la raison africaine au moment même où l'Occident la niait. Mais elle a une faiblesse décisive : elle présente la philosophie comme une pensée \textbf{anonyme, collective et unanimiste}, où tous les Africains penseraient la même chose sans débat ni critique. Or une pensée qui n'admet ni discussion ni désaccord ressemble davantage à une \emph{cosmologie} ou à une \emph{vision du monde} qu'à une philosophie au sens strict. C'est la critique que lui adresseront Hountondji et les philosophes professionnels.

\courssub{La philosophie sagace : la critique individuelle des sages}

\textbf{Intuition.} L'ethnophilosophie prétend que « tout le peuple pense ainsi ». Mais n'importe quel observateur d'un village camerounais sait que tous les anciens ne se valent pas : certains répètent les croyances communes, d'autres les interrogent, les critiquent, proposent des idées personnelles. C'est cette différence que le philosophe kényan Henry Odera \cle{Oruka} (1944-1995) a voulu saisir.

\begin{definition}[La philosophie sagace]
La \cle{philosophie sagace}, notion forgée par Henry Odera Oruka, désigne la pensée rationnelle, critique et \textbf{individuelle} de sages africains — lettrés ou illettrés — identifiée par \textbf{enquête de terrain}. Oruka la distingue soigneusement de la \emph{philosophie populaire}, qui n'est que l'ensemble des croyances communes, proverbes et maximes acceptés sans examen par la communauté.
\end{definition}

\textbf{La démarche d'investigation.} Oruka ne reste pas dans sa bibliothèque : il mène une véritable \emph{enquête de terrain}, comparable à une investigation scientifique. Le protocole est le suivant : repérer dans une communauté les personnes réputées sages ; les interroger longuement, non pas sur « ce que le peuple croit », mais sur \emph{ce qu'elles pensent personnellement} ; puis tester la cohérence rationnelle de leurs réponses, en les poussant à justifier, à argumenter, voire à douter des traditions. Deux résultats apparaissent alors :
\begin{itemize}
\item le \textbf{sage populaire} (ou « sage-folk ») : il maîtrise parfaitement la sagesse collective, mais la répète sans la critiquer ;
\item le \textbf{sage philosophique} (ou « sage-philosophe ») : il est capable de distance critique, d'argumentation personnelle, parfois de scepticisme à l'égard des croyances reçues. C'est lui, et lui seul, qui fait de la philosophie.
\end{itemize}

\begin{exemplebox}[Des sages camerounais]
Appliquée au Cameroun, la méthode d'Oruka consisterait à interroger des sages \textbf{Bamoun} autour de la cour de Foumban ou des notables \textbf{Bassa} de la région de Douala : certains se contenteraient de citer les proverbes des ancêtres (philosophie populaire), tandis que d'autres expliqueraient pourquoi tel proverbe est juste, dans quelles conditions il s'applique, et pourquoi tel autre doit être nuancé ou rejeté. Ces derniers sont des philosophes au sens plein, même s'ils n'ont jamais écrit une ligne.
\end{exemplebox}

\textbf{Apport décisif.} La philosophie sagace réconcilie deux exigences : elle reconnaît que la philosophie africaine peut être \emph{orale} (contre les tenants de l'écriture obligatoire), tout en exigeant qu'elle soit \emph{individuelle et critique} (contre l'unanimisme ethnophilosophique). Elle prouve concrètement que la rationalité critique existe en Afrique, y compris chez des illettrés.

\courssub{Les philosophies nationales-idéologiques : penser pour construire la nation}

\textbf{Intuition.} Dans les années 1960, les États africains accèdent à l'indépendance. Leurs dirigeants font face à une question immense : sur quelles idées construire la nation, l'économie et l'État ? Plusieurs chefs d'État-philosophes répondent en forgeant des doctrines politiques inspirées des traditions africaines.

\begin{definition}[Philosophie nationale-idéologique]
On appelle \cle{philosophies nationales-idéologiques} les doctrines élaborées par des chefs d'État africains après les indépendances, qui mobilisent des valeurs traditionnelles (communauté, solidarité, dialogue) pour légitimer un projet de construction nationale et de développement.
\end{definition}

Les figures majeures de cette tendance sont :
\begin{itemize}
\item Kwame \cle{Nkrumah} (Ghana) et le \textbf{consciencisme} : une synthèse entre l'héritage culturel africain, l'expérience coloniale et le socialisme, destinée à guider la révolution africaine ;
\item Julius \cle{Nyerere} (Tanzanie) et l'\textbf{Ujamaa} (« famille élargie » en swahili) : un socialisme communautaire fondé sur le partage, l'entraide villageoise et la dignité du travail, traduit en politique de villagisation et d'éducation populaire ;
\item Léopold Sédar \cle{Senghor} (Sénégal) et la \textbf{Négritude} prolongée en socialisme africain : la valorisation de l'émotion, du rythme et de la culture nègre comme contribution au « rendez-vous du donner et du recevoir » ;
\item Sékou \cle{Touré} (Guinée) et son \textbf{humanisme intégral}, visant l'épanouissement total de l'homme africain libéré de la domination.
\end{itemize}

\begin{attention}[Ne pas réduire ces philosophies à de la propagande]
Il serait injuste de traiter les philosophies nationales comme de simples slogans de partis uniques. Elles posent un véritable problème philosophique et politique : qu'est-ce qui \textbf{légitime} un État post-colonial ? Sur quelles valeurs fonder la justice et le développement quand les institutions ont été imposées de l'extérieur ? L'Ujamaa de Nyerere, par exemple, est une réflexion cohérente sur la propriété, l'égalité et la communauté. Leur limite est ailleurs : en s'identifiant au pouvoir d'État, ces doctrines ont parfois étouffé le débat critique, condition pourtant essentielle de toute philosophie.
\end{attention}

\courssub{La philosophie critique ou professionnelle : l'exigence de rigueur}

\textbf{Intuition.} Un mathématicien camerounais ne fait pas des « mathématiques bantoues » : il fait des mathématiques, avec les mêmes critères de preuve que partout dans le monde. Pourquoi la philosophie échapperait-elle à cette exigence ? C'est le point de départ de la tendance critique.

\begin{definition}[Philosophie critique ou professionnelle]
La \cle{philosophie critique} (ou professionnelle) est la tendance qui exige que la philosophie africaine obéisse aux standards universels de la discipline : \textbf{écriture}, \textbf{argumentation individuelle}, \textbf{discussion rationnelle} et rigueur académique. Elle refuse l'unanimisme (l'idée que « tous les Africains pensent pareil ») et le folklore (la réduction de la philosophie à l'ethnologie des coutumes).
\end{definition}

\textbf{Le manifeste de 1973.} L'ouvrage fondateur est \emph{Sur la « philosophie africaine »} (1973) du Béninois Paulin \cle{Hountondji}. Sa thèse est provocante : ce que Tempels appelait « philosophie bantoue » n'est pas de la philosophie, mais de l'\emph{ethnophilosophie}, c'est-à-dire une description ethnologique de croyances collectives. La philosophie, elle, est toujours une démarche \emph{individuelle}, \emph{écrite} et \emph{critique} : elle suppose un auteur qui prend position, argumente et s'expose à la contradiction. Hountondji définit donc la philosophie africaine non par un contenu culturel, mais simplement comme « l'ensemble des textes écrits par des Africains et qualifiés de philosophiques par leurs auteurs ».

\begin{exemplebox}[Un exemple chiffré]
En 1973, \emph{Sur la « philosophie africaine »} de Hountondji mobilise plus de \textbf{50 références} occidentales et africaines (de Husserl à Tempels, de Marx à Kagamé). Ce fait matériel est symboliquement décisif : il marque l'entrée du débat africain dans la conversation philosophique \emph{mondiale}, avec ses codes de citation, de réfutation et de publication.
\end{exemplebox}

\textbf{Les autres figures.} Le Ghanéen Kwasi \cle{Wiredu} plaide pour une « décolonisation conceptuelle » : examiner les concepts africains avec les outils de la logique analytique, sans complexe ni soumission. Le Nigérian Peter \cle{Bodunrin} demande aux philosophes africains de débattre des problèmes universels (vérité, connaissance, démocratie) plutôt que de décrire les coutumes. Le Camerounais Fabien \cle{Eboussi Boulaga} montre, dans \emph{La Crise du Muntu} (1977), que l'ethnophilosophie est un piège : en cherchant à prouver à l'Occident que « nous aussi, nous avons une philosophie », on reste prisonnier de son regard. Il faut au contraire partir de la \emph{crise} vécue par l'Africain concret pour produire une pensée libre et créatrice.

\begin{propriete}[La querelle de la philosophie africaine]
La « \cle{querelle de la philosophie africaine} » (années 1970) oppose les ethnophilosophes (Tempels, Kagamé, Mbiti) aux philosophes critiques (Hountondji, Wiredu, Bodunrin, Eboussi Boulaga) sur le \textbf{critère de scientificité} de la philosophie : pour les premiers, l'oralité, l'anonymat et la tradition collective suffisent ; pour les seconds, seules comptent l'écriture, l'argumentation individuelle et le débat rationnel. La philosophie sagace d'Oruka occupe une position intermédiaire : oralité possible, mais critique individuelle exigée.
\end{propriete}

\courssub{Tableau comparatif des quatre tendances}

\begin{popfigure}
\begin{tikzpicture}[x=1cm,y=1cm]
% Grille complète
\draw[popline, thick] (0,0) grid[step=1] (16.3,-6.8);
% En-tête
\fill[popblue] (0,0) rectangle (16.3,-0.8);
\node[white,font=\bfseries\small] at (1.55,-0.4) {Tendance};
\node[white,font=\bfseries\small] at (4.75,-0.4) {Objet d'étude};
\node[white,font=\bfseries\small] at (8.05,-0.4) {Méthode};
\node[white,font=\bfseries\small] at (11.35,-0.4) {Figures clés};
\node[white,font=\bfseries\small] at (14.65,-0.4) {Critique principale};
% Lignes de séparation verticales (déjà dans grid)
% Lignes horizontales (déjà dans grid)
% Contenu - Ligne 1 : Ethnophilosophie
\fill[popblueL] (0,-1.7) rectangle (16.3,-0.8);
\node[font=\small\bfseries,text=popdark,align=center] at (1.55,-1.25) {Ethno-\\philosophie};
\node[font=\scriptsize,align=center,text width=3cm] at (4.75,-1.25) {Vision du monde implicite : mythes, rites, langue};
\node[font=\scriptsize,align=center,text width=3cm] at (8.05,-1.25) {Analyse des cultures et des langues};
\node[font=\scriptsize,align=center,text width=3cm] at (11.35,-1.25) {Tempels, Kagamé, Mbiti};
\node[font=\scriptsize,align=center,text width=3cm] at (14.65,-1.25) {Unanimisme, anonymat, folklore};
% Ligne 2 : Sagacité
\fill[poporangeL] (0,-3.4) rectangle (16.3,-2.5);
\node[font=\small\bfseries,text=popdark,align=center] at (1.55,-2.95) {Philosophie\\sagace};
\node[font=\scriptsize,align=center,text width=3cm] at (4.75,-2.95) {Pensée critique individuelle des sages};
\node[font=\scriptsize,align=center,text width=3cm] at (8.05,-2.95) {Enquête de terrain, entretiens};
\node[font=\scriptsize,align=center,text width=3cm] at (11.35,-2.95) {Oruka (et sages Bamoun, Bassa)};
\node[font=\scriptsize,align=center,text width=3cm] at (14.65,-2.95) {Difficile à vérifier, corpus restreint};
% Ligne 3 : Nationale
\fill[popgreenL] (0,-5.1) rectangle (16.3,-4.2);
\node[font=\small\bfseries,text=popdark,align=center] at (1.55,-4.65) {Philosophies\\nationales-\\idéologiques};
\node[font=\scriptsize,align=center,text width=3cm] at (4.75,-4.65) {Construction de l'État et du développement};
\node[font=\scriptsize,align=center,text width=3cm] at (8.05,-4.65) {Doctrine politique, discours, programmes};
\node[font=\scriptsize,align=center,text width=3cm] at (11.35,-4.65) {Nkrumah, Nyerere, Senghor, Touré};
\node[font=\scriptsize,align=center,text width=3cm] at (14.65,-4.65) {Risque d'idéologie d'État, peu de débat};
% Ligne 4 : Critique
\fill[poppurpleL] (0,-6.8) rectangle (16.3,-5.9);
\node[font=\small\bfseries,text=popdark,align=center] at (1.55,-6.35) {Philosophie\\critique /\\professionnelle};
\node[font=\scriptsize,align=center,text width=3cm] at (4.75,-6.35) {Problèmes universels, textes argumentés};
\node[font=\scriptsize,align=center,text width=3cm] at (8.05,-6.35) {Écriture, logique, débat académique};
\node[font=\scriptsize,align=center,text width=3cm] at (11.35,-6.35) {Hountondji, Wiredu, Bodunrin, Eboussi Boulaga};
\node[font=\scriptsize,align=center,text width=3cm] at (14.65,-6.35) {Risque de couper la philosophie du vécu africain};
\end{tikzpicture}
\legende{Tableau comparatif des quatre tendances fondatrices de la philosophie africaine moderne : objet d'étude, méthode, figures clés et critique principale adressée à chacune.}
\end{popfigure}

\begin{methode}[Classer un auteur dans l'une des quatre tendances]
Pour situer un auteur ou un texte, pose-toi quatre questions dans l'ordre :
\begin{enumerate}
\item Cherche-t-il la philosophie dans la \textbf{tradition collective} (mythes, proverbes, langue) ? $\rightarrow$ \emph{Ethnophilosophie}.
\item Interroge-t-il des \textbf{individus sages} pour tester leur pensée critique personnelle ? $\rightarrow$ \emph{Philosophie sagace}.
\item Propose-t-il une \textbf{idéologie d'État} pour construire la nation post-indépendance ? $\rightarrow$ \emph{Philosophie nationale-idéologique}.
\item Exige-t-il les \textbf{standards académiques internationaux} (écriture, argumentation, débat) ? $\rightarrow$ \emph{Philosophie critique}.
\end{enumerate}
Un même penseur peut évoluer : Hountondji, par exemple, a nuancé ses positions de 1973 dans ses travaux ultérieurs.
\end{methode}

\courssub{Chronologie : de la réhabilitation à la critique}

\begin{popfigure}
\begin{tikzpicture}
\begin{axis}[
width=\linewidth, height=6.2cm,
axis lines=middle,
xmin=1940, xmax=2010,
ymin=0, ymax=3.2,
xtick={1945,1960,1970,1973,1980,1990,2000},
xticklabels={1945,1960,1970,1973,1980,1990,2000},
ytick=\empty,
xlabel={Années}, xlabel style={anchor=west},
every axis x label/.style={at={(ticklabel* cs:1.02)}},
axis line style={-Latex, thick, popdark},
tick style={popdark},
]
\addplot[only marks, mark=*, mark size=2.5pt, poporange] coordinates {(1945,1) (1969,1.6) (1973,2.2) (1980,1.6) (1990,2.2) (2000,1.6)};
\node[font=\scriptsize, align=center, text width=2.6cm, anchor=south, popdark] at (axis cs:1945,1.05) {Tempels, \emph{Philosophie bantoue}};
\node[font=\scriptsize, align=center, text width=2.6cm, anchor=south, popdark] at (axis cs:1969,1.65) {Mbiti ; Oruka lance la sagacité};
\node[font=\scriptsize, align=center, text width=2.8cm, anchor=south, popdark] at (axis cs:1973,2.25) {Hountondji, \emph{Sur la « philosophie africaine »}};
\node[font=\scriptsize, align=center, text width=2.6cm, anchor=south, popdark] at (axis cs:1980,1.65) {Wiredu, \emph{Philosophy and an African Culture}};
\node[font=\scriptsize, align=center, text width=2.6cm, anchor=south, popdark] at (axis cs:1990,2.25) {Mudimbe, \emph{The Invention of Africa}};
\node[font=\scriptsize, align=center, text width=2.6cm, anchor=south, popdark] at (axis cs:2000,1.65) {Essor du courant décolonial};
\end{axis}
\end{tikzpicture}
\legende{Repères chronologiques : de la réhabilitation ethnophilosophique (1945) à la critique professionnelle (1973-1980), puis à la réflexion sur l'invention coloniale de l'Afrique (Mudimbe) et à la décolonialité (années 2000).}
\end{popfigure}

\courssub{Focus Cameroun : l'Université de Yaoundé, foyer de la critique}

Dans les années 1970-1980, l'\cle{Université de Yaoundé} devient l'un des centres majeurs du débat philosophique africain. Trois noms doivent être retenus :
\begin{itemize}
\item Fabien \cle{Eboussi Boulaga}, avec \emph{La Crise du Muntu} (1977) : il démonte la logique de l'ethnophilosophie et appelle à une philosophie née de l'expérience vécue de la domination et de la liberté retrouvée ;
\item Jean-Marc \cle{Ela}, prêtre et sociologue-philosophe, auteur de \emph{Le Cri de l'homme africain} (1980) : il exige une pensée qui parte des conditions concrètes des paysans et des pauvres, non des spéculations abstraites ;
\item Patrice \cle{Nganang}, écrivain et penseur, prolonge plus tard cette exigence en interrogeant la mémoire, la langue et la citoyenneté camerounaises.
\end{itemize}
Le Cameroun illustre ainsi à lui seul le passage de la philosophie africaine du statut de curiosité ethnologique à celui de discipline critique engagée dans les problèmes de la cité.

\begin{exoresolu}[Identifier une tendance]
\textbf{Énoncé.} Un chercheur camerounais passe deux ans dans un village de l'Ouest. Il y interroge un notable réputé pour sa sagesse et découvre que celui-ci conteste personnellement plusieurs croyances de sa communauté sur la sorcellerie, qu'il juge irrationnelles, et qu'il en propose une explication sociale argumentée. Le chercheur publie les entretiens en les présentant comme de la philosophie. À quelle tendance cette démarche appartient-elle ? Justifie.
\tcblower
\textbf{Solution détaillée.} Il s'agit de la \textbf{philosophie sagace} d'Henry Odera Oruka. Trois indices convergent : (1) la méthode est une \emph{enquête de terrain} auprès d'un sage identifié par sa réputation ; (2) le chercheur ne décrit pas les croyances collectives du village (ce serait de la philosophie populaire, voire de l'ethnophilosophie), mais la pensée \emph{individuelle} du notable ; (3) ce sage manifeste une attitude \emph{critique} : il conteste, argumente et propose une explication rationnelle — c'est donc un « sage-philosophe » et non un simple « sage populaire ». La publication des entretiens confirme enfin l'entrée de cette pensée orale dans le débat écrit.
\end{exoresolu}

\begin{aretenir}[L'essentiel de la section]
\begin{itemize}
\item La philosophie africaine moderne naît d'une \textbf{réponse au procès de la raison africaine}, mais elle se divise en quatre tendances.
\item \textbf{Ethnophilosophie} (Tempels, Kagamé, Mbiti) : la philosophie est implicite dans la culture collective (force vitale) — mérite : la réhabilitation ; limite : l'unanimisme.
\item \textbf{Philosophie sagace} (Oruka) : la pensée critique \emph{individuelle} des sages, saisie par enquête de terrain, distincte des croyances populaires.
\item \textbf{Philosophies nationales-idéologiques} (Nkrumah, Nyerere, Senghor, Touré) : la pensée au service de la construction nationale — légitimes mais exposées au risque idéologique.
\item \textbf{Philosophie critique} (Hountondji, Wiredu, Bodunrin, Eboussi Boulaga) : exigence d'écriture, d'argumentation et de rigueur universelle ; \emph{Sur la « philosophie africaine »} (1973) en est le manifeste.
\item Le Cameroun, avec l'Université de Yaoundé (Eboussi Boulaga, Ela, Nganang), a joué un rôle de premier plan dans cette histoire.
\end{itemize}
\end{aretenir}

\coursec{La réorientation actuelle : décolonialité, développement et pertinence sociale}

Dans la section précédente, nous avons vu comment la philosophie africaine moderne s'est constituée en quatre tendances fondatrices (ethnophilosophie, philosophie nationaliste-idéologique, philosophie critique et philosophie herméneutique) pour répondre au procès fait à la raison africaine. Mais à force de vouloir \emph{prouver} que l'Afrique a une philosophie, les penseurs africains ne risquaient-ils pas de rester prisonniers d'une question posée par l'Occident ? C'est le constat dressé à partir des années 1980 : la querelle de la légitimité (\og existe-t-il une philosophie africaine ? \fg{}) doit céder la place à une question plus urgente et plus féconde : \emph{à quoi sert} la philosophie africaine ? C'est toute la réorientation actuelle du débat que nous étudions maintenant.

\courssub{De la querelle de légitimité à la querelle de pertinence}

\textbf{Intuition.} Imagine un élève de Première technique qui passe tout son temps à démontrer à ses camarades qu'il est intelligent, au lieu d'utiliser son intelligence pour résoudre des problèmes concrets. On dira qu'il gaspille son énergie dans une justification stérile. La philosophie africaine a longtemps été dans cette situation : obsédée par la preuve de son existence, elle risquait d'oublier sa mission.

\begin{definition}[Querelle de pertinence]
La \cle{querelle de pertinence} désigne le déplacement du débat philosophique africain : il ne s'agit plus de demander si une philosophie africaine \emph{existe} (querelle de légitimité), mais de déterminer ce qu'elle \emph{doit faire}, à quels problèmes concrets des sociétés africaines elle doit répondre (développement, justice, gouvernance, éthique).
\end{definition}

Ce déplacement repose sur un raisonnement simple. La question \og la philosophie africaine existe-t-elle ? \fg{} présupposait que l'Occident détenait le droit d'accorder ou de refuser le titre de \og philosophie \fg{}. Or, accepter de répondre à cette question, c'est déjà accepter le tribunal de l'autre. La nouvelle génération refuse donc ce tribunal et pose ses propres questions : comment penser la liberté, l'État, le travail, la technique, à partir des réalités africaines ?

\begin{aretenir}
La réorientation actuelle de la philosophie africaine consiste à passer de la \cle{légitimité} (prouver qu'elle existe) à la \cle{pertinence} (servir à transformer la société). La philosophie devient un outil critique et pratique, non un certificat d'humanité à décrocher devant l'Occident.
\end{aretenir}

\courssub{Le tournant décolonial : déconstruire l'« invention de l'Afrique »}

\textbf{Intuition.} Quand quelqu'un raconte ta propre histoire à ta place, il finit par te faire voir toi-même à travers ses yeux. Pour te libérer, il ne suffit pas de protester : il faut reprendre la parole et raconter ton histoire depuis ton propre point de vue. C'est exactement le geste décolonial.

\begin{definition}[Décolonialité]
La \cle{décolonialité} (Walter Mignolo, Achille Mbembe) est une démarche critique visant à délier la pensée des catégories imposées par la \cle{colonialité} du pouvoir, du savoir et de l'être, c'est-à-dire la survivance des schémas coloniaux après la fin politique de la colonisation. Elle invite à penser à partir de soi, depuis son propre lieu d'énonciation (\emph{locus of enunciation}).
\end{definition}

L'ouvrage fondamental de V. Y. Mudimbe, \emph{L'Invention de l'Afrique} (1988), montre que \og l'Afrique \fg{} telle qu'on la décrit (continent primitif, traditionnel, sans écriture ni histoire) est en grande partie une \emph{construction} du discours occidental : récits des explorateurs, anthropologie coloniale, manuels scolaires. Mudimbe appelle \cle{bibliothèque coloniale} (\emph{colonial library}) cet ensemble de textes et de savoirs qui ont fabriqué une image de l'Afrique pour mieux la dominer. Achille Mbembe (\emph{De la postcolonie}, 2000) prolonge cette analyse en montrant comment le pouvoir postcolonial lui-même reproduit parfois les gestes du pouvoir colonial. Sabelo Ndlovu-Gatsheni parle de \cle{colonialité de l'être} : même libérés politiquement, les Africains peuvent rester intérieurement colonisés, méprisant leurs langues, leurs savoirs et leurs institutions.

Face à cela, le tournant décolonial défend une \cle{épistémologie du Sud} : les savoirs produits depuis le Sud (Afrique, Amérique latine, Asie) ne sont pas des savoirs inférieurs ou locaux, mais des points de vue légitimes sur le monde, capables d'universalité.

\begin{exemplebox}[La carte scolaire comme exemple de bibliothèque coloniale]
Dans beaucoup de manuels, l'histoire de l'Afrique commence avec l'arrivée des Européens. Un élève peut alors croire que \og avant, il n'y avait rien \fg{}. La critique décoloniale montre que ce découpage est un choix idéologique : les royaumes du Sahel, le royaume Kongo, l'Égypte antique ou les civilisations du Grassfield camerounais existaient bien avant la colonisation. Repenser les programmes scolaires est donc un acte philosophique décolonial.
\end{exemplebox}

\begin{attention}[Ne pas confondre décolonisation et décolonialité]
La décolonisation (années 1960) est un événement politique : l'indépendance des États. La \cle{colonialité}, elle, survit dans les mentalités, les langues d'enseignement, les modèles économiques. La décolonialité est donc un travail de longue durée, intellectuel et culturel, qui continue aujourd'hui.
\end{attention}

\courssub{Philosophie du développement et de la libération : Eboussi Boulaga et Jean-Marc Ela}

\textbf{Intuition.} Un médecin qui se contente de décrire la maladie sans la soigner ne sert à rien. De même, une philosophie qui décrit l'Afrique sans chercher à transformer ses maux (pauvreté, corruption, injustice) manque sa vocation. C'est la thèse des philosophes camerounais de la libération.

Fabien Eboussi Boulaga (\emph{La Crise du Muntu}, 1977 ; \emph{Christianisme sans fétiche}, 1981) soutient que la philosophie africaine doit être une \cle{critique de l'idéologie} : elle doit démasquer les discours qui justifient la domination, qu'ils viennent de l'étranger ou des élites locales. Jean-Marc Ela (\emph{Ma foi d'Africain}, 1985 ; \emph{Quand l'État pénètre en brousse}, 1990) va plus loin : la philosophie doit partir des \emph{conditions réelles de vie} des paysans et des pauvres, et devenir une \cle{conscience critique} du développement. Pour Ela, penser, c'est refuser que le paysan camerounais soit réduit au silence et à la misère.

\begin{propriete}[Thèse de la philosophie engagée]
Pour Eboussi Boulaga et Jean-Marc Ela, la philosophie africaine n'a de sens que comme \cle{outil de transformation sociale} : elle doit combattre la corruption, la dictature et la pauvreté, en articulant la réflexion critique à la pratique (la \emph{praxis}). Une philosophie coupée des luttes du peuple est une philosophie trahie.
\end{propriete}

Le schéma suivant résume l'enchaînement logique de cette réorientation : partir de la déconstruction décoloniale pour aboutir à une philosophie de la libération.

\begin{popfigure}
\begin{center}
\begin{tikzpicture}[node distance=0.55cm,
box/.style={draw=popblue, fill=popblueL, rounded corners=3pt, text width=4.6cm, align=center, font=\small, minimum height=1.1cm},
arr/.style={-{Stealth[length=3mm]}, thick, poporange}]
\node[box, align=center] (a){\textbf{Décolonialité}\\ déconstruire la bibliothèque coloniale};
\node[box, below=of a, align=center] (b){\textbf{Critique de l'État}\\ démasquer l'idéologie et la corruption};
\node[box, below=of b, align=center] (c){\textbf{Éthique du développement}\\ penser les fins : justice, dignité};
\node[box, below=of c, align=center] (d){\textbf{Philosophie de la libération}\\ transformation sociale (Ela, Eboussi Boulaga)};
\draw[arr] (a) -- (b);
\draw[arr] (b) -- (c);
\draw[arr] (c) -- (d);
\end{tikzpicture}
\end{center}
\legende{Chaîne conceptuelle de la réorientation actuelle : la critique décoloniale (1) rend possible la critique de l'État postcolonial (2), qui fonde une éthique du développement (3), aboutissant à une philosophie de la libération (4).}
\end{popfigure}

\courssub{Les concepts opératoires camerounais et africains}

La réorientation actuelle ne se contente pas de critiquer : elle \emph{construit} à partir des ressources propres des cultures africaines. Quatre concepts sont ici essentiels.

\begin{definition}[Palabre]
La \cle{Palabre} (au sens philosophique, chez Eboussi Boulaga) est un espace-temps de parole publique, contradictoire et ritualisée, visant le consensus ou la décision juste. Elle constitue le fondement d'une démocratie africaine authentique, où la parole de chacun compte dans la décision commune.
\end{definition}

\begin{itemize}
\item Le \cle{Njangi} ou \cle{tontine} : association d'épargne et de crédit rotatif, très répandue au Cameroun. Philosophiquement, c'est une \emph{économie solidaire} : la richesse circule, la confiance remplace la garantie bancaire, et chacun est responsable de tous.
\item L'\cle{Ubuntu} (cultures bantu) : \og Je suis parce que nous sommes \fg{}. L'identité de la personne se construit dans la relation aux autres ; l'égoïsme est une défaillance de l'humanité même.
\item Le \cle{Muntu} : chez Eboussi Boulaga, la personne africaine pensée comme être-de-relation, dont la \og crise \fg{} vient de la négation coloniale de sa dignité.
\item \cle{Maât} (Égypte antique) : ordre du monde fondé sur la vérité, la justice, l'harmonie et la réciprocité.
\end{itemize}

\begin{savaistu}[Maât, la plus ancienne éthique connue]
Le concept de \cle{Maât} (Égypte antique, vers 2500 av. J.-C.) est le plus ancien système éthique connu de l'humanité : Vérité, Justice, Harmonie, Réciprocité. Le défunt était jugé en pesant son cœur face à la plume de Maât. Ce concept africain préfigure l'Ubuntu et rappelle que la réflexion éthique est née, entre autres, sur le sol africain, bien avant la Grèce.
\end{savaistu}

\begin{popfigure}
\begin{center}
\begin{tikzpicture}[scale=1,
conc/.style={draw=poppurple, fill=poppurpleL, rounded corners=3pt, font=\small, align=center, text width=2.6cm, minimum height=0.9cm},
val/.style={draw=popgreen, fill=popgreenL, rounded corners=3pt, font=\small, align=center, text width=2.6cm, minimum height=0.9cm},
lnk/.style={-{Stealth[length=2.5mm]}, thick, popmuted}]
\node[conc, align=center] (p) at (0,2.6){\textbf{Palabre}\\ délibération};
\node[conc, align=center] (n) at (4.2,2.6){\textbf{Njangi}\\ économie solidaire};
\node[conc, align=center] (u) at (0,0){\textbf{Ubuntu}\\ relation};
\node[conc, align=center] (m) at (4.2,0){\textbf{Muntu}\\ personne};
\node[val] (r) at (8.4,3.9) {\textbf{Rigueur}};
\node[val] (s) at (8.4,2.6) {\textbf{Solidarité}};
\node[val] (i) at (8.4,1.3) {\textbf{Innovation}};
\node[val] (e) at (8.4,0) {\textbf{Éthique}};
\draw[lnk] (p) -- (r);
\draw[lnk] (n) -- (s);
\draw[lnk] (u) -- (e);
\draw[lnk] (m) -- (e);
\draw[lnk] (p) -- (i);
\node[font=\small\itshape, popdark] at (2.1,3.5) {Concepts africains};
\node[font=\small\itshape, popdark] at (8.4,4.9) {Valeurs techniques};
\end{tikzpicture}
\end{center}
\legende{Carte conceptuelle : les concepts africains (Palabre, Njangi, Ubuntu, Muntu) nourrissent les valeurs du technicien et de l'ingénieur (rigueur, solidarité, innovation, éthique).}
\end{popfigure}

\courssub{Le débat actuel : universalisme concret contre relativisme culturel}

Reste une difficulté majeure : si chaque culture a sa vérité, comment condamner l'injustice, la corruption ou les pratiques néfastes ? Le \cle{relativisme culturel} (\og chaque culture a ses valeurs, aucune n'est supérieure \fg{}) protège contre l'ethnocentrisme, mais il paralyse le jugement moral. À l'inverse, l'universalisme abstrait de type colonial imposait les valeurs occidentales comme \emph{les} valeurs universelles.

Le philosophe ghanéen Kwasi Wiredu propose une voie médiane : l'\cle{universalisme concret}.

\begin{propriete}[Thèse de Wiredu : « l'universalisme n'est pas l'uniformité »]
Selon Wiredu, on peut partager des standards logiques et scientifiques universels (la non-contradiction, la preuve, l'honnêteté intellectuelle) tout en conservant des cultures particulières. La philosophie africaine doit donc produire des concepts \emph{universels} à partir de son \emph{particularisme} : partir de sa culture pour dire quelque chose de valable pour tous.
\end{propriete}

\begin{exemplebox}[La Palabre comme universalisme concret]
La Palabre est une pratique particulière, enracinée dans les villages africains. Mais l'idée qu'une décision juste exige la délibération contradictoire de tous les concernés est une idée \emph{universelle} : elle rejoint les exigences de toute démocratie. La Palabre illustre donc parfaitement la thèse de Wiredu : du particulier africain naît un concept universel.
\end{exemplebox}

\textbf{La question de la langue.} Peut-on philosopher en bassa, en bamiléké, en fulfulde ? Wiredu lui-même réfléchissait en akan avant de traduire en anglais. Philosopher en langues nationales permet de penser avec ses propres catégories et de toucher le peuple ; philosopher en français ou en anglais permet de dialoguer avec le monde. La position la plus féconde est la \emph{double appartenance} : traduire les concepts africains vers les langues officielles, et enrichir celles-ci de notions venues des langues nationales.

\begin{attention}[Deux pièges symétriques à éviter]
Éviter le \cle{culturalisme} (\og tout est bon dans la tradition \fg{} : certaines pratiques traditionnelles sont injustes et doivent être critiquées) et le \cle{modernisme aveugle} (\og tout vient de l'Occident \fg{} : c'est retomber dans la colonialité). La réorientation actuelle prône un \emph{dialogue critique des cultures} : on évalue chaque héritage, africain ou occidental, à l'épreuve de la raison et de la justice.
\end{attention}

\courssub{Enjeu pour le Baccalauréat technique : l'éthique professionnelle}

Pour l'élève de Première technique, cette réorientation n'est pas abstraite : elle fonde l'\cle{éthique professionnelle}. Le technicien ou l'ingénieur camerounais ne travaille pas seulement pour un salaire : il est responsable devant la communauté. L'Ubuntu (\og je suis parce que nous sommes \fg{}) signifie concrètement qu'un pont mal construit, un réseau électrique frauduleux ou un logiciel bâclé trahissent \emph{toute la communauté}. La Palabre inspire la concertation dans les équipes de travail ; le Njangi inspire la mutualisation des moyens dans les entreprises coopératives ; Maât inspire l'honnêteté dans les mesures, les devis et les rapports techniques. La déontologie professionnelle moderne (responsabilité, intégrité, service public) trouve ainsi des racines africaines qui la rendent plus vivante et mieux comprise.

\begin{methode}[Dissertation : « La philosophie est-elle utile au développement ? »]
\begin{enumerate}
\item \textbf{Problématiser} : la philosophie est-elle un luxe pour les pays pauvres, ou une condition de leur développement ?
\item \textbf{Thèse 1} : elle éclaire les \emph{fins} du développement (quelle justice ? quelle dignité ? quelle éthique ?).
\item \textbf{Thèse 2} : elle critique les \emph{moyens} (dénonciation de l'aliénation, de la corruption, des modèles importés sans examen).
\item \textbf{Synthèse} : la philosophie est la \emph{conscience critique} du développement (Jean-Marc Ela) : sans elle, le développement est aveugle ; avec elle, il devient libération.
\end{enumerate}
\end{methode}

\begin{exoresolu}[Analyse de sujet : « Peut-on développer l'Afrique sans philosopher ? »]
\textbf{Énoncé.} Un camarade affirme : \og L'Afrique a besoin de routes et d'hôpitaux, pas de philosophie. \fg{} Construisez une réponse argumentée en trois étapes, en mobilisant la réorientation actuelle de la philosophie africaine.
\tcblower
\textbf{Solution détaillée.}
\begin{enumerate}
\item \textbf{Dégager le présupposé.} L'affirmation oppose la matérialité (routes, hôpitaux) à la réflexion, comme si penser était inutile. Or toute construction suppose des choix : quelles routes, pour qui, financées comment ? Ces choix sont déjà philosophiques (justice, priorité, bien commun).
\item \textbf{Mobiliser les auteurs.} Jean-Marc Ela montre qu'un développement sans conscience critique reproduit la domination : des projets imposés de l'extérieur, détournés par la corruption, échouent. Eboussi Boulaga ajoute que la critique de l'idéologie démasque les discours qui justifient ces échecs. Wiredu rappelle que les standards techniques sont universels, mais que leurs finalités doivent être pensées depuis l'Afrique.
\item \textbf{Conclure.} On ne peut pas développer l'Afrique sans philosopher, car développer, c'est d'abord décider \emph{pour quoi} et \emph{pour qui} l'on développe. La philosophie n'est pas un luxe : elle est la condition d'un développement juste et durable. Le technicien lui-même, par son éthique professionnelle (Ubuntu, Maât), est un philosophe en acte.
\end{enumerate}
\end{exoresolu}

\begin{exercice}[Pour s'entraîner]
Sujet de dissertation : \og La tradition africaine est-elle un obstacle ou une ressource pour la modernité technique ? \fg{} Rédigez l'introduction (accroche, problématique, annonce du plan) en veillant à éviter les deux pièges du culturalisme et du modernisme aveugle.
\end{exercice}

\begin{corrige}[Exercice]
Éléments de correction. \textbf{Accroche} : citer une pratique concrète (la tontine qui finance un atelier, la palabre qui règle un conflit de chantier). \textbf{Définition des termes} : tradition (héritage culturel), modernité technique (transformation rationnelle du monde). \textbf{Problématique} : faut-il choisir entre la tradition et la technique, ou la tradition peut-elle fournir les valeurs (solidarité, délibération, responsabilité) qui rendent la technique humaine ? \textbf{Annonce du plan} : (1) la tradition peut être un obstacle quand elle est figée (culturalisme) ; (2) elle est une ressource quand elle inspire l'éthique (Ubuntu, Njangi, Palabre) ; (3) synthèse : le dialogue critique des cultures (Wiredu) permet une modernité technique enracinée et responsable.
\end{corrige}

\coursec{Atelier pratique : Analyse de documents et construction d'outils heuristiques (TP)}

\courssub{Transition depuis la réorientation actuelle}

Après avoir exploré la \textbf{réorientation actuelle} du débat philosophique africain — décolonialité, développement endogène, pertinence sociale — nous passons maintenant à l'\emph{opérationnalisation} de ces acquis. La philosophie ne saurait rester une contemplation abstraite ; pour le technicien que vous êtes en devenir, elle doit devenir un \textbf{outil heuristique} : un instrument de pensée qui structure l'analyse de situations concrètes, guide la décision technique et éclaire l'éthique professionnelle. Cet atelier propose cinq dispositifs progressifs : de l'analyse textuelle rigoureuse à la construction de votre propre définition opératoire, en passant par l'enquête de terrain simulée, la modélisation graphique et le débat argumenté.

\begin{methode}[Démarche globale de l'atelier]
\begin{enumerate}
\item \textbf{Analyser} : déconstruire des textes fondateurs pour en extraire la structure argumentative.
\item \textbf{Enquêter} : appliquer la méthode de la \cle{philosophie sagace} (Oruka) à des proverbes camerounais.
\item \textbf{Modéliser} : organiser les connaissances en carte mentale et en croquis heuristique (arbre conceptuel).
\item \textbf{Débattre} : confronter les thèses dans un cadre réglementé (débat contradictoire encadré).
\item \textbf{Synthétiser} : produire une fiche personnelle, utilisable en situation professionnelle et à l'examen.
\end{enumerate}
\end{methode}

\courssub{TP 1 : Corpus de textes fondateurs — Analyse comparative}

\begin{exemplebox}[Corpus fourni aux élèves (extraits courts)]
\textbf{Texte A — Hegel / Lévy-Bruhl (Thèses ethnocentristes).} \\
\textit{Hegel, \emph{Leçons sur la philosophie de l'histoire} (cours professés dans les années 1820--1830, publiés en 1837) :} « Le Nègre représente l'homme naturel dans toute sa sauvagerie et son indocilité [...] L'Afrique n'est pas une partie historique du monde ; elle n'a ni mouvement ni développement à exhiber. » \\
\textit{Lévy-Bruhl, \emph{Les Fonctions mentales dans les sociétés inférieures} (1910) :} « La mentalité primitive ne connaît pas la contradiction [...] Elle participe mystiquement aux êtres et aux choses. »

\textbf{Texte B — Tempels / Kagamé (Ethnophilosophie / Philosophie bantoue).} \\
\textit{Tempels, \emph{La Philosophie bantoue} (1945) :} « L'être, pour le Bantou, c'est la force [...] La hiérarchie des forces est la loi fondamentale de l'univers. » \\
\textit{Kagamé, \emph{La Philosophie bantu-rwandaise de l'être} (1956) :} « Le concept d'être, en langue rwandaise, implique toujours une relation : être, c'est être-en-relation, être-force agissante. »

\textbf{Texte C — Hountondji / Wiredu / Eboussi Boulaga (Philosophie critique).} \\
\textit{Hountondji, \emph{Sur la « philosophie africaine »} (1976) :} « L'ethnophilosophie fige la pensée africaine dans un système clos, anonyme, intemporel. La vraie philosophie est critique, individuelle, écrite, universelle. » \\
\textit{Wiredu, \emph{Philosophy and an African Culture} (1980) :} « Il faut décoloniser la pensée africaine : soumettre les concepts hérités de la colonisation à l'examen critique, et distinguer ce qui est universel de ce qui n'est que particulier à une culture. » \\
\textit{Eboussi Boulaga, \emph{La Crise du Muntu} (1977) :} « La philosophie africaine doit naître de l'analyse des contradictions réelles de nos sociétés : pouvoir, dépossession, aliénation culturelle. »
\end{exemplebox}

\begin{methode}[Grille d'analyse d'un texte philosophique (3 étapes)]
\begin{enumerate}
\item \textbf{Identifier la thèse centrale} : en une phrase, que soutient l'auteur ?
\item \textbf{Repérer le vocabulaire clé} : concepts techniques, néologismes, termes polysémiques (ex. : \emph{force}, \emph{participation}, \cle{ethnophilosophie}, \cle{décolonisation épistémique}).
\item \textbf{Reconstituer l'argument d'autorité} : sur quoi l'auteur s'appuie-t-il ? (tradition orale, observation ethnologique, logique formelle, expérience historique, engagement politique).
\end{enumerate}
\end{methode}

\begin{exoresolu}[Analyse guidée du Texte A (Hegel)]
\textbf{Thèse :} L'Afrique est hors de l'histoire universelle ; l'homme africain incarnerait l'état de nature pré-logique, pré-étatique, dépourvu de rationalité dialectique. \\
\textbf{Vocabulaire clé :} « sauvagerie », « partie historique », « mouvement », « développement » (au sens hégélien : procès dialectique de l'Esprit). \\
\textbf{Argument d'autorité :} La \emph{philosophie de l'histoire} comme science universelle qui juge les peuples à l'aune de la réalisation de l'Esprit. L'ethnologie naissante (Lévy-Bruhl) viendra plus tard « confirmer » ce jugement par une prétendue scientificité (mentalité pré-logique, loi de la participation mystique). \\
\textbf{Critique interne :} Confusion entre \emph{absence d'histoire écrite} et \emph{absence d'historicité} ; projection de catégories européennes (État, écriture, dialectique) érigées en normes universelles.
\tcblower
\textbf{Solution détaillée :} Ce texte illustre la \cle{thèse ethnocentriste} fondatrice : le déni de rationalité et d'historicité à l'Afrique. Hegel ne nie pas l'existence de sociétés africaines, mais leur \emph{pertinence philosophique}. Lévy-Bruhl transpose ce jugement en psychologie collective : la « mentalité primitive » ne respecterait pas le principe de non-contradiction. Signalons que Lévy-Bruhl lui-même, dans ses \emph{Carnets} posthumes (1949), reviendra partiellement sur cette thèse — preuve que la science progresse par autocritique. Ces thèses ont légitimé la colonisation comme « mission civilisatrice » (apporter l'histoire, la raison, l'État). En TP, vous devez montrer comment ce double discours (philosophique + pseudo-scientifique) structure encore certains préjugés contemporains présentant l'« oralité » ou la « tradition » comme obstacles au développement technique.
\end{exoresolu}

\begin{exercice}[À faire individuellement (15 min)]
Appliquez la grille aux Textes B et C. Pour chaque texte : thèse, 3 concepts clés, argument d'autorité. Puis comparez : en quoi Tempels et Kagamé répondent-ils à Hegel ? En quoi Hountondji, Wiredu et Eboussi Boulaga rompent-ils avec Tempels ?
\end{exercice}

\begin{attention}[Pièges fréquents en analyse de corpus]
\begin{itemize}
\item Ne pas confondre \textbf{résumé} et \textbf{analyse} : ne retracez pas le texte, déconstruisez-le.
\item Ne pas projeter vos opinions : restez sur \emph{ce que le texte dit} et \emph{comment il le dit}.
\item Attention à l'anachronisme : Tempels (1945) écrit \emph{avant} les indépendances ; Hountondji (années 1970) écrit \emph{après}. Le contexte change la portée de l'argument d'autorité.
\item « Argument d'autorité » ne signifie pas « citation d'un grand nom » : c'est la \emph{source de légitimité} invoquée (tradition, science, logique, vécu).
\end{itemize}
\end{attention}

\courssub{TP 2 : Enquête « Philosophie sagace » simulée — Proverbes camerounais}

\begin{definition}[Philosophie sagace (Henry Odera Oruka)]
La \cle{philosophie sagace} désigne la pensée rationnelle, critique et argumentée de \emph{sages} identifiés (et non anonymes), capables de justifier leurs positions par un raisonnement explicite, et distincts du « sens commun » ou des proverbes figés. Elle vise à prouver l'existence d'une philosophie africaine \emph{individuelle, rationnelle et contemporaine}, non réductible à l'ethnophilosophie.
\end{definition}

\begin{methode}[Analyse de proverbe — 5 étapes (Philosophie sagace)]
\begin{enumerate}
\item \textbf{Traduire littéralement} : mot à mot, sans interprétation.
\item \textbf{Identifier les acteurs et les termes} : qui agit ? quels objets ? quelles relations ? (sujet, prédicat, connecteurs logiques).
\item \textbf{Reconstituer le raisonnement implicite} : prémisse majeure (règle générale) + prémisse mineure (fait particulier) $\rightarrow$ conclusion (norme ou valeur).
\item \textbf{Extraire la norme éthique ou épistémologique} : quelle valeur technique, sociale ou cognitive le proverbe institue-t-il ?
\item \textbf{Tester l'universalité} : la structure logique tient-elle dans un autre contexte culturel ? Si oui, c'est une \cle{structure anthropologique invariante} ; si non, c'est une \cle{norme culturellement située}.
\end{enumerate}
\end{methode}

\begin{popfigure}
\begin{center}
\begin{tabularx}{\linewidth}{|l|X|X|X|}
\hline
\rowcolor{popblueL}
\textbf{Proverbe (aire culturelle)} & \textbf{Traduction littérale} & \textbf{Structure logique} & \textbf{Valeur technique / éthique} \\
\hline
\textbf{Ouest (bamiléké)} : « Un seul doigt ne lave pas le visage » & Un doigt unique ne peut nettoyer une face entière. & Majeure : nettoyer une surface complexe requiert une action coordonnée. Mineure : un doigt seul ne peut coordonner. Conclusion : la coopération est nécessaire. & \textbf{Division du travail, travail d'équipe} : en maintenance, un technicien seul ne peut diagnostiquer un système complexe (électrique + hydraulique + informatique). \\
\hline
\textbf{Littoral (bassa)} : « Le crocodile ne mord pas le crocodile » & Entre membres d'un même groupe, pas d'agression. & Majeure : les membres d'un groupe partagent une vulnérabilité commune. Mineure : l'agression interne affaiblit le groupe. Conclusion : solidarité intra-groupe obligatoire. & \textbf{Éthique professionnelle, corps de métier} : les techniciens d'une même filière ne se sabotent pas ; ils partagent des normes de sécurité. \\
\hline
\textbf{Littoral (douala)} : « Le pays ne se construit pas tout seul » & L'œuvre commune exige l'apport de chacun. & Majeure : toute construction collective requiert l'apport de chaque partie. Mineure : l'inaction d'un membre retarde l'œuvre. Conclusion : engagement citoyen requis. & \textbf{Responsabilité technique, service public} : l'ingénieur qui refuse de signer un rapport de conformité retarde la mise en service d'un hôpital. \\
\hline
\textbf{Nord (fulfuldé)} : « La patience est le lait de la réussite » & La réussite se nourrit de patience. & Majeure : les processus complexes ont une temporalité propre. Mineure : la précipitation brise la chaîne causale. Conclusion : respecter les temps techniques. & \textbf{Qualité, respect des procédés} : en BTP, respecter le temps de prise du béton ; en agroalimentaire, respecter la fermentation. \\
\hline
\textbf{Centre-Sud (beti)} : « L'oreille entend pour les deux » & L'écoute profite à celui qui parle et à celui qui écoute. & Majeure : l'information circule mal si un seul pôle fonctionne. Mineure : l'écoute active permet la compréhension. Conclusion : communication bidirectionnelle requise. & \textbf{Communication technique, sécurité} : en maintenance, confirmer l'ordre reçu (boucle de rétroaction) évite l'erreur humaine. \\
\hline
\end{tabularx}
\end{center}
\legende{Grille d'analyse de proverbes camerounais — structure logique et valeur technique. Les formulations des proverbes varient selon les localités : l'important est la méthode d'analyse, non la version retenue. À reproduire au tableau ou sur papier A3.}
\end{popfigure}

\begin{exoresolu}[Application complète : « Un seul doigt ne lave pas le visage »]
\textbf{1. Traduction littérale :} un doigt, seul, ne peut pas laver le visage. \\
\textbf{2. Acteurs et termes :} sujet = « un doigt » (partie) ; objet = « le visage » (totalité complexe) ; verbe = « laver » (action d'entretien) ; négation = « ne peut pas ». \\
\textbf{3. Raisonnement reconstitué :} \\
\quad P1 (majeure) : \emph{pour nettoyer une surface aux reliefs multiples (le visage), une action à plusieurs points de contact simultanés est requise.} \\
\quad P2 (mineure) : \emph{un seul doigt ne fournit qu'un point de contact.} \\
\quad C (conclusion) : \emph{donc un seul doigt ne peut pas laver le visage ; la coopération (les cinq doigts, la main) est nécessaire.} \\
\textbf{4. Valeur technique et éthique :} \textbf{division du travail, interdépendance des compétences, travail d'équipe.} En atelier de maintenance industrielle : le mécanicien, l'électricien, l'automaticien et le soudeur sont les « doigts » ; la machine est le « visage ». L'absence d'un « doigt » (une compétence) laisse une zone « non lavée » (une défaillance non détectée). \\
\textbf{5. Test d'universalité :} la structure logique est valide partout (couverture d'un ensemble par plusieurs éléments). Différence culturelle : la tradition occidentale insiste souvent sur l'\emph{individu} qui coordonne (le chef de projet) ; ici, la \emph{main} (le collectif) est l'agent premier.
\tcblower
\textbf{Solution détaillée :} Ce proverbe n'est pas une simple image : il encode une \textbf{théorie de l'action collective} applicable à l'ingénierie des systèmes. La « main » = le système technique ; les « doigts » = les sous-systèmes et les opérateurs ; la « saleté » = la panne, le défaut ; laver = la maintenance corrective ou préventive. Le proverbe affirme qu'aucun sous-système isolé ne peut assurer la maintenance du système global. C'est le principe de la \textbf{maintenance pluridisciplinaire} et de la \textbf{gestion des interfaces}.
\end{exoresolu}

\begin{exercice}[En binôme (20 min)]
Choisissez un proverbe de votre langue maternelle (ou de votre quartier). Appliquez la méthode en 5 étapes. Présentez au tableau : raisonnement reconstitué + valeur technique + test d'universalité.
\end{exercice}

\begin{attention}[Erreur fréquente : proverbe et parole de sage]
\textbf{Proverbe} = phrase figée, anonyme, transmissible, souvent métaphorique, \emph{non argumentée} en elle-même. Ex. : « Un seul doigt ne lave pas le visage. » \\
\textbf{Parole de sage (philosophie sagace)} = parole \emph{attribuée} à une personne identifiée, \emph{contextualisée}, \emph{argumentée} (le sage \emph{explique} pourquoi). Ex. : « Mon père, le forgeron, disait : "Quand je forge une houe, je frappe le fer chaud à trois endroits précis, car si je ne frappe qu'un seul, la lame se tord. C'est comme pour gouverner : il faut toucher à l'économie, à la justice et à l'éducation ensemble." » \\
\emph{La philosophie sagace vise la parole de sage, pas le proverbe seul.} Le proverbe est la \emph{trace} ; la parole argumentée du sage est l'\emph{acte philosophique}.
\end{attention}

\courssub{TP 3 : Carte mentale géante — « Philosophie africaine » (papier A3 / tableau)}

\begin{definition}[Carte mentale (mind map)]
Représentation graphique \emph{schématique, non réaliste}, qui organise les connaissances sous forme \emph{radiale} (centre $\rightarrow$ branches $\rightarrow$ sous-branches) pour faire apparaître la \textbf{structure logique} d'un concept complexe, ses liens transversaux et ses zones de tension. Outil de \textbf{révision}, de \textbf{plan de dissertation} et de \textbf{communication technique}.
\end{definition}

\begin{popfigure}
\centering
\begin{tikzpicture}[
  mindmap,
  concept color=popblue,
  level 1 concept/.append style={level distance=4.5cm, sibling angle=72},
  level 2 concept/.append style={level distance=3.2cm, sibling angle=45},
  every node/.style={font=\small},
]
\node[concept, text width=3.2cm, align=center] {Philosophie\\ africaine}
  child[concept color=popblue!30] { node[concept, text width=2.8cm, align=center] {1. Ethno-\\philosophie}
    child { node[concept, text width=2.2cm] {Tempels (force vitale)} }
    child { node[concept, text width=2.2cm] {Kagamé (catégories de l'être)} }
    child { node[concept, text width=2.2cm] {Mbiti (temps et religion)} }
    child { node[concept, text width=2.2cm] {Critique : figement, anonymat} }
  }
  child[concept color=poporange!30] { node[concept, text width=2.8cm, align=center] {2. Philosophie\\ nationaliste-\\idéologique}
    child { node[concept, text width=2.2cm] {Nkrumah (consciencisme)} }
    child { node[concept, text width=2.2cm] {Nyerere (ujamaa)} }
    child { node[concept, text width=2.2cm] {Senghor (négritude)} }
    child { node[concept, text width=2.2cm] {Critique : instrumentalisation politique} }
  }
  child[concept color=popgreen!30] { node[concept, text width=2.8cm, align=center] {3. Philosophie\\ critique}
    child { node[concept, text width=2.2cm] {Hountondji (critique de l'ethnophilosophie)} }
    child { node[concept, text width=2.2cm] {Wiredu (décolonisation conceptuelle)} }
    child { node[concept, text width=2.2cm] {Eboussi Boulaga (contradictions réelles)} }
    child { node[concept, text width=2.2cm] {Mudimbe (invention de l'Afrique)} }
  }
  child[concept color=poppurple!30] { node[concept, text width=2.8cm, align=center] {4. Philosophie\\ sagace\\ (Oruka)}
    child { node[concept, text width=2.2cm] {Sages identifiés} }
    child { node[concept, text width=2.2cm] {Rationalité explicite} }
    child { node[concept, text width=2.2cm] {Proverbe $\neq$ parole de sage} }
    child { node[concept, text width=2.2cm] {Critique : représentativité ?} }
  }
  child[concept color=poppink!30] { node[concept, text width=2.8cm, align=center] {5. Réorientation\\ actuelle}
    child { node[concept, text width=2.2cm] {Décolonialité (Mbembe)} }
    child { node[concept, text width=2.2cm] {Développement endogène} }
    child { node[concept, text width=2.2cm] {Éthique publique, gouvernance} }
    child { node[concept, text width=2.2cm] {Pertinence sociale} }
  };
\end{tikzpicture}
\legende{Modèle de carte mentale — structure radiale à 5 branches principales. Les élèves reproduisent cette structure sur A3 en ajoutant : auteurs (nom + date), concepts clés (définition en une ligne), critiques (un argument), actualité camerounaise (ex. : bilinguisme officiel, débats sur la gouvernance, innovation numérique).}
\end{popfigure}

\begin{methode}[Construction de la carte mentale (consignes pour l'élève)]
\begin{enumerate}
\item \textbf{Centre} : « Philosophie africaine » (couleur vive, écriture grande).
\item \textbf{5 branches principales} : les 4 tendances fondatrices + la réorientation actuelle (couleurs distinctes).
\item \textbf{Sous-branches niveau 1} : auteurs majeurs (nom + œuvre clé + date).
\item \textbf{Sous-branches niveau 2} : concepts opératoires (définition courte).
\item \textbf{Sous-branches niveau 3} : critiques internes ou externes (une phrase).
\item \textbf{Sous-branches niveau 4} : \textbf{actualité camerounaise} — relier chaque tendance à un fait concret : loi, institution, débat public, innovation technique.
\item \textbf{Liens transversaux} : traits pointillés entre branches (ex. : Wiredu $\leftrightarrow$ Oruka sur le langage ; Eboussi Boulaga $\leftrightarrow$ Mbembe sur le pouvoir).
\end{enumerate}
\end{methode}

\begin{aretenir}[La carte mentale comme plan de dissertation]
Chaque branche principale = une \textbf{partie} de dissertation. Les sous-branches = les \textbf{arguments} (pour et contre). Les liens transversaux = les \textbf{transitions} et la \textbf{problématisation}. Exemple de sujet : « La philosophie africaine est-elle une philosophie de la libération ? » $\rightarrow$ Partie 1 : ethnophilosophie et philosophie nationaliste = libération identitaire ; Partie 2 : philosophie critique et philosophie sagace = libération épistémologique ; Partie 3 : réorientation actuelle = libération politique et économique.
\end{aretenir}

\courssub{TP 4 : Croquis heuristique — « L'arbre de la philosophie africaine au Cameroun »}

\begin{definition}[Croquis heuristique]
Représentation graphique schématique (non réaliste) qui organise les connaissances sous forme d'\emph{arbre, de réseau ou de flux} pour faire apparaître la structure logique d'un concept complexe. Ici : métaphore organique (racines $\rightarrow$ tronc $\rightarrow$ branches $\rightarrow$ fruits) pour montrer la \textbf{généalogie}, la \textbf{continuité} et la \textbf{fécondité} de la philosophie africaine \emph{au Cameroun spécifiquement}.
\end{definition}

\begin{popfigure}
\centering
\begin{tikzpicture}[
  scale=0.85,
  every node/.style={font=\small, align=center},
  root/.style={draw=popdark, thick, fill=popgold!25, rounded corners=4pt, minimum width=3cm, minimum height=1.2cm},
  trunk/.style={draw=popdark, thick, fill=poporange!20, rounded corners=4pt, minimum width=4cm, minimum height=1.5cm},
  branch/.style={draw=popblue, thick, fill=popblueL, rounded corners=3pt, minimum width=3.5cm, minimum height=1cm},
  fruit/.style={draw=popgreen, thick, fill=popgreenL, circle, minimum size=1.8cm},
  arrow/.style={->, thick, popline},
]
% Racines
\node[root, align=center] (rac1) at (-5, -3){Traditions orales\\ (proverbes, mythes,\\ initiations, sagesses)};
\node[root, align=center] (rac2) at (0, -3.5){Résistances\\ (anticoloniales,\\ culturelles, linguistiques)};
\node[root, align=center] (rac3) at (5, -3){Terroirs et\\ cosmologies\\ (bamiléké, bassa,\\ beti, peul, douala...)};

% Tronc
\node[trunk, align=center] (tronc) at (0, 1){Universités camerounaises\\ (années 1970--1990)\\ Eboussi Boulaga, Owona,\\ Mveng, Ngoa...};

% Branches
\node[branch, align=center] (br1) at (-6, 5){Éthique publique\\ et gouvernance};
\node[branch, align=center] (br2) at (-2, 5.5){Épistémologie\\ et décolonialité};
\node[branch, align=center] (br3) at (2, 5.5){Esthétique\\ et arts vivants};
\node[branch, align=center] (br4) at (6, 5){Philosophie des\\ sciences et techniques};

% Fruits
\node[fruit, align=center] (fr1) at (-7, 8){Politiques\\ publiques};
\node[fruit, align=center] (fr2) at (-3.5, 8.5){Éducation\\ (programmes,\\ manuels)};
\node[fruit, align=center] (fr3) at (0, 9){Résolution\\ des conflits\\ (médiation)};
\node[fruit, align=center] (fr4) at (3.5, 8.5){Innovation\\ frugale\\ (low-tech)};
\node[fruit, align=center] (fr5) at (7, 8){Souveraineté\\ numérique\\ (données, IA)};

% Flèches racines -> tronc
\draw[arrow] (rac1.north) -- (tronc.south west);
\draw[arrow] (rac2.north) -- (tronc.south);
\draw[arrow] (rac3.north) -- (tronc.south east);

% Flèches tronc -> branches
\draw[arrow] (tronc.north west) -- (br1.south);
\draw[arrow] (tronc.north) -- (br2.south);
\draw[arrow] (tronc.north) -- (br3.south);
\draw[arrow] (tronc.north east) -- (br4.south);

% Flèches branches -> fruits
\draw[arrow] (br1.north) -- (fr1.south);
\draw[arrow] (br1.north east) -- (fr2.south west);
\draw[arrow] (br2.north) -- (fr2.north);
\draw[arrow] (br2.north east) -- (fr3.north west);
\draw[arrow] (br3.north east) -- (fr3.north east);
\draw[arrow] (br4.north west) -- (fr4.south west);
\draw[arrow] (br4.north) -- (fr5.south);

% Légende intégrée
\node[draw=popdark, fill=popcream, rounded corners=4pt, inner sep=6pt, font=\footnotesize, align=left] at (10, -2) {
\textbf{Légende :}\\
\textcolor{popgold}{■} Racines = sources vives (orales, résistances)\\
\textcolor{poporange}{■} Tronc = institutionnalisation (université)\\
\textcolor{popblue}{■} Branches = champs disciplinaires actuels\\
\textcolor{popgreen}{■} Fruits = retombées sociétales concrètes
};
\end{tikzpicture}
\legende{Croquis heuristique : l'arbre de la philosophie africaine au Cameroun. À reproduire en grand format (A2/A1) en groupe. Chaque élève ajoute un fruit concret et vérifiable (ex. : une loi, un projet, une association, un mémoire de fin d'études).}
\end{popfigure}

\begin{exemplebox}[Lecture technique de l'arbre — Exemples concrets camerounais]
\textbf{Racine « Résistances » $\rightarrow$ Branche « Éthique publique » $\rightarrow$ Fruit « Politiques publiques » :} la \emph{résistance} des mouvements nationalistes camerounais (dont l'UPC et Ruben Um Nyobè) a nourri une \emph{éthique de la responsabilité} (Eboussi Boulaga) qui inspire aujourd'hui les débats sur la \textbf{déontologie de la fonction publique} et la lutte contre la corruption (déclaration des biens et avoirs, CONAC).
\\
\textbf{Racine « Traditions orales » $\rightarrow$ Branche « Philosophie des sciences et techniques » $\rightarrow$ Fruit « Innovation frugale » :} le \emph{savoir-faire} des forgerons (traitement thermique empirique du fer) $\rightarrow$ analyse épistémologique (Wiredu : rationalité des savoirs locaux) $\rightarrow$ \textbf{fab labs et ateliers d'innovation} de Douala et Yaoundé qui hybrident techniques numériques et savoirs locaux (ex. : séchoirs solaires pour le cacao co-conçus avec les producteurs).
\\
\textbf{Racine « Cosmologies beti » $\rightarrow$ Branche « Esthétique » $\rightarrow$ Fruit « Résolution des conflits » :} le \emph{mvet} (épopée chantée) comme espace de parole ritualisée $\rightarrow$ philosophie de la \emph{parole partagée} $\rightarrow$ pratiques de \textbf{médiation communautaire} et de conciliation traditionnelle reconnues dans le règlement des litiges fonciers et successoraux.
\end{exemplebox}

\begin{exercice}[Groupe de 4 (30 min + 10 min de restitution)]
Sur papier A2, dessinez l'arbre. Chaque membre prend une branche. Cherchez \textbf{2 fruits vérifiables} (loi, projet, association, article de presse, mémoire de fin d'études) par branche. Collez des post-its « fruits » sur les branches. Présentez : « Cette branche porte ce fruit parce que... »
\end{exercice}

\courssub{TP 5 : Débat structuré — Débat contradictoire encadré}

\begin{methode}[Débat contradictoire encadré (inspiré de la méthode critique de Karl Popper)]
\begin{enumerate}
\item \textbf{Sujet} : « La philosophie africaine doit-elle s'écrire obligatoirement en langues africaines pour être authentique ? »
\item \textbf{Équipes} : 3 élèves \emph{Pour} (affirmative), 3 élèves \emph{Contre} (négative), 1 \emph{jury} (3 élèves), 1 \emph{modérateur} (le professeur).
\item \textbf{Temps} : 5 min de préparation en équipe $\rightarrow$ 3 min d'exposé initial (Pour) $\rightarrow$ 3 min d'exposé initial (Contre) $\rightarrow$ 8 min de \emph{questions croisées} (questions directes, pas de discours) $\rightarrow$ 2 min de conclusion (Pour) $\rightarrow$ 2 min de conclusion (Contre) $\rightarrow$ 5 min de délibération du jury + verdict argumenté.
\item \textbf{Règles} : pas d'interruption ; arguments étayés (auteur, concept, fait) ; pas d'argument d'autorité nu (« Senghor a dit » $\rightarrow$ « Senghor soutient que... parce que... »). L'esprit de la méthode de Popper : une thèse n'est recevable que si elle accepte d'être \emph{mise à l'épreuve} et discutée.
\end{enumerate}
\end{methode}

\begin{exemplebox}[Arguments types attendus — Pour préparer les équipes]
\textbf{POUR (l'authenticité passe par la langue) :}
\begin{itemize}
\item Wiredu : les catégories grammaticales (sujet/objet, temps, genre) structurent la pensée. Les langues indo-européennes imposent une ontologie « substance/attribut » étrangère à bien des langues africaines (processuelles, relationnelles).
\item Ngũgĩ wa Thiong'o : « décoloniser l'esprit » passe par l'écriture dans sa propre langue (le gikuyu pour lui). La langue est porteuse d'une \emph{vision du monde}.
\item Exemple technique : dans plusieurs langues camerounaises, le verbe « faire » implique toujours un \emph{patient} et un \emph{résultat} : on ne « fait » pas dans le vide. Cela structure une éthique technique de la \emph{traçabilité} (qui fait quoi, pour quel résultat).
\item UNESCO : droit à l'éducation en langue maternelle (recommandations sur le multilinguisme, 2003).
\end{itemize}

\textbf{CONTRE (l'authenticité passe par la rationalité et l'universalité) :}
\begin{itemize}
\item Hountondji : la philosophie est \emph{universelle} par sa méthode (logique, critique, argumentation), non par sa langue. Écrire en français ou en anglais permet le dialogue critique global.
\item Wiredu lui-même écrit en anglais pour être lu et discuté ; il \emph{traduit} ses concepts akan en les soumettant à la critique universelle.
\item Réalité camerounaise : environ \textbf{250 langues}, dont une minorité seulement dispose d'une orthographe standardisée. Quelle langue choisir ? Risque de \emph{fragmentation} intellectuelle.
\item Exemple technique : le code source, les normes internationales (ISO), les brevets s'écrivent majoritairement en anglais. Un ingénieur camerounais \emph{doit} aussi philosopher en anglais pour participer aux standards mondiaux.
\end{itemize}

\textbf{SYNTHÈSE possible (jury) :} \textbf{bilinguisme épistémologique.} Écrire \emph{les deux} : concepts fondateurs en langue africaine (pour l'ancrage, la nuance, la décolonisation cognitive) + publication en français ou en anglais (pour la circulation, la critique, l'interopérabilité technique). Le défi technique : créer des \textbf{terminologies standardisées} en langues africaines (néologismes, encodage informatique, claviers adaptés).
\end{exemplebox}

\begin{exoresolu}[Simulation de questions croisées (extrait)]
\textbf{Équipe Contre (question à Pour) :} « Vous citez Wiredu. Mais Wiredu écrit en anglais. N'est-ce pas une contradiction ? » \\
\textbf{Équipe Pour (réponse) :} « Wiredu distingue la \emph{langue de pensée} et la \emph{langue de publication}. Il pense en akan et traduit en anglais pour être discuté. Il milite \emph{pour} l'écriture philosophique en akan \emph{aussi}. La contradiction n'est qu'apparente : c'est une stratégie tactique, pas un renoncement. » \\
\textbf{Jury (note) :} distinction opératoire claire (pensée / publication). Argument recevable.
\tcblower
\textbf{Solution détaillée :} Ce débat forme à l'\textbf{argumentation technique} : distinguer le \emph{plan ontologique} (la langue structure-t-elle la pensée ?), le \emph{plan sociologique} (quelle langue pour quelle communauté ?), le \emph{plan technique} (interopérabilité, standards). Le verdict du jury doit être \emph{argumenté}, pas un vote à main levée.
\end{exoresolu}

\begin{exemplebox}[Le défi linguistique en chiffres (Cameroun)]
\begin{itemize}
\item Environ \textbf{250 à 280} langues vivantes recensées (Ethnologue, SIL).
\item Une \textbf{minorité} seulement dispose d'une \emph{orthographe standardisée} (travaux de l'ALCAM, du PROPELCA, du SIL et des universités).
\item \textbf{Aucune} langue africaine n'est langue d'enseignement général à l'université (français et anglais dominent), même si des langues nationales sont enseignées comme matières.
\item Les écritures africaines (bamum, bassa vah, etc.) sont progressivement intégrées à la norme \textbf{Unicode}, mais les \textbf{claviers physiques et les outils de saisie} restent rares pour l'usage courant.
\item \textbf{Enjeu technique :} former des \emph{ingénieurs linguistes} (traitement automatique des langues, corpus, reconnaissance optique, synthèse vocale) pour construire les infrastructures numériques des langues africaines. La philosophie africaine en langues africaines est \emph{aussi} un projet d'ingénierie logicielle.
\end{itemize}
\end{exemplebox}

\begin{savaistu}[Achille Mbembe — Figure majeure contemporaine]
Achille \textbf{Mbembe} (né en 1957 au Cameroun), philosophe et politologue, a enseigné notamment à l'\textbf{université du Witwatersrand} (Johannesburg) et dans plusieurs universités américaines. Dans \emph{Critique de la raison nègre} (2013), il déconstruit le « discours sur le Nègre » (de Hegel aux discours coloniaux, mais aussi certaines réponses de la négritude) comme une production historique. Concept clé : la \textbf{nécropolitique} — le pouvoir souverain contemporain s'exercerait aussi par le \emph{droit de faire mourir ou de laisser mourir} certaines populations (migrants, quartiers précaires, zones de conflit). Son œuvre relie \textbf{philosophie politique, épistémologie décoloniale et analyse du pouvoir d'État}, et illustre la réorientation actuelle du débat philosophique africain vers les problèmes concrets des sociétés.
\end{savaistu}

\courssub{Production attendue : Fiche de synthèse individuelle (1 page A4)}

\begin{propriete}[Fiche « Ma définition opérationnelle de la philosophie africaine pour mon futur métier de technicien »]
\textbf{Format :} 1 page A4, recto seul, police 11 pt, interligne 1,15. \\
\textbf{Structure imposée (5 zones) :}
\begin{enumerate}
\item \textbf{Définition personnelle (3 lignes max)} : « Pour moi, la philosophie africaine, c'est \emph{[verbe d'action]} \emph{[objet]} par \emph{[moyen]} en vue de \emph{[finalité technique/éthique]}. » Ex. : « ... c'est \textbf{interroger les présupposés techniques} de nos projets par \textbf{l'analyse critique des savoirs locaux} en vue de \textbf{concevoir des solutions adaptées, justes et souveraines}. »
\item \textbf{3 auteurs-référents (1 ligne chacun)} : nom + concept clé + pourquoi il guide mon métier.
\item \textbf{1 proverbe ou parole de sage opératoire (analysé au TP2)} : proverbe + raisonnement reconstitué + application technique concrète (mon stage, mon projet de fin d'études).
\item \textbf{1 branche de l'arbre (TP4) où je me situe} : éthique publique / épistémologie / esthétique / philosophie des sciences et techniques + 1 fruit visé.
\item \textbf{1 engagement écrit (signature)} : « Je m'engage à \emph{[action concrète]} dans mes 2 premières années d'exercice. » Ex. : « Intégrer une revue de conception participative (sages, usagers) avant tout cahier des charges. »
\end{enumerate}
\end{propriete}

\begin{aretenir}[Sujet classique Probatoire / Bac technique : « L'Afrique a-t-elle besoin de philosophie ? »]
\textbf{Plan type (3 parties dialectiques) :}
\begin{enumerate}
\item \textbf{Thèse : Non, l'urgence vitale prime.} « Manger, soigner, loger, éduquer » sont des nécessités matérielles ; la philosophie serait un luxe de lettrés. Arguments : l'urgence du développement (Hountondji lui-même insiste sur les besoins concrets), les faibles moyens consacrés à la recherche.
\item \textbf{Antithèse : Oui, pour choisir \emph{comment} manger, soigner, construire.} Une technique sans éthique produit l'aliénation (technologies imposées, dette infrastructurelle, usages prédateurs des données). La philosophie est une \textbf{boussole éthique} (justice, consentement, souveraineté). Auteurs : Eboussi Boulaga (contradictions réelles), Mbembe (critique du pouvoir), les sages (proverbes comme éthique pratique).
\item \textbf{Synthèse : Oui, pour ne pas subir la technique venue d'ailleurs (souveraineté épistémologique).} La philosophie africaine n'est pas « décorative » : elle produit des \textbf{critères d'évaluation technique endogènes} (ex. : « Cette technologie respecte-t-elle le temps de la terre ? » ; « Qui répare quand ça casse ? »). Elle transforme le technicien en \textbf{sujet technique souverain}, et non en simple opérateur de boîtes noires.
\end{enumerate}
\textbf{Accroche conseillée :} partir d'une situation concrète (un chantier, une panne, un appel d'offres) et montrer que la question « faut-il philosopher ? » est déjà une question philosophique.
\end{aretenir}

\begin{attention}[Consignes de rendu pour la fiche de synthèse]
\begin{itemize}
\item \textbf{Interdit :} copier-coller du cours, définitions toutes faites, phrases vagues (« c'est important », « ça aide à réfléchir »).
\item \textbf{Exigé :} verbes d'action (interroger, concevoir, arbitrer, médier, auditer, standardiser), référence précise à \emph{votre} spécialité technique (génie civil, électrotechnique, informatique, maintenance, agroalimentaire, etc.), lisibilité (mots-clés en gras).
\item \textbf{Évaluation :} cette fiche compte dans la \textbf{note continue} du chapitre et peut être versée à votre \textbf{portfolio technique} (utile en soutenance de projet).
\end{itemize}
\end{attention}

\courssub{Bilan de l'atelier — Compétences validées}

\begin{aretenir}[Compétences transversales mobilisées (cadre MINESEC)]
\begin{itemize}
\item \textbf{C1 — Analyser} : déconstruire un corpus, identifier thèse, arguments et source de légitimité (TP1).
\item \textbf{C2 — Enquêter} : appliquer une méthode heuristique (philosophie sagace) à un terrain culturel propre (TP2).
\item \textbf{C3 — Modéliser} : construire des représentations graphiques structurantes (carte mentale, croquis heuristique) pour synthétiser un champ conceptuel complexe (TP3, TP4).
\item \textbf{C4 — Argumenter} : participer à un débat contradictoire réglé, distinguer arguments forts et faibles, produire une synthèse (TP5).
\item \textbf{C5 — Transférer} : produire une définition \emph{opératoire} mobilisable en situation professionnelle réelle (fiche de synthèse).
\end{itemize}
Ces compétences sont \textbf{directement évaluables} à l'écrit (dissertation, commentaire de texte) et à l'oral (soutenance de projet technique, entretien professionnel).
\end{aretenir}

\begin{methode}[Prochaine étape : Intégration (section 5)]
\begin{enumerate}
\item Reprendre votre fiche de synthèse individuelle.
\item La confronter à des \textbf{cas professionnels réels} : appels d'offres, normes environnementales, gestion des données, intelligence artificielle, résolution de conflits sur chantier.
\item Démontrer, cas par cas, que la philosophie africaine n'est pas un « cours de plus », mais un \textbf{levier de compétence technique supérieure}.
\end{enumerate}
\end{methode}

\coursec{Intégration : Éthique technique, gouvernance et philosophie africaine aujourd'hui}

L'atelier pratique précédent nous a permis de construire des outils heuristiques : grilles d'analyse de textes, cartes des tendances, tableaux comparatifs des auteurs. Mais un outil qui reste dans la boîte ne sert à rien. Cette dernière séance répond donc à la question que tout élève de Première technique est en droit de poser : \emph{à quoi sert la philosophie africaine dans ma vie de futur technicien, ingénieur ou entrepreneur ?} Nous allons montrer que la philosophie africaine n'est pas un musée où l'on admire des concepts derrière une vitrine, mais une \cle{boîte à outils conceptuels} pour résoudre des problèmes concrets : la corruption dans les marchés publics, la gouvernance d'une entreprise, le financement d'une startup, la protection des savoirs traditionnels. En un mot : penser le Cameroun de 2035.

\courssub{L'éthique du technicien face à la corruption : Ubuntu contre Kant ?}

Commençons par une situation que tout technicien camerounais rencontrera tôt ou tard. Vous êtes chargé de valider un marché de travaux (construction d'un hangar, installation d'un réseau électrique). Le chef de service vous glisse : « Ferme les yeux sur les malfaçons, signe le procès-verbal de réception, et tu auras ta part. » Que faire ?

Deux grilles d'analyse philosophiques s'offrent à vous.

\textbf{La grille kantienne (éthique déontologique universelle).} Pour Kant, la moralité repose sur le \cle{devoir} et sur un impératif catégorique : « Agis de telle sorte que tu traites l'humanité, en ta personne comme en toute autre, toujours en même temps comme une fin, et jamais simplement comme un moyen. » Accepter le pot-de-vin, c'est traiter les usagers du bâtiment (qui risquent leur vie si l'ouvrage s'effondre) comme de simples moyens au service de mon enrichissement. C'est donc moralement condamnable, quelle que soit la culture.

\textbf{La grille Ubuntu (éthique relationnelle africaine).} L'\cle{Ubuntu} — « je suis parce que nous sommes » — affirme que ma dignité dépend de la tienne. Corrompre, c'est rompre le tissu relationnel qui me constitue comme personne : je trahis la communauté qui m'a formé, je déshonore ma famille, je détruis la confiance sans laquelle aucune vie commune n'est possible. Le corrupteur ne viole pas seulement une règle abstraite : il se \emph{déshumanise} lui-même en coupant le lien qui le fait exister.

\begin{exemplebox}[Deux raisonnements, une même conclusion]
Face au pot-de-vin : Kant dit « tu ne peux pas vouloir que la corruption devienne une loi universelle, car alors plus aucun marché ne serait fiable ». L'Ubuntu dit « en acceptant, tu appauvris la communauté qui t'a fait homme, donc tu t'appauvris toi-même ». Les chemins diffèrent (raison universelle / relation communautaire), mais la conclusion converge : \textbf{refuser}. C'est précisément cette convergence qui fonde une éthique professionnelle contextualisée.
\end{exemplebox}

\begin{popfigure}
\begin{center}
\begin{tikzpicture}[scale=0.95]
\draw[fill=popblue!25,draw=popblue,thick] (-1.6,0) circle (2.6);
\draw[fill=poporange!25,draw=poporange,thick] (1.6,0) circle (2.6);
\node[text=popdark,align=center,font=\bfseries] at (-2.6,1.6) {Éthique\\universelle};
\node[text=popdark,align=center] at (-2.6,0.7) {\small Droits de l'Homme\\ \small Kant, devoir};
\node[text=popdark,align=center,font=\bfseries] at (2.6,1.6) {Éthique\\africaine};
\node[text=popdark,align=center] at (2.6,0.7) {\small Ubuntu, Palabre\\ \small solidarité, parole};
\node[text=popdark,align=center,font=\bfseries\small] at (0,-0.1) {Éthique\\professionnelle\\contextualisée};
\node[text=popdark,align=center] at (0,-1.15) {\small Ingénieur\\ \small camerounais};
\end{tikzpicture}
\end{center}
\legende{Diagramme de Venn : l'éthique professionnelle du technicien camerounais se construit à l'intersection de l'exigence universelle (droits humains, devoir kantien) et des valeurs africaines (Ubuntu, palabre, honneur de la parole donnée).}
\end{popfigure}

\begin{attention}[Ne pas essentialiser « l'Africain »]
Il n'existe pas « un » Africain qui serait naturellement communautaire, ni « un » Occidental naturellement individualiste : ce serait retomber dans les préjugés ethnocentristes dénoncés dans la première section. Le Cameroun lui-même est pluriel : anglophone et francophone, Nord et Sud, urbain et rural, plus de 250 groupes ethniques. Une philosophie africaine sérieuse doit penser cette \cle{pluralité interne} au lieu d'opposer un bloc « Afrique » à un bloc « Occident ».
\end{attention}

\courssub{La Palabre : un modèle de gouvernance d'entreprise et de projet}

La \cle{palabre} est l'institution traditionnelle de délibération publique : sous l'arbre à palabres, tous les concernés par un problème prennent la parole, les anciens écoutent, et l'on cherche un \cle{consensus} plutôt qu'un vote majoritaire qui écrase la minorité.

Transposée dans une entreprise ou un projet technique, la palabre devient un modèle de \cle{gouvernance participative} :
\begin{itemize}
\item \textbf{Inclusion des parties prenantes} : avant de lancer un projet (forage d'un puits, construction d'une route), on consulte tous ceux qui seront affectés — usagers, riverains, techniciens, autorités locales — et pas seulement le financeur.
\item \textbf{Recherche du consensus} : la décision n'est adoptée que lorsqu'aucune objection majeure ne subsiste, ce qui garantit l'adhésion durable de tous.
\item \textbf{Le temps long} : la palabre prend du temps, mais ce temps « perdu » en amont évite les conflits, les sabotages et les abandons de projet en aval.
\end{itemize}

\begin{exemplebox}[Palabre contre efficacité court-termiste]
Un projet d'électrification rurale décidé « depuis Yaoundé » sans consultation peut installer des poteaux sur des terres sacrées : le projet sera boycotté et l'investissement perdu. À l'inverse, une réunion de palabre préalable (deux journées de discussion) permet de déplacer le tracé de 200 mètres et d'obtenir l'engagement du village à protéger les installations. Le modèle occidental court-termiste optimise le \emph{délai de décision} ; le modèle de la palabre optimise la \emph{durabilité de la décision}. Le bon gestionnaire de projet combine les deux.
\end{exemplebox}

\courssub{Le Njangi : une économie sociale fondée sur la parole}

Le \cle{Njangi} (tontine, appelée aussi \emph{esusu} dans les régions anglophones) est une association d'épargne et de crédit rotatif : chaque membre cotise périodiquement une somme fixe, et le total est versé à tour de rôle à l'un des membres. Aucune banque, aucun contrat notarié, aucune garantie matérielle : tout repose sur la \cle{confiance} et la \cle{parole donnée} — un véritable \emph{contrat moral}.

Philosophiquement, le Njangi démontre qu'une institution économique efficace peut reposer sur des valeurs relationnelles (honneur, réciprocité, solidarité) plutôt que sur la seule méfiance institutionnalisée (garanties, hypothèques, contentieux). C'est une forme authentiquement africaine d'\cle{économie sociale et solidaire} (ESS).

\begin{exemplebox}[Financer un projet technique sans banque]
Huit jeunes diplômés d'un lycée technique de Douala créent un Njangi de \num{100000} FCFA mensuels par membre. Chaque mois, un membre reçoit \num{800000} FCFA. Le premier achète un poste à souder et ouvre un atelier ; le deuxième achète un stock de composants électroniques ; le troisième finance le dépôt de sa startup de maintenance informatique. En huit mois, huit micro-entreprises sont lancées, sans aucun crédit bancaire (souvent refusé aux jeunes sans garantie) et sans intérêts usuraires.
\end{exemplebox}

\begin{exemplebox}[Un exemple chiffré qui oblige à philosopher]
Selon la BEAC, le secteur informel — largement porté par la logique Njangi/Esusu — représente environ 90\,\% de l'emploi au Cameroun. Conséquence philosophique : toute réflexion sur l'économie camerounaise qui ignorerait ces pratiques endogènes serait une philosophie hors sol. Penser le développement, c'est d'abord penser \emph{à partir de} ce que les Camerounais font réellement.
\end{exemplebox}

\courssub{Les défis actuels : cerveaux, brevets et algorithmes}

\textbf{La fuite des cerveaux} (\emph{brain drain}). Des milliers de médecins, ingénieurs et techniciens formés au Cameroun — donc financés par la collectivité — partent exercer en Europe ou en Amérique du Nord. Du point de vue de l'Ubuntu, un débat s'ouvre : l'individu a le droit de chercher son épanouissement, mais ne doit-il rien à la communauté qui a payé sa formation ? La question n'a pas de réponse simple : elle exige de concilier liberté individuelle (valeur universelle) et dette envers la communauté (valeur africaine).

\textbf{La biopiraterie.}

\begin{definition}[Biopiraterie]
La \cle{biopiraterie} est l'appropriation illégale ou inéquitable de ressources génétiques et des savoirs traditionnels associés — par exemple le dépôt d'un brevet industriel sur les propriétés du \emph{Prunus africana} (écorce médicinale des montagnes de l'Ouest camerounais) ou de l'Iboga — sans consentement des communautés détentrices du savoir et sans partage des avantages. Le \emph{Protocole de Nagoya} (2010) encadre internationalement l'accès aux ressources génétiques et impose le partage juste et équitable des avantages.
\end{definition}

Le scandale est double : économique (les laboratoires gagnent des milliards, les villages rien) et épistémologique (le savoir traditionnel, validé par des siècles d'usage, n'est reconnu comme « savoir » que lorsqu'un laboratoire occidental le brevète — survivance directe des thèses ethnocentristes étudiées dans la première section).

\textbf{L'intelligence artificielle et les biais algorithmiques.} Les IA sont entraînées sur des données massivement occidentales. Conséquences concrètes : reconnaissance faciale moins fiable sur les visages africains, traduction automatique défaillante pour les langues camerounaises, crédit scoring qui pénalise les profils informels (donc les acteurs du Njangi). Pour le futur technicien africain, l'enjeu est de produire des \cle{données locales} et de former des algorithmes qui ne reproduisent pas l'ethnocentrisme sous forme numérique.

\begin{propriete}[L'éthique de la responsabilité appliquée à l'Afrique]
Hans Jonas formule le principe de l'éthique de la responsabilité : « Agis de sorte que les effets de ton action soient compatibles avec la permanence d'une vie authentiquement humaine sur terre. » Cette maxime résonne remarquablement avec la sagesse africaine selon laquelle \emph{la terre ne nous appartient pas : elle nous est prêtée par nos enfants} (et confiée par les ancêtres). Déforestation, pollution minière, épuisement des sols : le technicien africain est ainsi doublement interpellé — par l'éthique universelle de la responsabilité et par la cosmologie africaine de l'intergénérationnalité. (La formulation de Jonas est à connaître ; sa confrontation avec la pensée des ancêtres est un excellent point d'appui pour une dissertation.)
\end{propriete}

\begin{popfigure}
\begin{center}
\begin{tikzpicture}[
box/.style={draw=popdark,thick,rounded corners=4pt,fill=popblueL,align=center,text width=3.6cm,minimum height=2.6cm,font=\small},
box2/.style={draw=popdark,thick,rounded corners=4pt,fill=poporangeL,align=center,text width=3.6cm,minimum height=2.6cm,font=\small},
box3/.style={draw=popdark,thick,rounded corners=4pt,fill=popgreenL,align=center,text width=3.6cm,minimum height=2.6cm,font=\small},
fleche/.style={-{Stealth[length=3mm]},very thick,popdark}]
\node[box, align=center] (a) at (0,0){\textbf{Savoir endogène}\\[2pt] Plantes médicinales\\ Architecture en banco\\ Njangi / Esusu};
\node[box2, align=center] (b) at (5.6,0){\textbf{Valorisation scientifique}\\[2pt] Recherche locale\\ Brevet, norme ISO\\ Protocole de Nagoya};
\node[box3, align=center] (c) at (11.2,0){\textbf{Innovation technique}\\[2pt] Startups locales\\ Emplois qualifiés\\ Développement endogène};
\draw[fleche] (a) -- (b);
\draw[fleche] (b) -- (c);
\end{tikzpicture}
\end{center}
\legende{Chaîne de valorisation : les savoirs endogènes ne produisent de l'innovation et de l'emploi que s'ils passent par une étape de valorisation scientifique et juridique (brevets, normes, protection contre la biopiraterie).}
\end{popfigure}

\courssub{Méthode : résoudre un cas d'éthique technique}

\begin{methode}[La résolution d'un cas pratique d'éthique technique]
\begin{enumerate}
\item \textbf{Énoncer le dilemme} technique et économique : quels intérêts s'opposent ? (profit / sécurité, délai / qualité, individu / communauté).
\item \textbf{Identifier les valeurs africaines en jeu} : solidarité, honneur de la parole donnée, respect des aînés et de la communauté, Ubuntu.
\item \textbf{Identifier les normes professionnelles} applicables : normes ISO, code de déontologie du corps de métier, lois en vigueur.
\item \textbf{Proposer une issue médiane} par la « palabre technique » : une solution concertée qui respecte à la fois les normes universelles et les valeurs relationnelles locales.
\end{enumerate}
\end{methode}

\begin{exoresolu}[Cas pratique : le forage défectueux]
\textbf{Énoncé.} Une entreprise forée un puits dans un village du Nord. Le technicien camerounais constate que la profondeur contractuelle n'est pas atteinte (nappe non captée correctement), mais le chef de projet lui demande de signer le procès-verbal de réception contre une enveloppe de \num{200000} FCFA, car « le village a soif, on ne peut pas attendre un nouveau financement ». Analyser le cas selon la méthode en quatre étapes.
\tcblower
\textbf{Solution détaillée.}
\begin{enumerate}
\item \textbf{Dilemme} : d'un côté, l'urgence humanitaire (le village a besoin d'eau) et l'intérêt financier personnel ; de l'autre, la qualité technique et la durabilité de l'ouvrage (un forage mal exécuté tarira en saison sèche).
\item \textbf{Valeurs africaines en jeu} : l'Ubuntu — accepter l'enveloppe, c'est trahir la communauté villageoise qui me fait confiance et dégrader ma dignité de personne relationnelle ; l'honneur de la parole donnée — mon engagement professionnel est un contrat moral, comme dans le Njangi.
\item \textbf{Normes professionnelles} : les normes de forage (profondeur, débit, qualité de l'eau) et le code de déontologie interdisent la certification d'un ouvrage non conforme ; Kant interdit de traiter les villageois comme de simples moyens de clôturer le marché.
\item \textbf{Issue médiane (palabre technique)} : refuser de signer, mais ne pas se contenter d'un refus sec. Convoquer une réunion avec le chef de projet, les autorités villageoises et le bailleur : exposer techniquement le problème, proposer un avenant au marché (approfondissement du forage) ou une solution transitoire (adduction provisoire). La vérité dite publiquement — c'est l'esprit de la palabre — protège à la fois le village, l'entreprise et le technicien.
\end{enumerate}
\textbf{Conclusion} : les deux traditions (universelle et africaine) convergent vers le refus de la corruption, et la palabre fournit la procédure concrète pour transformer ce refus en solution constructive.
\end{exoresolu}

\courssub{Perspective : l'« Afrotopie » de Felwine Sarr}

Comment conclure sans tomber ni dans le pessimisme ni dans la naïveté ? Le philosophe sénégalais \cle{Felwine Sarr} propose le concept d'\cle{Afrotopie} (dans \emph{Afrotopia}, 2016) : non pas une utopie irréaliste, mais un \emph{espace réel} — le continent africain lui-même — où l'Afrique pense son avenir \emph{par elle-même}, en mobilisant ses propres ressources culturelles (langues, savoirs, institutions comme la palabre et le Njangi) comme leviers d'innovation technique et sociale, au lieu d'imiter des modèles de développement importés qui ont montré leurs limites.

L'Afrotopie n'est ni le repli identitaire (refuser la science universelle) ni la soumission mimétique (copier l'Occident) : c'est la \emph{synthèse créatrice} — exactement ce que notre diagramme de Venn figurait pour l'éthique.

\begin{savaistu}[La philosophie africaine dans la SND30]
La Stratégie Nationale de Développement 2020-2030 (SND30), feuille de route du Cameroun vers l'émergence en 2035, mentionne explicitement « la promotion des valeurs culturelles » parmi ses leviers de développement. Autrement dit : ce que vous étudiez en philosophie n'est pas un luxe scolaire — c'est le cadre théorique d'une politique nationale. Le Njangi, la palabre, les savoirs endogènes sont officiellement reconnus comme des ressources de développement.
\end{savaistu}

\begin{aretenir}[Synthèse finale de la leçon]
\begin{itemize}
\item La philosophie africaine est née d'une \textbf{contestation} des thèses ethnocentristes (Hegel, Lévy-Bruhl) qui niaient la raison africaine.
\item Elle s'est construite en \textbf{quatre tendances} (ethnophilosophie, philosophie sagace, philosophie nationaliste-idéologique, philosophie critique-professionnelle).
\item Elle se \textbf{réoriente} aujourd'hui vers les problèmes concrets : développement, gouvernance, décolonialité.
\item Elle est enfin une \textbf{boîte à outils} pour le technicien : Ubuntu et Kant pour l'éthique, palabre pour la gouvernance, Njangi pour l'économie solidaire, Afrotopie pour l'avenir.
\end{itemize}
La philosophie africaine n'est pas un musée : c'est un atelier où se fabrique le Cameroun de 2035.
\end{aretenir}

\begin{exercice}[Pour s'entraîner]
\begin{enumerate}
\item Votre startup technique cherche un financement de \num{2000000} FCFA. Comparez deux options — crédit bancaire classique et Njangi — en identifiant pour chacune les valeurs philosophiques sous-jacentes (méfiance institutionnelle vs confiance relationnelle).
\item Sujet de dissertation : « La philosophie africaine peut-elle contribuer au développement technique du Cameroun ? » (Vous mobiliserez la palabre, le Njangi, l'Ubuntu et l'Afrotopie, et vous éviterez le piège de l'essentialisation.)
\end{enumerate}
\end{exercice}

\begin{corrige}[Exercice 1]
\textbf{Éléments de corrigé.}
\begin{enumerate}
\item Le crédit bancaire repose sur une philosophie de la \emph{méfiance institutionnalisée} : garanties matérielles, intérêts, sanctions juridiques ; il suppose que l'individu est un calculateur rationnel qu'il faut contraindre. Le Njangi repose sur une philosophie de la \emph{confiance relationnelle} : la parole donnée vaut contrat, la pression du groupe remplace l'hypothèque, et l'honneur est le capital. Avantages du Njangi : accessibilité aux jeunes sans garantie, absence d'intérêts, renforcement du réseau social. Limites : montants plafonnés, risque de défaut d'un membre, rotation parfois lente. La solution réaliste est souvent hybride : démarrer par le Njangi, puis accéder au crédit formel une fois l'activité lancée — synthèse concrète de l'éthique contextualisée étudiée dans cette section.
\item \textbf{Plan possible} : I. Le développement technique exige plus que des machines : il exige des valeurs (confiance, responsabilité, délibération) — or les modèles importés ignorent les ressources culturelles locales. II. Les concepts africains offrent des outils opérationnels : Ubuntu (éthique anti-corruption), palabre (gouvernance participative des projets), Njangi (financement solidaire), protection des savoirs endogènes contre la biopiraterie. III. Conditions et limites : éviter l'essentialisation, articuler ces valeurs aux normes universelles (ISO, droits humains, Protocole de Nagoya) — c'est le sens de l'Afrotopie de Felwine Sarr et de la SND30.
\end{enumerate}
\end{corrige}

\newpage

\coursec{Activité d'intégration}

\coursec{Leçon 1 : La Philosophie africaine — Atelier d'intégration : Le GIC « Mboa Manioc »}

\courssub{Objectifs de l'activité}
Cette activité d'intégration vise à mobiliser l'ensemble des savoirs et savoir-faire de l'unité « La Philosophie africaine » dans une situation professionnelle authentique. Elle permet à l'élève de :
\begin{itemize}
    \item \textbf{Analyser} des documents hétérogènes (juridiques, proverbiaux, philosophiques, techniques) pour en extraire des principes éthiques et de gouvernance.
    \item \textbf{Argumenter} philosophiquement le bien-fondé d'un projet technique en s'appuyant sur l'anthropologie africaine (\cle{Ubuntu}, \cle{Force vitale}, \cle{Palabre}) et en réfutant les thèses ethnocentristes.
    \item \textbf{Concevoir} un livrable professionnel (Charte Éthique et de Gouvernance) structuré, juridiquement cohérent (droit OHADA), techniquement réaliste et culturellement ancré.
    \item \textbf{Communiquer} et \textbf{défendre} collectivement une proposition lors d'une soutenance orale face à un jury mixte (enseignant, acteur socio-économique).
\end{itemize}

\courssub{Situation}
\textbf{Contexte réel :} Le \textbf{Lycée Technique de Douala-Bassa} (Littoral, Cameroun) s'engage dans un projet d'insertion professionnelle et de développement local. En partenariat avec l'ONG \emph{« Femme et Développement Durable »} (FDD), l'établissement met à disposition un terrain de 2 hectares attenant aux ateliers de Génie Mécanique et de Transformation Agroalimentaire. L'objectif est d'installer une \textbf{unité semi-industrielle de transformation du manioc en farine (gari, foufou, amidon) et en éthanol}, gérée par un \textbf{Groupement d'Initiative Commune (GIC)} composé de 35 femmes transformatrices du quartier Bassa et d'anciennes élèves de la filière Agroalimentaire.

\textbf{Données techniques du projet (Document 4 - Schéma simplifié) :}
\begin{itemize}
    \item \textbf{Capacité :} 5 tonnes de racines fraîches/jour (200 jours/an $\rightarrow$ 1000 t/an).
    \item \textbf{Flux matière :} Racines $\rightarrow$ Lavage $\rightarrow$ Épluchage $\rightarrow$ Broyage $\rightarrow$ Pressage (déshydratation 65\% $\rightarrow$ 35\% matière sèche) $\rightarrow$ Tamisage $\rightarrow$ Torréfaction (Gari) / Séchage (Fécule) $\rightarrow$ Conditionnement.
    \item \textbf{Flux énergie :} Chaudières à biomasse (écorces de manioc + bois de chauffe durable) pour la torréfaction ; groupe électrogène 50 kVA (secours) ; panneaux solaires 10 kWc (éclairage/bureaux).
    \item \textbf{Déchets/Valorisation :} Épluchures (15\% masse) $\rightarrow$ compostage / alimentation porcine ; Jus de pressage (effluents riches en amidon, DBO$_5$ $\approx$ 12 000 mg/L) $\rightarrow$ lagunage naturel 3 bassins $\rightarrow$ irrigation maraîchère ; Gaz de distillation (éthanol) $\rightarrow$ récupération énergie.
    \item \textbf{Investissement initial :} 85 millions FCFA (Subvention ONG 60\%, Apport GIC 10\%, Prêt microfinance 30\%).
    \item \textbf{Emplois directs :} 12 permanents + 25 saisonniers ; \textbf{Chiffre d'affaires prévisionnel an 1 :} 45 millions FCFA.
\end{itemize}

\textbf{Problème éthique et de gouvernance :} La réussite technique ne suffit pas. Des projets similaires au Cameroun ont échoué par : détournement de fonds (corruption fournisseurs/acheteurs), conflits de leadership (absence de règles claires), non-respect environnemental (rejet d'effluents dans le Wouri), marginalisation des membres analphabètes dans les décisions. Le GIC et l'ONG exigent donc, \textbf{avant tout décaissement}, la rédaction d'une \textbf{Charte Éthique et de Gouvernance} signée par tous les membres.

\textbf{Documents ressources fournis aux groupes :}
\begin{enumerate}
    \item \textbf{Doc 1 (Juridique) :} Extrait de l'Acte uniforme OHADA relatif au droit des sociétés coopératives (Arts 10, 22, 45, 67) : principe « un homme = une voix », réserves légales (10\% excédent), responsabilité limitée aux apports, obligations de tenue de comptes et d'assemblée générale annuelle.
    \item \textbf{Doc 2 (Culturel) :} Trois proverbes Duala/Bassa : 
    \begin{itemize}
        \item « \emph{Mboa na mboa, ndolo na ndolo} » (La communauté avant l'individu / L'union fait la force).
        \item « \emph{Nyongo a mwane te mwane a nyongo} » (L'argent va et vient, la confiance reste).
        \item « \emph{Moto a tele bato, a tele mba} » (Qui travaille pour les autres travaille pour lui-même).
    \end{itemize}
    \item \textbf{Doc 3 (Philosophique) :} Extrait de \textbf{Fabien Eboussi Boulaga}, \emph{La Crise du Muntu}, « La Palabre comme espace de vérité et de gouvernance » : « La palabre n'est pas une discussion oisive ; c'est l'institution par laquelle le groupe se constitue comme sujet moral. Elle exige la parole libre, l'écoute de l'adversaire, la recherche du consensus non comme compromis mou, mais comme vérité partagée. Le chef n'est pas celui qui commande, mais celui qui synthétise la volonté commune. »
    \item \textbf{Doc 4 (Technique) :} Schéma flux matière/énergie/déchets (données ci-dessus).
\end{enumerate}

\courssub{Consignes}
Travaillant en groupes de 4 élèves (rôles attribués : \textbf{Ingénieur Procédés}, \textbf{Sociologue/Philosophe}, \textbf{Juriste Droit OHADA}, \textbf{Gestionnaire/Comptable}), vous devez produire le livrable suivant :

\begin{enumerate}
    \item \textbf{Analyse documentaire (30 min) :} Chaque rôle analyse le document correspondant (Ingénieur $\rightarrow$ Doc 4 ; Philosophe $\rightarrow$ Doc 2 \& 3 ; Juriste $\rightarrow$ Doc 1 ; Gestionnaire $\rightarrow$ Doc 1 \& 4 coûts). Remplir une fiche de synthèse : « Principes extraits $\rightarrow$ Application au GIC ».
    \item \textbf{Rédaction de la Charte (1h30) :} Rédiger collectivement la \textbf{Charte Éthique et de Gouvernance du GIC « Mboa Manioc »} (2 pages max, police Times New Roman 12, interligne 1,15). Structure imposée :
    \begin{itemize}
        \item \textbf{Préambule philosophique (1 paragraphe, 150 mots max) :} Justifier le projet par une anthropologie africaine (Ubuntu / Force vitale / Muntu). Refuser explicitement la vision ethnocentriste (Lévy-Bruhl, Hegel) qui nie la rationalité technique et organisationnelle africaine.
        \item \textbf{Titre I : Gouvernance par la Palabre (5 Articles) :}
        \begin{itemize}
            \item Art 1 : Mode de décision (Consensus / Majorité qualifiée ? Rôle du « Chef de Palabre » / Président).
            \item Art 2 : Gestion des conflits (Médiation interne, sages, avant tribunal).
            \item Art 3 : Transparence \& Reddition des comptes (Accès aux livres, affichage mensuel).
            \item Art 4 : Rotation des responsabilités \& Formation continue.
            \item Art 5 : Adhésion / Exclusion (Critères éthiques, non-discrimination genre/âge).
        \end{itemize}
        \item \textbf{Titre II : Responsabilité Environnementale (3 Articles) :}
        \begin{itemize}
            \item Art 6 : Terre des ancêtres / Héritage des enfants (Gestion déchets, lagunage, zéro rejet brut).
            \item Art 7 : Économie circulaire \& Énergie (Valorisation épluchures, biomasse, solaire).
            \item Art 8 : Veille écologique (Comité environnement, indicateurs DBO$_5$, sols).
        \end{itemize}
        \item \textbf{Titre III : Éthique de l'Ingénieur \& Lutte Anti-Corruption (Clause unique, Art 9) :} Refus pots-de-vin fournisseurs/acheteurs ; déclaration d'intérêts ; clause de non-concurrence loyale ; sanction : exclusion + poursuite pénale (Code pénal camerounais, Art 134 ss).
    \end{itemize}
    \item \textbf{Préparation soutenance (30 min) :} Préparer un diaporama 5 slides (Contexte, Philosophie, Gouvernance, Environnement, Business Model Éthique) et désigner 2 orateurs.
    \item \textbf{Soutenance orale (10 min + 5 min questions) :} Devant le jury (Professeur de Philosophie + Chef d'entreprise agroalimentaire / Chef de quartier Bassa). Grille d'évaluation : Pertinence philosophique (4 pts), Faisabilité technique/juridique (4 pts), Ancrage culturel/proverbes (4 pts), Qualité rédactionnelle/orale (4 pts), Esprit d'équipe/réponses (4 pts). Total /20.
\end{enumerate}

\courssub{Ressources à mobiliser}
\begin{itemize}
    \item \textbf{Concepts philosophiques (Leçon) :} \cle{Ubuntu} (« Je suis parce que nous sommes »), \cle{Force vitale} (Tempels, Kagamé), \cle{Muntu} (Eboussi Boulaga), \cle{Palabre} (gouvernance par le verbe), \cle{Ethnocentrisme} (thèses à réfuter : Lévy-Bruhl « pré-logique », Hegel « Afrique sans histoire »), \cle{Décolonialité} (Mudimbe, Wiredu, sortie de l'imitation).
    \item \textbf{Savoir-faire :} Dissertation (thèse/antithèse/synthèse), Commentaire de texte (analyse concepts/arguments), Analyse de documents (croiser sources), Argumentation juridique (droit OHADA), Calculs de base (seuils de rentabilité, ratios déchets).
    \item \textbf{Références légales :} Acte uniforme OHADA Coopératives (2010), Loi n°96/12 du 5 août 1996 (Environnement), Code pénal (Corruption), Code du travail.
\end{itemize}

\begin{exoresolu}[Démarche et corrigé commenté]
\textbf{Étape 1 : Analyse croisée des documents (Exemple rôle Philosophe)}
\begin{itemize}
    \item \textbf{Doc 2 (Proverbes) $\rightarrow$ Principes :} « Mboa na mboa » $\rightarrow$ Primauté du bien commun (Art 1, 5). « Nyongo a mwane... » $\rightarrow$ Confiance $>$ Profit court terme (Art 3, 9). « Moto a tele bato... » $\rightarrow$ Réciprocité, travail décent (Préambule, Art 4).
    \item \textbf{Doc 3 (Eboussi Boulaga) $\rightarrow$ Principes :} Palabre = Institution de vérité (Art 1, 2). Chef = Synthétiseur (Art 1). Parole libre/Écoute (Art 2). Consensus = Vérité partagée (Art 1).
    \item \textbf{Synthèse :} La gouvernance ne s'importe pas (modèle occidental pyramidal), elle s'ancre dans la \cle{Palabre} comme \emph{technique sociale} de cohésion.
\end{itemize}

\textbf{Étape 2 : Rédaction du Préambule (Modèle corrigé)}
\begin{quote}
\textit{« Considérant que l'être humain africain (Muntu) est un être-de-relation, porteur d'une Force vitale qui ne se déploie pleinement que dans la communauté (Ubuntu) ; que la technique n'est pas l'apanage de l'Occident mais une expression de l'intelligence africaine à transformer la nature pour le bien-vivre ensemble ; que les thèses ethnocentristes (Hegel, Lévy-Bruhl) niant notre rationalité technique et éthique sont historiquement falsifiées et ontologiquement violentes ; les membres du GIC "Mboa Manioc" affirment que cette unité de transformation n'est pas une simple usine, mais un \textbf{lieu de vie} où le travail du manioc, don de la terre des ancêtres, devient service de la vie pour les enfants d'aujourd'hui et de demain. La Palabre, notre méthode de gouvernance, garantit que chaque voix compte, que la terre est respectée, et que l'argent sert l'humain, non l'inverse. »}
\end{quote}
\textbf{Commentaire :} Ancre le projet (anthropologie), réfute l'ethnocentrisme (exigence programme), lie technique/vie, annonce la Palabre comme méthode.

\textbf{Étape 3 : Rédaction Articles Gouvernance (Exemples corrigés)}
\begin{itemize}
    \item \textbf{Art 1 (Décision) :} « Les décisions stratégiques (investissements $>$ 500 000 FCFA, admission/exclusion, modification charte) sont prises en Assemblée Générale par \textbf{consensus} (recherche unanimité par palabre). À défaut, majorité qualifiée des 2/3 des membres présents. Le Président (Chef de Palabre) ne vote qu'en cas d'égalité ; son rôle est de formuler la synthèse. » $\rightarrow$ \textit{Justification : Eboussi Boulaga (synthèse) + OHADA Art 22 (AG) + Proverbe "Mboa na mboa".}
    \item \textbf{Art 3 (Transparence) :} « Les comptes (recettes/dépenses, stocks, ventes) sont tenus selon le plan comptable OHADA (SYSCOHADA). Un tableau de bord mensuel (CA, marge, trésorerie, déchets) est affiché au local du GIC et lu en langue duala/bassa lors de la réunion mensuelle. Deux commissaires aux comptes (membres non élus au bureau) vérifient trimestriellement. » $\rightarrow$ \textit{Justification : OHADA Art 45, 67 + "Nyongo a mwane" (confiance par preuve).}
    \item \textbf{Art 9 (Anti-corruption Ingénieur) :} « Tout membre, notamment l'Ingénieur Procédés et le Gestionnaire, s'interdit de solliciter ou accepter tout avantage (espèces, nature, emploi parentèle) d'un fournisseur (intrants, emballages, maintenance) ou acheteur (grossistes, industries). Tout conflit d'intérêt est déclaré par écrit au Bureau. La violation entraîne l'exclusion immédiate (Art 5) et la saisine du Procureur (Code pénal Art 134, 134-1). Un registre des cadeaux/hospitalités (seuil 10 000 FCFA) est tenu. » $\rightarrow$ \textit{Justification : Éthique ingénieur (déontologie) + Droit pénal camerounais + "Moto a tele bato" (intérêt collectif).}
\end{itemize}

\textbf{Étape 4 : Articles Environnement (Exemples corrigés)}
\begin{itemize}
    \item \textbf{Art 6 :} « La terre est un dépôt sacré des ancêtres et un prêt des enfants. Les effluents (jus de pressage) sont traités par le système de lagunage 3 bassins (rétention 21 jours, DBO$_5$ sortie $<$ 100 mg/L) avant toute irrigation. Aucun rejet direct dans le cours d'eau (Wouri). Analyse mensuelle par laboratoire agréé (MINEE). » $\rightarrow$ \textit{Justification : Force vitale (terre vivante) + Loi 96/12 Art 17 (pollution) + Données techniques Doc 4.}
    \item \textbf{Art 7 :} « L'unité vise l'autonomie énergétique 80\% à 3 ans. Chaudières 100\% biomasse (écorces + bois certifié reboisement GIC). Installation solaire 10 kWc an 1 $\rightarrow$ 30 kWc an 3. Épluchures : 50\% compost (retour sol), 50\% alimentation porcs (convention éleveurs locaux). » $\rightarrow$ \textit{Justification : Économie circulaire + Décolonialité (souveraineté énergétique) + Données techniques.}
\end{itemize}

\textbf{Étape 5 : Soutenance (Grille jury - Extraits attendus)}
\begin{itemize}
    \item \textbf{Question Jury (Chef traditionnel) :} « Comment la Palabre gère-t-elle une urgence technique (panne chaudière) sans attendre l'AG ? »
    \item \textbf{Réponse type (Philosophe/Ingénieur) :} « La Charte (Art 1 al. 3) prévoit un "Conseil des Sages Opérationnels" (Président, Ingénieur, Trésorière, Doyenne) décideur d'urgence (plafond 200 000 FCFA), redevable à la palabre mensuelle. C'est la Palabre adaptée à la contrainte technique : réactivité + reddition. »
    \item \textbf{Question Jury (Entrepreneur) :} « Votre marge nette prévisionnelle est de 12\%. Comment financez-vous le lagunage (coût 8 millions) sans augmenter le prix aux productrices ? »
    \item \textbf{Réponse type (Gestionnaire) :} « Subvention ONG (environnement) 4M + Fonds de réserve OHADA (10\% excédent an 1 $\approx$ 0,5M) + Prêt vert microfinance (taux bonifié 5\%) 3,5M. Prix productrices garanti par contrat-cadre 3 ans. » $\rightarrow$ \textit{Mobilisation : Juridique (réserves), Technique (coûts), Éthique (juste prix).}
\end{itemize}
\end{exoresolu}

\courssub{Bilan de l'activité}
Cette activité d'intégration a permis de vérifier l'appropriation effective des trois piliers de la leçon :
\begin{enumerate}
    \item \textbf{La réfutation opératoire de l'ethnocentrisme :} Les élèves n'ont pas seulement cité Lévy-Bruhl ou Hegel ; ils ont \textbf{produit} une rationalité technique (schéma flux, calculs) et éthique (Charte) africaine, prouvant par l'acte même que la philosophie africaine est une \emph{philosophie de l'action} (Eboussi Boulaga, Hountondji).
    \item \textbf{L'articulation des quatre tendances :} 
    \begin{itemize}
        \item \textit{Ethnophilosophie} (Proverbes, Ubuntu, Force vitale $\rightarrow$ Préambule, Art 6).
        \item \textit{Nationaliste/Idéologique} (Souvenirs lutte, développement endogène $\rightarrow$ Souveraineté énergétique Art 7, GIC local).
        \item \textit{Professionnelle/Académique} (Rigueur conceptuelle, Eboussi Boulaga, Wiredu $\rightarrow$ Structure Palabre Art 1-2, distinction consensus/compromis).
        \item \textit{Critique/Décoloniale} (Sortie de l'imitation, pertinence sociale $\rightarrow$ Charte hybride OHADA/Proverbes, clause anti-corruption adaptée au réel camerounais).
    \end{itemize}
    \item \textbf{La réorientation actuelle (Décolonialité/Développement) :} Le livrable n'est pas un exercice scolaire clos, mais un \textbf{outil juridique et managérial réel}, utilisable immédiatement par le GIC. L'évaluation par un jury mixte (académique + terrain) sanctionne la \cle{pertinence sociale}, critère suprême de la philosophie africaine contemporaine (Mudimbe, Mbembe).
    \item \textbf{Transversalité technique :} L'ingénieur n'est pas un simple exécutant ; l'Art 9 fait de lui un \textbf{garant éthique}. La gestion des déchets (Art 6-7) mobilise les compétences de Génie Chimique/Procédés. La Charte devient ainsi le \textbf{contrat social} liant technique, éthique, droit et culture.
\end{enumerate}
\textbf{Compétences certifiées (Probatoire) :} Rédiger un texte argumenté structuré (Dissertation/Commentaire) ; Analyser un corpus documentaire ; Mobiliser des connaissances philosophiques pour éclairer une situation concrète ; Travailler en équipe ; Communiquer oralement avec pertinence.

\begin{attention}[Pièges à éviter]
\begin{enumerate}
\item Confondre \emph{consensus} (recherche d'un accord unanime par la parole) et \emph{compromis} (renoncements réciproques). La Palabre vise le consensus comme vérité partagée (Eboussi Boulaga).
\item Oublier d'ancrer les articles juridiques (OHADA) dans les proverbes et la philosophie (Ubuntu, Palabre) : la Charte doit être un texte \emph{hybride}, pas un copier-coller du droit uniforme.
\item Négliger la contrainte technique réelle (coût lagunage, DBO$_5$, capacité chaudière) dans les articles environnementaux : l'éthique sans technique est vaine.
\item Réduire la réfutation de l'ethnocentrisme à une citation d'auteurs sans montrer en quoi le projet \emph{incarne} la rationalité africaine (organisation, calcul, éthique).
\end{enumerate}
\end{attention}

\begin{savaistu}[Le GIC au Cameroun : cadre juridique et ancrage culturel]
Le \textbf{Groupement d'Initiative Commune (GIC)} est une forme juridique souple (Loi n°92/006 du 14 août 1992), très utilisée par les femmes rurales et les jeunes diplômés pour mutualiser des moyens de production. Contrairement à la coopérative (OHADA), le GIC n'a pas de capital minimum ni de réserve légale obligatoire, mais il peut adhérer à l'Acte uniforme coopératif pour sécuriser ses fonds. Le choix du GIC « Mboa Manioc » reflète une stratégie \emph{décoloniale} : s'approprier le droit moderne (OHADA pour la comptabilité, la transparence) tout en légitimant la gouvernance par le droit coutumier (Palabre, proverbes, autorité des aînés). C'est une illustration concrète de la \cle{décolonialité} prônée par Mudimbe et Wiredu : « penser par soi-même » en hybridant les normes.
\end{savaistu}

\begin{methode}[Réussir la soutenance orale (5 slides, 10 min)]
\begin{enumerate}
    \item \textbf{Slide 1 - Contexte \& Problématique :} Localisation (Douala-Bassa), Partenaires (Lycée, ONG FDD), Objectif (Unité manioc 5t/j), Enjeu (Charte avant décaissement).
    \item \textbf{Slide 2 - Fondement Philosophique :} Anthropologie (Muntu, Force vitale, Ubuntu), Réfutation ethnocentrisme (Hegel/Lévy-Bruhl), Méthode : Palabre (Eboussi Boulaga) + Proverbes Duala.
    \item \textbf{Slide 3 - Gouvernance Opérationnelle :} Organigramme (AG, Bureau, Conseil Sages Opérationnels), Règles décision (Consensus/2/3), Transparence (SYSCOHADA, affichage bilingue), Anti-corruption (Art 9, registre cadeaux).
    \item \textbf{Slide 4 - Responsabilité Environnementale :} Lagunage 3 bassins (normes MINEE), Zéro rejet Wouri, Économie circulaire (épluchures $\rightarrow$ compost/porcs, biomasse $\rightarrow$ énergie), Objectif autonomie 80\% (solaire + biomasse).
    \item \textbf{Slide 5 - Business Model Éthique \& Durable :} CA 45M FCFA, Marge 12\%, Financement lagunage (Subvention + Réserves + Prêt vert), Prix productrices garanti 3 ans, Emplois (12 perm + 25 sais), Impact social (femmes, jeunes).
\end{enumerate}
\textbf{Astuce :} Les 2 orateurs se partagent : Philosophe/Juriste (Slides 1, 2, 3) et Ingénieur/Gestionnaire (Slides 3, 4, 5). Prévoir une fiche « Réponses types » aux 5 questions probables du jury.
\end{methode}

\newpage

\coursec{Méthodes et exercices résolus}

\courssub{Méthodes et exercices résolus}

\begin{methode}[Analyser une thèse ethnocentriste occidentale (Hegel, Lévy-Bruhl, Kant)]
\textbf{Objectif :} Démontrer comment un texte philosophe nie la rationalité ou l'historicité à l'Afrique.
\begin{enumerate}
    \item \textbf{Identifier l'auteur et le contexte :} Situer l'œuvre (ex. \emph{La Raison dans l'histoire} et les \emph{Leçons sur la philosophie de l'histoire} de Hegel, \emph{Les fonctions mentales dans les sociétés inférieures} de Lévy-Bruhl) et l'intention (justification coloniale, anthropologie évolutionniste).
    \item \textbf{Repérer les concepts-clés opératoires :} \cle{Esprit subjectif} / \cle{Esprit objectif} (Hegel) ; \cle{Mentalité prélogique} / \cle{Loi de participation} / \cle{Mystique} (Lévy-Bruhl) ; \cle{Caractère national} / \cle{Tempérament} (Kant).
    \item \textbf{Isoler les arguments explicites :} Citer les passages où l'auteur affirme l'absence d'histoire, de logique formelle, de morale universelle ou de technique chez l'Africain.
    \item \textbf{Mettre au jour le présupposé ethnocentrique :} Montrer que le critère d'évaluation est exclusivement occidental (État moderne, logique aristotélicienne, écriture, industrie). C'est un \cle{projectionnisme} : l'Autre est mesuré à l'aune du Même.
    \item \textbf{Conclure sur la portée idéologique :} Expliquer en quoi cette thèse légitime la \cle{mission civilisatrice} et le déni de subjectivité politique.
\end{enumerate}
\end{methode}

\begin{exoresolu}[Analyse d'un extrait de Hegel : « L'Afrique n'est pas une partie historique du monde »]
\textbf{Énoncé :} À partir de l'extrait suivant des \emph{Leçons sur la philosophie de l'histoire} (leçons professées dans les années 1820, publiées en 1837) : \emph{« L'Afrique proprement dite... est le pays de l'or et de l'enfance... Ce que nous comprenons en somme sous le nom d'Afrique, c'est un monde anhistorique non développé, entièrement prisonnier de l'esprit naturel... Le Nègre représente l'homme naturel dans toute sa sauvagerie et son absence de frein... »}, analysez la thèse ethnocentrique en suivant la méthode ci-dessus.
\tcblower
\textbf{Solution détaillée :}
\begin{enumerate}
    \item \textbf{Auteur/Contexte :} Georg Wilhelm Friedrich Hegel (1770-1831), philosophe allemand de l'idéalisme absolu. Ce texte provient de ses leçons sur la philosophie de l'histoire (publiées après sa mort, en 1837). Contexte : essor de la colonisation européenne, diffusion des théories hiérarchisant les races, besoin idéologique de justifier la traite négrière puis la conquête coloniale (qui culminera avec la Conférence de Berlin, 1884-1885).
    \item \textbf{Concepts-clés :} \cle{Esprit universel} (Weltgeist), \cle{Esprit subjectif} (individuel), \cle{Esprit objectif} (État, Droit, Morale), \cle{Histoire universelle} comprise comme progrès de la conscience de la liberté. L'Afrique est exclue du mouvement dialectique de l'Esprit.
    \item \textbf{Arguments explicites :}
        \begin{itemize}
            \item Isolement géographique : l'Afrique serait « repliée sur elle-même », séparée du mouvement de l'histoire universelle (déterminisme géographique).
            \item Négation historique : « monde anhistorique non développé », l'Afrique « n'est pas une partie historique du monde ».
            \item Essentialisation anthropologique : « l'homme naturel », « sauvagerie », « absence de frein », absence supposée de morale, de droit et de religion vraie (réduite au « fétichisme »).
        \end{itemize}
    \item \textbf{Présupposé ethnocentrique :} Hegel prend l'\cle{État moderne occidental} (prussien) comme telos (fin) de l'Histoire. La \cle{raison} est identifiée à la \cle{logique dialectique hégélienne} et à l'\cle{écriture}. L'oralité, la parenté, la gestion consensuelle des conflits (la palabre) sont invisibilisées car non conformes au modèle hégélien. C'est un \cle{critère unique} (l'État bureaucratique) érigé en universalité.
    \item \textbf{Portée idéologique :} Ce discours transforme une \cle{contingence historique} (rapports de force, traite négrière) en \cle{nécessité ontologique} (nature immuable de l'Africain). Il fournit la \cle{légitimation philosophique} à l'entreprise coloniale : « civiliser » celui qui serait « hors de l'histoire ».
\end{enumerate}
\textbf{Conclusion :} La thèse de Hegel n'est pas une observation anthropologique neutre, mais une \cle{construction philosophique} servant à fonder l'altérité radicale de l'Africain pour justifier sa domination.
\end{exoresolu}

\begin{methode}[Identifier et caractériser les quatre tendances de la philosophie africaine moderne]
\textbf{Objectif :} Classer un auteur, un texte ou une thèse dans l'une des quatre familles identifiées par la critique (Oruka, Hountondji, Bodunrin, Mudimbe).
\begin{enumerate}
    \item \textbf{Rappeler les quatre tendances canoniques (taxinomie d'Oruka/Hountondji) :}
        \begin{itemize}
            \item \cle{Ethnophilosophie} : Étude des visions du monde traditionnelles (mythes, proverbes, coutumes) présentées comme une philosophie implicite, collective, anonyme. (Ex. Tempels, Kagame, Mbiti).
            \item \cle{Philosophie nationaliste-idéologique} : Pensée politique des pères des indépendances, mise au service de la libération et de la construction nationale. (Ex. Nkrumah, Nyerere, Senghor, Sékou Touré).
            \item \cle{Philosophie professionnelle (académique)} : Philosophie écrite, critique, argumentée, publiée dans des revues, enseignée à l'université, utilisant les standards logiques universels. (Ex. Hountondji, Wiredu, Bodunrin, Oruka, Mudimbe).
            \item \cle{Philosophie herméneutique} : Interprétation critique des textes et des oralités traditionnelles avec les outils de la phénoménologie et de l'herméneutique (Ricœur, Heidegger) pour en extraire une sagesse opératoire. (Ex. Serequeberhan, Okere, Wamba dia Wamba).
        \end{itemize}
    \item \textbf{Analyser l'objet d'étude :} Est-ce la \cle{tradition orale} (Ethno), le \cle{projet politique} (Nationaliste), l'\cle{argumentation logique} (Professionnelle), l'\cle{interprétation textuelle} (Herméneutique) ?
    \item \textbf{Analyser la méthode :} \cle{Descriptive/collective} (Ethno) vs \cle{Prescriptive/politique} (Nat.) vs \cle{Critique/argumentative individuelle} (Pro) vs \cle{Interprétative/réflexive} (Herm.).
    \item \textbf{Identifier la visée :} \cle{Identité culturelle} (Ethno), \cle{Libération/Développement} (Nat.), \cle{Vérité scientifique/Universalisme critique} (Pro), \cle{Compréhension/Sagesse} (Herm.).
    \item \textbf{Nuancer :} La plupart des auteurs majeurs traversent plusieurs tendances (ex. Wiredu : philosophie professionnelle + critique de l'ethnophilosophie ; Senghor : nationaliste-idéologique + ethnophilosophie).
\end{enumerate}
\end{methode}

\begin{exoresolu}[Classification d'extraits selon les quatre tendances]
\textbf{Énoncé :} Rattachez chacun des extraits suivants à l'une des quatre tendances de la philosophie africaine. Justifiez votre réponse en mobilisant deux critères (objet, méthode, visée).
\begin{enumerate}
    \item \emph{« La philosophie bantoue est une philosophie de la force vitale... La force vitale est l'être même. »} (Placide Tempels, \emph{La Philosophie bantoue}, 1945).
    \item \emph{« Il n'y a pas de philosophie africaine traditionnelle... Il n'y a que des pensées traditionnelles. La philosophie est une activité critique, individuelle et écrite. »} (Paulin Hountondji, \emph{Sur la « philosophie africaine »}, 1976).
    \item \emph{« Le socialisme africain n'est pas une idéologie importée... Il découle de notre structure sociale traditionnelle : l'ujamaa (famille élargie). »} (Julius Nyerere, \emph{Ujamaa. Essais sur le socialisme}, 1962-1968).
    \item \emph{« Il faut lire les proverbes yoruba non comme des vérités figées, mais comme des invitations à un dialogue critique sur l'éthique. »} (approche herméneutique, type Sophie Olúwọlé).
\end{enumerate}
\tcblower
\textbf{Solution détaillée :}
\begin{enumerate}
    \item \textbf{Tendance : Ethnophilosophie.}
        \begin{itemize}
            \item \textit{Objet :} « Philosophie bantoue », « force vitale », concepts traditionnels collectifs.
            \item \textit{Méthode :} Description systématisante d'une vision du monde implicite, anonyme, présentée comme un système cohérent (ontologie de la force).
            \item \textit{Visée :} Prouver l'existence d'une philosophie africaine \emph{avant} l'écriture et la colonisation (identité).
        \end{itemize}
    \item \textbf{Tendance : Philosophie professionnelle (académique).}
        \begin{itemize}
            \item \textit{Objet :} La définition même de la philosophie (critique, individuelle, écrite).
            \item \textit{Méthode :} Argumentation logique, distinction conceptuelle rigoureuse (pensée vs philosophie), réfutation de l'ethnophilosophie (Tempels).
            \item \textit{Visée :} Épistémologique : fonder une philosophie africaine \cle{scientifique} et \cle{universelle} selon les standards académiques.
        \end{itemize}
    \item \textbf{Tendance : Philosophie nationaliste-idéologique.}
        \begin{itemize}
            \item \textit{Objet :} L'organisation politique et économique (socialisme, ujamaa).
            \item \textit{Méthode :} Récupération et réécriture de la tradition (famille élargie) pour légitimer un projet politique moderne (développement, décolonisation).
            \item \textit{Visée :} \cle{Pragmatique} et \cle{émancipatrice} : construire la nation postcoloniale.
        \end{itemize}
    \item \textbf{Tendance : Philosophie herméneutique.}
        \begin{itemize}
            \item \textit{Objet :} Les proverbes (textes/oralités) yoruba.
            \item \textit{Méthode :} \cle{Interprétation} (herméneutique) : ne pas figer le sens (vérité dogmatique) mais ouvrir un \cle{dialogue critique} (horizon de compréhension, fusion des horizons).
            \item \textit{Visée :} Extraire une \cle{sagesse opératoire} pour l'éthique contemporaine, sans essentialiser la tradition.
        \end{itemize}
\end{enumerate}
\end{exoresolu}

\begin{methode}[Construire un plan de dissertation philosophique (sujet type Probatoire/Bac technique)]
\textbf{Objectif :} Répondre à une question philosophique par une démonstration structurée, dialectique et argumentée.
\begin{enumerate}
    \item \textbf{Analyser le sujet (5 min) :}
        \begin{itemize}
            \item Définir les \cle{mots-clés} (ex. « Tradition », « Modernité », « Développement », « Raison », « Communauté »).
            \item Identifier la \cle{tension} ou le \cle{problème} (opposition, paradoxe, évolution).
            \item Formuler la \cle{problématique} (question unique à laquelle le plan répondra).
        \end{itemize}
    \item \textbf{Construire le plan dialectique (3 parties) :}
        \begin{itemize}
            \item \textbf{Thèse (I) :} L'évidence, le sens commun, la thèse adverse que l'on va critiquer. (Ex. « La tradition est un frein à la modernité technique »).
            \item \textbf{Antithèse (II) :} La critique de I, l'apport de la philosophie africaine, la nuance. (Ex. « La tradition est le socle indispensable d'une modernité endogène »).
            \item \textbf{Synthèse (III) :} Le dépassement, la réconciliation conceptuelle, la proposition nouvelle. (Ex. « Pour une \cle{modernité négociée} : l'innovation technique encadrée par l'éthique communautaire »).
        \end{itemize}
    \item \textbf{Sous-parties (A, B, C) :} Une idée directrice par sous-partie = un argument + un exemple/référence auteur + une transition.
    \item \textbf{Rédiger l'introduction (AMI) :} \cle{Amorce} (citation/fait social) $\rightarrow$ \cle{Mise au point} (définitions) $\rightarrow$ \cle{Intérêt/Problématique} $\rightarrow$ \cle{Annonce du plan}.
    \item \textbf{Rédiger la conclusion :} \cle{Bilan} (réponse à la problématique) $\rightarrow$ \cle{Ouverture} (autre champ, actualité technique, autre auteur).
\end{enumerate}
\end{methode}

\begin{exoresolu}[Sujet : « La tradition africaine est-elle un obstacle au développement technique ? »]
\textbf{Énoncé :} Faites une introduction complète (AMI) et un plan détaillé (I, A, B, C ; II, A, B, C ; III, A, B, C) pour ce sujet de dissertation.
\tcblower
\textbf{Solution détaillée :}

\textbf{Introduction (AMI) :}
\emph{Amorce :} « La tradition n'est pas le culte des cendres, mais la conservation du feu » (Gustav Mahler). Au Cameroun, l'ingénieur conçoit des ponts en béton pendant que le chef de village consulte les ancêtres avant les travaux.
\emph{Mise au point :} \cle{Tradition} : ensemble de coutumes, savoirs, valeurs transmis oralement. \cle{Développement technique} : maîtrise de la nature par la science et la technique pour améliorer les conditions de vie. \cle{Obstacle} : frein, blocage, mais aussi résistance nécessaire.
\emph{Problématique :} La tradition africaine, souvent perçue comme statique et irrationnelle (thèse ethnocentrique), ne constitue-t-elle pas au contraire, par son éthique communautaire et sa connaissance du milieu, le \cle{levier culturel} indispensable à un développement technique \cle{soutenable} et \cle{approprié} ?
\emph{Annonce du plan :} Nous montrerons d'abord que la tradition peut apparaître comme un frein épistémologique et social (I), avant de démontrer qu'elle recèle des ressources épistémologiques et éthiques pour une technique endogène (II), pour enfin proposer le concept de \cle{modernité négociée} comme dépassement (III).

\textbf{Plan détaillé :}

\textbf{I. La tradition comme frein épistémologique et structurel au développement technique (Thèse)}
\begin{itemize}
    \item \textbf{A. L'argument ethnocentrique :} L'oralité et le mythe s'opposeraient au \cle{logos} scientifique (écriture, calcul, expérience reproductible). Réf. : Lévy-Bruhl (mentalité prélogique), Hegel (absence d'histoire).
    \item \textbf{B. Le poids du sacré et de l'ancestralité :} La technique moderne exige la \cle{rupture} et l'\cle{innovation} ; la tradition sacralise la \cle{répétition} et l'\cle{immuable}. Ex. : interdits coutumiers compliquant l'exploitation minière ou l'aménagement hydroélectrique (sites sacrés à préserver autour des grands barrages camerounais).
    \item \textbf{C. Le communautarisme face à l'individualisme technique :} L'innovation technique suppose un \cle{sujet autonome} (brevet, profit, responsabilité légale). La tradition valorise le \cle{collectif} et la \cle{redistribution}, ce qui peut freiner l'entrepreneuriat technique individuel. Réf. : Wiredu (critique du consensus autoritaire).
\end{itemize}

\textbf{II. La tradition comme ressource épistémologique et éthique pour une technique endogène (Antithèse)}
\begin{itemize}
    \item \textbf{A. Savoirs endogènes et « science indigène » :} L'agroforesterie (zaï, haies vives), la pharmacopée, l'architecture en banco, la gestion traditionnelle de l'eau sont des \cle{technologies vernaculaires} efficaces, écologiques, adaptées au milieu. Réf. : Paulin Hountondji (savoirs endogènes), Cheikh Anta Diop.
    \item \textbf{B. L'éthique de la \cle{Force Vitale} / de l'\cle{Ubuntu} comme régulateur :} La technique sans éthique mène à l'écocide et à l'aliénation. La tradition impose la \cle{responsabilité intergénérationnelle} (« Nous n'héritons pas de la terre de nos parents, nous l'empruntons à nos enfants »). Critique du « développement » purement quantitatif.
    \item \textbf{C. La \cle{palabre} comme gouvernance technique participative :} La prise de décision technique (implantation d'usine, OGM, IA) gagne à intégrer la \cle{concertation} traditionnelle (palabre, tontines/njangi) pour la légitimité sociale et l'acceptabilité. Réf. : Kwame Gyekye (personne et communauté), Francis Nyamnjoh.
\end{itemize}

\textbf{III. Vers une modernité négociée : l'hybridation critique comme condition du développement réel (Synthèse)}
\begin{itemize}
    \item \textbf{A. Dépasser le binaire Tradition/Modernité :} Concept de \cle{tradition vivante} : la tradition n'est pas un musée, elle s'invente dans la rencontre. La technique n'est pas neutre, elle porte des valeurs.
    \item \textbf{B. L'\cle{appropriation technique} (endogénéisation) :} Former l'ingénieur africain à \cle{penser en langues africaines} les concepts techniques (Wiredu : décolonisation conceptuelle). Ex. : traduire « algorithmique » en langues locales pour en saisir la logique dans un cadre conceptuel propre.
    \item \textbf{C. Proposition concrète : \cle{ingénierie communautaire} :} Cahier des charges technique intégrant : 1. l'impact sur la \cle{Force Vitale} (environnement, sites sacrés) ; 2. le \cle{transfert de compétences} local (formation des jeunes) ; 3. la \cle{gouvernance par la palabre} (comité de suivi villageois/entreprise/État). Modèle : projets d'énergie solaire décentralisée gérés par des coopératives villageoises (Cameroun, Sénégal).
\end{itemize}

\textbf{Conclusion :} La tradition n'est obstacle que si on la fige (ethnophilosophie) ou si on importe la technique sans éthique. Elle devient levier si on la \cle{critique} et l'\cle{hybride}. \emph{Ouverture :} Quel rôle pour l'IA générative dans la préservation et la création de cette tradition vivante ?
\end{exoresolu}

\begin{methode}[Appliquer les concepts éthiques africains (Ubuntu, Force Vitale, Palabre) à une situation technique concrète]
\textbf{Objectif :} Résoudre un dilemme éthique professionnel (ingénieur, technicien, gestionnaire) en mobilisant la philosophie africaine.
\begin{enumerate}
    \item \textbf{Identifier le dilemme technique :} Conflit entre \cle{efficacité/rendement} (norme technique importée) et \cle{valeurs humaines/écologiques} (contexte local).
    \item \textbf{Mobiliser le concept opératoire pertinent :}
        \begin{itemize}
            \item \cle{Ubuntu / Bumuntu} : « Je suis parce que nous sommes » $\rightarrow$ Priorité au \cle{bien commun}, à la \cle{dignité} de la personne dans la relation.
            \item \cle{Force Vitale} (Tempels/Kagame) : Tout être est force. La technique doit \cle{renforcer} la vie, non la détruire (environnement, santé).
            \item \cle{Palabre} : Art du dialogue inclusif, de la \cle{parole partagée}, du consensus non unanime mais accepté. Outil de \cle{gouvernance} de projet.
            \item \cle{Responsabilité intergénérationnelle} : Gestion des ressources pour les ancêtres (passé) et les enfants (futur).
        \end{itemize}
    \item \textbf{Construire l'argumentation en 3 temps :}
        \begin{enumerate}
            \item \textbf{Diagnostic éthique :} En quoi la solution purement technique viole-t-elle le concept africain ? (Ex. déforestation = atteinte à la Force Vitale des ancêtres).
            \item \textbf{Proposition alternative :} Comment le concept réoriente-t-il le cahier des charges ? (Ex. reboisement obligatoire, circuit court, comité de vigilance villageois).
            \item \textbf{Justification de la viabilité :} Montrer que l'alternative est \cle{techniquement réalisable} et \cle{économiquement viable} à long terme (résilience, acceptabilité sociale).
        \end{enumerate}
    \item \textbf{Formuler la conclusion professionnelle :} « En tant qu'ingénieur [spécialité], je décide de... car cela respecte le principe de... »
\end{enumerate}
\end{methode}

\begin{exoresolu}[Cas pratique : Ingénieur Gestion de l'Environnement / Projet minier]
\textbf{Énoncé :} Une multinationale obtient un permis d'exploitation de cobalt dans l'Est-Cameroun. Le projet technique standard prévoit : déforestation de 500 ha, relocalisation de 3 villages sans concertation préalable, rejet de résidus dans la rivière. En tant qu'ingénieur conseil chargé de l'Étude d'Impact Environnemental et Social (EIES), proposez une alternative éthique fondée sur la philosophie africaine.
\tcblower
\textbf{Solution détaillée :}

\textbf{1. Diagnostic éthique (violation des principes africains) :}
\begin{itemize}
    \item \cle{Force Vitale} : La forêt n'est pas une « ressource » inerte mais le siège des forces des ancêtres et des esprits tutélaires (totems). Sa destruction pure et simple est un \cle{sacrilège ontologique}, pas seulement une perte de biodiversité.
    \item \cle{Ubuntu} : « Umuntu ngumuntu ngabantu » (« un homme est un homme par les autres hommes »). Déplacer des villages sans \cle{consentement libre, préalable et éclairé (CLIP)} brise le lien social. L'individu déraciné perd sa dignité (perte des tombeaux, des lieux d'initiation).
    \item \cle{Palabre} : L'absence de concertation vraie remplace le dialogue par l'injonction administrative. Risque de conflit latent (instabilité du projet).
\end{itemize}

\textbf{2. Proposition alternative : le « Protocole Palabre-Cobalt »}
\begin{itemize}
    \item \textbf{A. Phase 1 : la Grande Palabre pré-opérationnelle (gouvernance).} 6 mois de dialogue inclusif : entreprise, État, chefferies, jeunes, femmes, ONG locales. \textbf{Sortie :} accord signé (contrat social) définissant : zone d'exploitation réduite (200 ha), couloirs écologiques et sites sacrés préservés, fonds de développement communautaire géré par un \cle{Conseil des Sages} paritaire.
    \item \textbf{B. Phase 2 : technique « Force Vitale » (environnement).} Remplacer le rejet en rivière par un \cle{procédé hydrométallurgique en circuit fermé} (zéro rejet liquide). Reboisement immédiat avec essences locales (moabi, iroko, plantes médicinales) sur 300 ha compensatoires. \textbf{Indicateur :} taux de survie des plants supérieur à 80\,\% à 2 ans (suivi par le comité villageois).
    \item \textbf{C. Phase 3 : Ubuntu économique (social).} Pas de relocalisation forcée. \cle{Co-activité} : emploi prioritaire des jeunes des villages (formation à la maintenance des engins). Parts sociales (5\,\% des dividendes) versées dans une caisse de type \cle{njangi} villageoise (épargne solidaire) pour la santé et l'éducation.
\end{itemize}

\textbf{3. Justification de la viabilité technique et économique :}
\begin{itemize}
    \item \textbf{Coût initial +15\,\%} (circuit fermé, palabre), mais \textbf{risque opérationnel fortement réduit} (pas de blocage des routes, pas de grève, licence sociale acquise).
    \item \textbf{Conformité :} respecte les \cle{Principes de l'Équateur} (banques), les principes directeurs de l'\cle{OCDE} (devoir de vigilance), le principe \cle{CLIP} (ONU, droits des peuples autochtones).
    \item \textbf{Ingénierie :} solution technique \cle{supérieure} (le circuit fermé permet la récupération des métaux résiduels, donc un profit supplémentaire).
\end{itemize}

\textbf{Conclusion professionnelle :} « En tant qu'ingénieur Environnement, je valide l'alternative "Palabre-Cobalt". Elle transforme la contrainte éthique africaine (Ubuntu, Force Vitale, Palabre) en \cle{performance technique} (circuit fermé) et \cle{sécurité juridique} (licence sociale). C'est l'application concrète de la \cle{décolonialité épistémique} dans l'ingénierie. »
\end{exoresolu}

\begin{methode}[Rédiger une synthèse de documents philosophiques (commentaire composé / explication de texte)]
\textbf{Objectif :} Articuler plusieurs documents (textes, graphiques, données) pour répondre à une question problème.
\begin{enumerate}
    \item \textbf{Lecture active (15 min) :} Pour chaque document : auteur ? date ? thèse principale (une phrase) ? concepts-clés ? ton (polémique, descriptif, programmatique) ?
    \item \textbf{Construire le fil conducteur :} Regrouper les documents par \cle{convergence} (même thèse) ou \cle{divergence} (thèses opposées) par rapport à la question problème.
    \item \textbf{Plan en 2 ou 3 parties thématiques (pas par document !) :}
        \begin{itemize}
            \item I. Le diagnostic partagé / le point de rupture (ex. : le constat de l'aliénation culturelle).
            \item II. Les réponses divergentes / les voies proposées (ex. : retour à la tradition vs invention moderne).
            \item III. La perspective de réorientation actuelle (ex. : décolonialité, épistémologies du Sud).
        \end{itemize}
    \item \textbf{Rédiger en « je » analytique :} « Le document 1 montre que... », « Contrairement à l'auteur 2, l'auteur 3 argue que... », « Cette tension illustre le débat entre... ».
    \item \textbf{Citation courte, analyse longue :} Citer \emph{entre guillemets} (max 1 ligne), puis \textbf{expliquer} le sens, la portée, les limites.
\end{enumerate}
\end{methode}

\begin{exoresolu}[Synthèse : « La réorientation actuelle du débat philosophique africain »]
\textbf{Énoncé :} À partir du corpus suivant, montrez comment le débat philosophique africain est passé de la « preuve d'existence » à la « pertinence sociale ».
\textbf{Doc 1 :} Hountondji (1976) : \emph{Sur la « philosophie africaine »} (critique de l'ethnophilosophie, appel à la rigueur professionnelle).
\textbf{Doc 2 :} Wiredu (1996) : \emph{Cultural Universals and Particulars} (décolonisation conceptuelle, distinction tradition/modernité).
\textbf{Doc 3 :} Ndlovu-Gatsheni (2018) : \emph{Epistemic Freedom in Africa} (décolonialité, épistémologies du Sud, justice cognitive).
\textbf{Doc 4 :} Données (fictives, à titre d'illustration pédagogique) : nombre de revues philosophiques africaines indexées (Scopus/Web of Science) : 2000 : 3 $\rightarrow$ 2023 : 27. Part d'articles co-écrits Nord-Sud : 65\,\%.
\tcblower
\textbf{Solution détaillée :}

\textbf{Introduction :} Le débat a connu trois temps : 1. \cle{Légitimation} (années 1950-1970 : « Y a-t-il une philosophie africaine ? »). 2. \cle{Professionnalisation} (années 1970-1990 : « Comment faire une philosophie rigoureuse ? »). 3. \cle{Décolonialité/Pertinence} (depuis les années 2000 : « Pour qui et pour quoi philosopher ? »). La synthèse montre ce glissement.

\textbf{I. De la défense d'identité à l'exigence de rigueur : la rupture professionnelle (Doc 1 et Doc 2).}
\begin{itemize}
    \item \textbf{Convergence :} Rejet commun de l'\cle{ethnophilosophie} (Tempels) jugée essentialiste, anonyme, non critique.
    \item \textbf{Doc 1 (Hountondji) :} Pose la condition \cle{épistémologique} : « il n'y a de philosophie que professionnelle ». Critique le « mythe » d'une philosophie collective. Visée : \cle{universalisme critique} (standards logiques assumés).
    \item \textbf{Doc 2 (Wiredu) :} Prolonge par la \cle{décolonisation conceptuelle}. Il ne suffit pas d'écrire en anglais ou en français ; il faut \cle{traduire} les concepts (ex. « personne » vs « individu », « vérité » vs « opinion ») pour éviter l'impensé colonial. Visée : \cle{pluralisme épistémique} rigoureux.
\end{itemize}

\textbf{II. De l'académisme à la justice cognitive : le tournant décolonial (Doc 3).}
\begin{itemize}
    \item \textbf{Rupture :} Ndlovu-Gatsheni critique l'universalisme abstrait de Hountondji et Wiredu comme encore trop centré sur les canons occidentaux (revues indexées, langue anglaise).
    \item \textbf{Nouveau paradigme :} \cle{Décolonialité} (Quijano, Mignolo) appliquée à l'Afrique. \cle{Liberté épistémique} : droit de penser \emph{depuis} l'Afrique (langues africaines, oralité, spiritualité) et non plus seulement \emph{sur} l'Afrique.
    \item \textbf{Visée :} \cle{Pertinence sociale}. La philosophie doit répondre aux urgences : pauvreté, changement climatique, gouvernance, IA.
\end{itemize}

\textbf{III. Les indicateurs structurels d'une réorientation inachevée (Doc 4).}
\begin{itemize}
    \item \textbf{Progression quantitative :} multiplication par 9 des revues indexées en 23 ans $\rightarrow$ \cle{institutionnalisation} réussie (validation par les pairs internationaux).
    \item \textbf{Limite structurelle :} 65\,\% de co-écriture Nord-Sud $\rightarrow$ \cle{dépendance épistémique} persistante. Les agendas de recherche, les financements, les standards de publication restent largement définis au Nord.
    \item \textbf{Tension actuelle :} Comment passer de la \cle{reconnaissance académique} (indexation) à la \cle{souveraineté épistémique} (revues en langues africaines, financement endogène, impact politique local) ?
\end{itemize}

\textbf{Conclusion :} Le débat s'est réorienté de la \cle{reconnaissance ontologique} (« Nous existons comme philosophes ») vers la \cle{responsabilité politique} (« Nous pensons pour transformer »). Le défi actuel n'est plus de \emph{prouver} qu'on philosophe, mais de \emph{décider} quelle philosophie faire : une philosophie de « rattrapage » (indexation) ou une philosophie de \cle{libération} (décolonialité, langues africaines, ingénierie sociale).
\end{exoresolu}

\begin{methode}[Rédiger une conclusion argumentée et une ouverture pertinente (savoir-faire : justifier une conclusion)]
\textbf{Objectif :} Clore une dissertation ou un commentaire par une réponse ferme à la problématique et une ouverture qui ne soit pas un « sujet pour une prochaine fois ».
\begin{enumerate}
    \item \textbf{Le Bilan (synthèse dialectique) :} Reprendre la \cle{problématique} et y répondre en une phrase affirmative qui résume le chemin parcouru (Thèse $\rightarrow$ Antithèse $\rightarrow$ Synthèse). \textbf{Interdit :} « En conclusion, nous avons vu que... » (récit de vie). \textbf{Exigé :} « Au terme de cette analyse, il apparaît que... »
    \item \textbf{La portée (enjeu) :} Dire \cle{quoi} cela change concrètement (pour l'ingénieur, le citoyen, la politique, l'épistémologie).
    \item \textbf{L'Ouverture (élargissement contrôlé) :} Trois types validés à l'examen :
        \begin{enumerate}
            \item \textbf{Extension disciplinaire :} même problème vu par une autre discipline (sociologie, droit, économie). Ex. : « Ce débat sur la tradition rejoint la sociologie du changement social... »
            \item \textbf{Actualité technique/technologique :} lien direct avec la spécialité de l'élève (génie civil, informatique, agriculture). Ex. : « Cette éthique ubuntu trouve une application urgente dans la gouvernance de l'IA générative au Cameroun (biais algorithmiques, données locales)... »
            \item \textbf{Auteur/Œuvre non vu :} citer un auteur pertinent non étudié dans le devoir. Ex. : « On pourrait prolonger par Achille Mbembe (\emph{Critique de la raison nègre}) qui montre que... »
        \end{enumerate}
    \item \textbf{Formule de chute :} une phrase courte, mémorable, qui laisse résonner la pensée.
\end{enumerate}
\end{methode}

\begin{exoresolu}[Rédaction d'une conclusion type Bac pour le sujet : « Tradition et Développement technique »]
\textbf{Énoncé :} Rédigez la conclusion complète (Bilan + Portée + Ouverture + Chute) pour le sujet de dissertation traité à l'exercice précédent.
\tcblower
\textbf{Solution détaillée :}

\textbf{Bilan :} Au terme de cette analyse, il apparaît que la tradition africaine n'est un obstacle au développement technique que si l'on la fige dans un passé mythique (ethnophilosophie) ou si l'on importe la technique comme un paquet clé en main sans éthique. Elle devient, au contraire, le \cle{levier indispensable} d'une modernité négociée lorsqu'elle est soumise à l'exigence critique de la philosophie professionnelle (Hountondji, Wiredu) et mobilisée comme ressource éthique (Ubuntu, Force Vitale) et politique (Palabre).

\textbf{Portée :} Pour l'ingénieur technicien camerounais, cela signifie que la \cle{compétence technique} (calcul de structure, code informatique, dimensionnement réseau) est \cle{nécessaire mais insuffisante}. La \cle{compétence éthique} — savoir pour qui, pour quoi, avec quels impacts sur la Force Vitale communautaire — devient le cœur du métier. L'ingénieur n'est plus un simple exécutant de normes, mais un \cle{médiateur culturel} entre la science universelle et les réalités vécues locales.

\textbf{Ouverture (type Actualité technique - IA) :} Cette réorientation trouve une urgence inédite avec l'avènement de l'\cle{Intelligence Artificielle générative}. Si les données d'entraînement (Big Data) ignorent les langues bassa, bamoun, fulfuldé ou les savoirs agricoles oraux, l'IA reproduira le mépris épistémique colonial. La philosophie africaine actuelle nous somme donc de concevoir une \cle{IA « ubuntu-compatible »} : open-source, entraînée sur des corpus locaux, gouvernée par des comités d'éthique locaux (palabre numérique), pour que la technique serve enfin la \cle{Force Vitale} du continent.

\textbf{Chute :} « La technique sans tradition est aveugle ; la tradition sans technique est impuissante. L'avenir de l'Afrique se joue dans leur \cle{traduction mutuelle}. »
\end{exoresolu}

\begin{methode}[Comparer deux conceptions de la rationalité (occidentale vs africaine) pour un commentaire de texte]
\textbf{Objectif :} Structurer une explication de texte comparatif (ex. extrait de Descartes vs extrait de Wiredu/Kagame).
\begin{enumerate}
    \item \textbf{Thèse du texte A (occidentale classique) :} Identifier le modèle : \cle{rationalité calculante} (Descartes), \cle{sujet transcendantal} (Kant), \cle{logique formelle binaire} (Aristote). Critères : universalité abstraite, décontextualisation, objectivité (« vue de nulle part »), écriture, individualisme méthodologique.
    \item \textbf{Thèse du texte B (africaine) :} Identifier le modèle : \cle{rationalité vitale/relationnelle} (Kagame, Tempels), \cle{rationalité communicative/consensuelle} (Wiredu, Gyekye), \cle{sagesse pratique} (sagesse orale). Critères : \cle{contextualisation} (situation concrète), \cle{finalité vitale} (bien-vivre), \cle{oralité/performativité}, \cle{communautarisme épistémique} (la vérité se construit \emph{ensemble}).
    \item \textbf{Tableau comparatif mental (à avoir en tête, pas nécessairement à reproduire dans la copie) :}
        \begin{center}
        \begin{tabular}{|l|l|l|}
        \hline
        \textbf{Critère} & \textbf{Rationalité occidentale (classique)} & \textbf{Rationalité africaine (moderne)} \\ \hline
        \textbf{Sujet} & \cle{Cogito} isolé, autonome & \cle{Nous} relationnel (Ubuntu) \\ \hline
        \textbf{Objet} & Nature objectifiée, mathématisable & \cle{Force Vitale}, entités relationnelles \\ \hline
        \textbf{Méthode} & Doute hyperbolique, déduction & \cle{Palabre}, consensus, proverbes \\ \hline
        \textbf{Vérité} & Adéquation intellect/réalité (statique) & \cle{Convenance} / cohérence vitale (dynamique) \\ \hline
        \textbf{Écriture} & Condition de la philosophie & Support secondaire (mémoire vive) \\ \hline
        \end{tabular}
        \end{center}
    \item \textbf{Problématiser la rencontre :} Ce n'est pas « qui a raison ? » mais « \cle{comment articuler les deux ?} » (complémentarité, hybridation, décolonialité). Éviter le relativisme (« chacun sa vérité ») et l'universalisme abstrait (« la science est une »).
    \item \textbf{Structure du commentaire :} Intro (auteurs, thèses, problématique commune) $\rightarrow$ I. Analyse texte A (logique interne) $\rightarrow$ II. Analyse texte B (logique interne + réponse à A) $\rightarrow$ III. Confrontation/Synthèse (apports réciproques, limites) $\rightarrow$ Conclusion.
\end{enumerate}
\end{methode}

\begin{exoresolu}[Commentaire comparé : Descartes (Discours de la méthode) vs Wiredu (philosophie et culture africaine)]
\textbf{Énoncé :} « Le bon sens est la chose du monde la mieux partagée... » (Descartes) vs « La vérité n'est pas une correspondance mais une cohérence au sein d'un système de croyances partagées... » (Wiredu, résumé). Problématiser : « La rationalité est-elle universelle ou culturellement située ? »
\tcblower
\textbf{Solution détaillée :}

\textbf{Introduction :}
\emph{Amorce :} Deux figures du philosophe : le solitaire dans son poêle (Descartes, 1637) et l'aîné sous l'arbre à palabres (Wiredu, Ghana, 1996).
\emph{Thèses :} Descartes fonde la rationalité sur le \cle{sujet pensant universel} (cogito), méthode du doute, clarté et distinction. Wiredu fonde la rationalité sur la \cle{pratique langagière communautaire}, le consensus rationnel, la décolonisation des concepts.
\emph{Problématique :} La rationalité est-elle une \cle{faculté naturelle universelle} (Descartes) identique en tout homme, ou une \cle{compétence culturelle située} (Wiredu) qui prend des formes historiques distinctes sans pour autant être relativiste ?
\emph{Plan :} 1. L'universalisme abstrait du sujet cartésien. 2. Le particularisme contextualisé de la rationalité akan. 3. Vers une rationalité \cle{pluriverselle} : universalité concrète par la traduction.

\textbf{I. Descartes : l'universalisme du sujet désincarné.}
\begin{itemize}
    \item \textbf{A. Le bon sens comme nature universelle :} « Puissance de bien juger » distribuée équitablement. Présupposé : l'erreur vient du \cle{mauvais usage} (précipitation, prévention), pas de la culture.
    \item \textbf{B. La méthode comme garantie d'objectivité :} Doute méthodique $\rightarrow$ \cle{cogito} (certitude absolue) $\rightarrow$ règles (analyse, synthèse, énumération). La vérité = \cle{clarté et distinction} (intuition intellectuelle).
    \item \textbf{C. Conséquence :} \cle{abstraction totale}. Le sujet épistémique n'a ni sexe, ni race, ni langue, ni histoire. La raison est \cle{une} et \cle{mathématique}. La pensée africaine (orale, vitale) est d'emblée exclue comme « pré-logique » ou « pré-scientifique ».
\end{itemize}

\textbf{II. Wiredu : la rationalité comme compétence culturelle située.}
\begin{itemize}
    \item \textbf{A. Critique de l'universalisme abstrait :} Le « bon sens » n'est pas inné, il s'\cle{apprend} dans une langue, une culture (akan). Les concepts de « vérité », « personne », « causalité » diffèrent structurellement (ex. \cle{nyansa} = sagesse pratique vs \cle{epistèmè} = science théorique).
    \item \textbf{B. La rationalité communicationnelle (palabre) :} La vérité émerge du \cle{dialogue} (argumentation, contre-argumentation) visant le \cle{consensus} (pas l'unanimité forcée). Le sujet épistémique est \cle{relationnel} : « Je pense \emph{avec} les autres ».
    \item \textbf{C. Décolonisation conceptuelle :} Il ne s'agit pas de nier la logique formelle (le modus ponens est valide partout), mais de refuser l'\cle{impérialisme conceptuel} (imposer « individu », « droit subjectif », « atome » comme seuls cadres de pensée). La rationalité africaine est \cle{pratico-vitale} (résolution de conflits, santé, agriculture).
\end{itemize}

\textbf{III. Confrontation : ni relativisme, ni universalisme vide — la rationalité pluriverselle.}
\begin{itemize}
    \item \textbf{A. Complémentarité :} Descartes donne la \cle{rigueur formelle} (mathématiques, physique, code informatique). Wiredu donne la \cle{finalité éthique} (pourquoi cette technique ? pour quel « nous » ?). L'ingénieur a besoin des deux.
    \item \textbf{B. Le pont : la \cle{traduction} (Wiredu) / l'\cle{herméneutique} (Serequeberhan).} Universaliser n'est pas \emph{imposer} un modèle unique, mais \emph{traduire} : montrer comment le concept akan de \cle{sunsum} (esprit/force vitale) éclaire la bioéthique moderne (consentement communautaire vs individuel).
    \item \textbf{C. Enjeu technique :} Un algorithme d'IA « cartésien » (optimisation mathématique pure) peut détruire un écosystème social « wireduien ». La rationalité technique doit intégrer la \cle{rationalité vitale} comme contrainte de conception (contrainte Ubuntu).
\end{itemize}

\textbf{Conclusion :} La rationalité n'est pas \emph{une} essence figée (Descartes) ni \emph{multiple} incommensurable (relativisme). Elle est \cle{pluriverselle} : une capacité universelle (tout homme raisonne) qui s'\cle{incarne} historiquement en formes diverses. Le défi pour l'Afrique technique n'est pas de choisir entre Descartes et Wiredu, mais de \cle{les faire dialoguer} dans l'atelier de l'ingénieur.
\end{exoresolu}

\begin{attention}[Les 5 erreurs qui coûtent des points à l'examen (Philosophie africaine)]
\begin{enumerate}
    \item \textbf{Confondre « Ethnophilosophie » et « Philosophie africaine » :} Dire « La philosophie africaine, c'est l'ethnophilosophie (Tempels) » est une \textbf{erreur épistémologique majeure}. L'ethnophilosophie est \emph{une tendance} (dépassée par la critique professionnelle). La philosophie africaine \emph{aujourd'hui} est multiple (professionnelle, herméneutique, décoloniale). \textit{Parade :} Toujours qualifier : « L'ethnophilosophie \emph{affirme que}... », « La philosophie professionnelle \emph{répond que}... ».
    \item \textbf{Essentialiser « L'Africain » ou « La Tradition » :} Écrire « L'Africain pense par participation » ou « La tradition est immuable ». C'est reprendre à son compte le \cle{racisme ontologique} de Lévy-Bruhl et Hegel. \textit{Parade :} Parler de « \cle{conceptions} africaines », « \cle{dynamiques} traditionnelles », « \cle{débats} internes ». Utiliser le conditionnel ou l'attribution : « Selon Kagame... », « Pour Wiredu... ».
    \item \textbf{Oublier la \cle{décolonialité} dans la réorientation actuelle :} Traiter la partie 3 du programme comme une simple suite chronologique (après Hountondji vient Mudimbe) sans saisir la \cle{rupture épistémologique} : passage de la \cle{reconnaissance} (exister dans l'académie) à la \cle{libération} (penser depuis soi, langues africaines, justice cognitive). \textit{Parade :} Mots-clés obligatoires : \cle{décolonialité}, \cle{épistémologies du Sud}, \cle{traduction}, \cle{justice cognitive}, \cle{souveraineté épistémique}.
    \item \textbf{Dissertation : plan « catalogue d'auteurs » au lieu de plan \cle{dialectique/thématique}.} Faire I. Tempels, II. Hountondji, III. Wiredu = \textbf{hors sujet}. Le plan doit répondre à la \cle{problématique} par des \cle{arguments} (Thèse/Antithèse/Synthèse), les auteurs sont des \cle{autorités} citées \emph{dans} les arguments. \textit{Parade :} Vérifier que les titres I, II, III sont des \cle{phrases affirmatives} (arguments), pas des noms propres.
    \item \textbf{Commentaire : le « résumé paraphrasé » au lieu de l'\cle{explication}.} Recopier le texte avec des mots différents sans analyser la \cle{structure argumentative} (thèse, arguments, concepts, présupposés, portée). \textit{Parade :} Pour chaque paragraphe analysé : « L'auteur \cle{soutient} que... », « Il \cle{argumente} en montrant que... », « Le concept clé ici est... », « La portée de cet argument est... ».
\end{enumerate}
\end{attention}

\newpage

\coursec{Exercices}

\courssub{Je vérifie mes connaissances}

\begin{exercice}[QCM : Thèses ethnocentristes et naissance de la philosophie africaine]
Parmi les affirmations suivantes, cochez celle qui est \textbf{exacte} :
\begin{enumerate}
    \item L'ethnophilosophie (ex. Tempels, Kagamé) est critiquée par Hountondji pour son caractère \emph{individuel} et \emph{critique}.
    \item La « philosophie sagace » (Oruka) consiste à interviewer des sages illettrés pour en extraire une pensée cohérente, individuelle et critique.
    \item Le débat sur l'existence de la philosophie africaine est encore la question centrale de la « réorientation actuelle » (années 2000+).
    \item L'ethnocentrisme hégélien affirme que l'Afrique a une histoire dialectique autonome depuis l'Antiquité.
\end{enumerate}
\end{exercice}

\begin{exercice}[Vrai / Faux justifié : Tendances et concepts]
Pour chaque affirmation, répondez \textbf{Vrai} ou \textbf{Faux} et justifiez en une phrase précise (auteur ou concept clé).
\begin{enumerate}
    \item Paulin Hountondji (\emph{Sur la philosophie africaine}) défend l'idée que la philosophie africaine traditionnelle est un bloc homogène, anonyme et immuable.
    \item Le concept d'\cle{Ubuntu} (« Je suis parce que nous sommes ») fonde une éthique relationnelle applicable à la gouvernance d'entreprise et au travail d'équipe.
    \item Kwasi Wiredu prône la « décolonisation conceptuelle » : il s'agit de rejeter toute langue européenne pour ne philosopher qu'en langues africaines.
    \item Fabien Eboussi Boulaga (\emph{La Crise du Muntu}) analyse l'aliénation de l'Africain comme une perte de subjectivité historique et ontologique.
\end{enumerate}
\end{exercice}

\begin{exercice}[Définitions opérationnelles : Concepts techniques et philosophiques]
Donnez une définition rigoureuse (2 à 3 lignes) pour chacun des termes suivants en les reliant au programme de Première technique :
\begin{enumerate}
    \item \cle{Ethnophilosophie} (tendance fondatrice).
    \item \cle{Philosophie sagace} (Henry Odera Oruka).
    \item \cle{Décolonialité} (réorientation actuelle).
    \item \cle{Palabre} (comme méthode de résolution de conflits et de gouvernance).
\end{enumerate}
\end{exercice}

\begin{exercice}[Correspondances : Auteurs / Œuvres / Thèses centrales]
Reliez chaque auteur à son œuvre majeure et à la thèse qui le caractérise (tracez des flèches ou faites un tableau) :
\begin{center}
\begin{tabularx}{\linewidth}{|X|X|X|}\hline
\textbf{Auteur} & \textbf{Œuvre majeure} & \textbf{Thèse / Apport central} \\ \hline
Paulin Hountondji & \emph{La Crise du Muntu} & Critique radicale de l'ethnophilosophie ; philosophie comme discours critique, scientifique et individuel. \\ \hline
Fabien Eboussi Boulaga & \emph{Philosophie et langue africaine} & Analyse de l'aliénation du « Muntu » ; appel à une refondation anthropologique et historique. \\ \hline
Kwasi Wiredu & \emph{Sur la philosophie africaine} & Décolonisation conceptuelle ; distinction philosophie / pensée traditionnelle ; logique culturelle. \\ \hline
Henry Odera Oruka & \emph{Sagesse africaine : philosophie sagace} & Identification de sages individuels critiques ; distinction « sagace populaire » / « sagace philosophique ». \\ \hline
\end{tabularx}
\end{center}
\end{exercice}

\courssub{Je m'entraîne}

\begin{exercice}[Analyse d'argumentaire : Hountondji contre l'ethnophilosophie]
Reprenez les \textbf{trois arguments principaux} de Paulin Hountondji contre l'ethnophilosophie (dans \emph{Sur la philosophie africaine}, 1976). Pour chacun, donnez un exemple concret (proverbe, pratique sociale ou discours courant au Cameroun) qui illustre ce que Hountondji dénonce comme « pensée de groupe, anonyme, non critique ».
\begin{methode}[Structure attendue]
\begin{enumerate}
\item Argument 1 (ex: confusion philosophie / vision du monde) + Exemple camerounais.
\item Argument 2 (ex: anonymat / absence de sujet pensant) + Exemple camerounais.
\item Argument 3 (ex: caractère figé, immuable, non historique) + Exemple camerounais.
\end{enumerate}
\end{methode}
\end{exercice}

\begin{exercice}[Application technique : L'Ubuntu dans la gestion de projet BTP]
Vous êtes ingénieur en chef sur un chantier de construction de route (Marché public). Un sous-traitant local a un retard de 3 semaines pour cause de deuil familial prolongé (rites funéraires obligatoires). Le contrat prévoit des pénalités de retard strictes.
\begin{enumerate}
    \item En vous appuyant sur le concept d'\cle{Ubuntu} (\emph{Umuntu ngumuntu ngabantu}), proposez une \textbf{décision managériale argumentée} (sanction, aménagement, médiation) qui concilie rigueur contractuelle et éthique relationnelle africaine.
    \item Identifiez le risque éthique (selon le code de déontologie de l'ingénieur) si vous appliquez \emph{strictement} la clause de pénalité sans tenir compte du contexte social.
\end{enumerate}
\end{exercice}

\begin{exercice}[Débat : Philosophie sagace vs Ethnophilosophie sur le terrain]
Un sociologue et un philosophe enquêtent dans un village Bamileke sur la notion de « succession patrimoniale ».
\begin{enumerate}
    \item Décrivez la \textbf{méthode} du sociologue (proche de l'ethnophilosophie) : que cherche-t-il ? Comment traite-t-il les paroles recueillies ?
    \item Décrivez la \textbf{méthode} du philosophe sagace (selon Oruka) : quel type d'interlocuteur cherche-t-il ? Quelle question lui pose-t-il ? Comment valide-t-il le propos ?
    \item Quelle méthode permet d'identifier une \cle{philosophie} (au sens strict : discours rationnel, argumenté, individuel) ? Justifiez.
\end{enumerate}
\end{exercice}

\begin{exercice}[Synthèse : La réorientation actuelle (Décolonialité / Développement)]
Lisez l'extrait suivant (fictif, inspiré de la littérature actuelle) : 
\begin{quote}
« La question n'est plus "Y a-t-il une philosophie africaine ?" mais "Quelle philosophie pour l'Afrique aujourd'hui ?". Il faut penser l'État, la technique, l'économie, l'écologie avec nos propres catégories, sans attendre la validation occidentale. »
\end{quote}
\begin{enumerate}
    \item Identifiez \textbf{deux ruptures} majeures avec les débats des années 1970-1990 (existence, identité, légitimation).
    \item Citez \textbf{deux concepts opératoires} de cette réorientation (ex: \cle{décolonialité}, \cle{epistemicide}, \cle{développement endogène}, \cle{biopolitique}) et définissez-les brièvement.
    \item Donnez un \textbf{exemple camerounais actuel} (politique publique, innovation technique, mouvement citoyen) qui illustre cette « philosophie pour l'Afrique ».
\end{enumerate}
\end{exercice}

\courssub{Je me prépare au Bac}

\begin{exercice}[Sujet de Dissertation Type Bac Probatoire]
\textbf{Sujet :} « La philosophie africaine est-elle une philosophie de la libération ? »

\begin{enumerate}
    \item \textbf{Analyse du sujet :} Définissez les termes « philosophie africaine » (tendances, évolution) et « libération » (épistémique, politique, existentielle, technique). Quelle est la tension centrale ?
    \item \textbf{Problématisation :} Formulez la problématique en une question directe qui oppose au moins deux thèses fortes.
    \item \textbf{Plan dialectique détaillé (3 parties, 2 sous-parties chacune) :} Rédigez l'intitulé exact de chaque partie et sous-partie. Indiquez pour chaque sous-partie : l'auteur de référence, le concept clé, l'argument principal et un exemple camerounais.
    \begin{itemize}
        \item \textbf{Thèse :} La philosophie africaine comme libération épistémique (décolonisation du savoir, critique de l'ethnocentrisme).
        \item \textbf{Antithèse :} Les limites et risques d'une libération politique confisquée (néocolonialisme, dictatures, instrumentalisation de la « tradition »).
        \item \textbf{Synthèse :} Une libération existentielle et technique par l'innovation endogène (Ubuntu, Palabre, éthique du développement, appropriation technique).
    \end{itemize}
    \item \textbf{Introduction rédigée :} Accroche, définition, problématique, annonce du plan (15 lignes max).
\end{enumerate}
\end{exercice}

\begin{exercice}[Sujet d'Explication de Texte Type Bac Probatoire]
\textbf{Texte :} Extrait de \emph{La Crise du Muntu} (Fabien Eboussi Boulaga, 1977) — \emph{sur l'aliénation du Muntu}.
\begin{quote}
« Le Muntu est l'homme africain dans sa situation historique concrète : un être qui a été dépossédé de son histoire, de sa culture, de son langage, de ses dieux, de sa terre, de son travail, de sa famille, de son corps même. [...] L'aliénation n'est pas seulement une dépossession, elle est une aliénation \emph{dans} l'être : le Muntu devient étranger à lui-même. Il ne se reconnaît plus dans ce qu'il fait, dans ce qu'il dit, dans ce qu'il pense. Il vit une vie qui n'est pas la sienne. »
\end{quote}

\begin{enumerate}
    \item \textbf{Thèse de l'auteur :} Résumez en une phrase la thèse défendue dans cet extrait.
    \item \textbf{Concepts clés :} Définissez philosophiquement \cle{Muntu} et \cle{aliénation} (au sens ontologique/existentiel chez Eboussi Boulaga, distinct du sens marxien pur).
    \item \textbf{Argumentation :} Relevez \textbf{trois étapes} de l'argumentation dans le texte (dépossession multiple $\rightarrow$ aliénation dans l'être $\rightarrow$ vie inauthentique). Expliquez la logique.
    \item \textbf{Portée critique :} En quoi cette analyse dépasse-t-elle la simple critique coloniale pour interroger la \textbf{post-colonialité} (État, élites, éducation, technique) au Cameroun aujourd'hui ?
    \item \textbf{Application contextuelle :} Donnez un exemple précis (système éducatif, administration, entreprise, rapport à la technique) où le « Muntu » camerounais contemporain vit encore cette « vie qui n'est pas la sienne ».
\end{enumerate}
\end{exercice}

\begin{exercice}[Cas Pratique / Épreuve Pratique : Éthique de l'ingénieur et Philosophie africaine]
\textbf{Situation professionnelle :}
Vous êtes \textbf{Chef de Projet BTP} (diplômé ENSP/ENSET) pour une entreprise locale. Le marché de construction d'un pont (financement public) arrive à réception. Le Contrôleur des Travaux (fonctionnaire camerounais) vous convoque et vous dit textuellement : 
\begin{quote}
« Le dossier technique est bon. Mais pour signer le PV de réception sans délai, il faut "voir le dossier" : 5\% du montant du marché. C'est la coutume ici, on ne peut pas changer le système. Pense à ta famille, pense à l'avenir de ton entreprise. Si tu refuses, le dossier dort 6 mois. »
\end{quote}

\textbf{Consignes :}
Rédigez votre \textbf{réponse argumentée écrite} (lettre ou mail formel) adressée à votre Directeur Général (avec copie au Contrôleur), pour justifier votre refus de payer le pot-de-vin. Votre réponse doit :
\begin{enumerate}
    \item Mobiliser \textbf{deux concepts de la philosophie africaine} (choisis parmi : \cle{Palabre}, \cle{Ubuntu}, \cle{Honneur / Dignité}, \cle{Responsabilité communautaire}, \cle{Justice réparatrice}) pour fonder votre éthique professionnelle.
    \item Faire référence explicite à \textbf{deux articles du Code de déontologie de l'ingénieur} (ex: intégrité, intérêt public, refus de la corruption, compétence).
    \item Proposer une \textbf{issue procédurale concrète} (Palabre formelle, saisine hiérarchique, ARMP, Tribunal administratif) qui respecte la loi et la relation.
    \item Tenir en \textbf{15 à 20 lignes maximum}, ton ferme, respectueux, professionnel.
\end{enumerate}
\begin{attention}[Pièges à éviter]
\begin{enumerate}
\item Ne pas confondre « coutume » (pratique sociale) et « tradition philosophique » (norme éthique critique).
\item Ne pas tomber dans le relativisme (« c'est la culture ») ni dans le moralisme abstrait sans ancrage africain.
\item La solution technique (saisine ARMP) doit être portée par l'argument éthique philosophique.
\end{enumerate}
\end{attention}
\end{exercice}

\newpage

\coursec{Corrigés des exercices}

\coursec{Exercices corrigés et applications — La Philosophie africaine}

\begin{corrige}[Exercice 1 — QCM : Thèses ethnocentristes et naissance de la philosophie africaine]
La réponse exacte est la \textbf{2}.

\textbf{Justification de chaque item :}
\begin{enumerate}
    \item \textbf{Faux.} Hountondji reproche à l'ethnophilosophie (Tempels, Kagamé) exactement l'inverse : son caractère \emph{collectif}, \emph{anonyme} et \emph{non critique}. Elle présente la pensée africaine comme un « système » de groupe, sans auteur identifié ni débat interne.
    \item \textbf{Exact.} La philosophie sagace d'Henry Odera Oruka consiste à interroger des \cle{sages} traditionnels (souvent illettrés) capables d'une pensée \emph{individuelle}, cohérente et critique, distincte de la simple sagesse populaire (proverbes, coutumes). C'est la « sagace philosophique ».
    \item \textbf{Faux.} La réorientation actuelle (années 2000 et après) dépasse la question de l'\emph{existence} de la philosophie africaine, jugée réglée, pour poser la question de sa \emph{pertinence} : « Quelle philosophie pour l'Afrique aujourd'hui ? » (développement, État, technique, écologie).
    \item \textbf{Faux.} Hegel, dans la \emph{Raison dans l'histoire}, exclut au contraire l'Afrique de l'histoire dialectique universelle : l'Afrique serait « sans histoire », immobile, en dehors de l'esprit. C'est le cœur de la thèse ethnocentriste.
\end{enumerate}
\end{corrige}

\begin{attention}[Points de vigilance — Exercice 1]
Bien distinguer \emph{sagace populaire} (savoir collectif, anonyme) et \emph{sagace philosophique} (pensée critique d'un individu). Ne pas attribuer à Hountondji une critique de l'individualisme : c'est précisément l'individualité du discours philosophique qu'il revendique.
\end{attention}

\begin{corrige}[Exercice 2 — Vrai / Faux justifié : Tendances et concepts]
\begin{enumerate}
    \item \textbf{Faux.} C'est exactement la thèse que Hountondji \emph{combat} : dans \emph{Sur la philosophie africaine} (1976), il dénonce l'ethnophilosophie qui présente la pensée africaine comme un bloc homogène, anonyme et immuable ; pour lui, la philosophie est un discours écrit, individuel et critique, comme la science.
    \item \textbf{Vrai.} L'\cle{Ubuntu} (« Je suis parce que nous sommes », \emph{Umuntu ngumuntu ngabantu}) fonde une éthique relationnelle où la personne se constitue par la communauté ; elle inspire des modèles de gouvernance d'entreprise (management participatif, solidarité dans le travail d'équipe) et de résolution consensuelle des conflits.
    \item \textbf{Faux.} La « décolonisation conceptuelle » de Kwasi Wiredu ne consiste pas à rejeter les langues européennes, mais à \emph{examiner critiquement} les concepts philosophiques hérités de la colonisation pour les éprouver à la lumière des catégories des langues africaines (ex. l'akan) et de la logique universelle. C'est une critique conceptuelle, non un rejet linguistique.
    \item \textbf{Vrai.} Dans \emph{La Crise du Muntu} (1977), Eboussi Boulaga analyse l'aliénation de l'Africain (le \emph{Muntu}) comme dépossession de son histoire, de sa culture et de son être : une perte de subjectivité à la fois historique et ontologique, appelant une refondation de l'existence africaine.
\end{enumerate}
\end{corrige}

\begin{attention}[Points de vigilance — Exercice 2]
La justification exigée doit citer un auteur ou un concept : une réponse « Vrai/Faux » seule ne rapporte aucun point. Attention à ne pas confondre décolonisation \emph{conceptuelle} (Wiredu) et rejet pur et simple de l'Occident.
\end{attention}

\begin{corrige}[Exercice 3 — Définitions opérationnelles]
\begin{enumerate}
    \item \textbf{\cle{Ethnophilosophie} :} tendance fondatrice (Tempels, \emph{La Philosophie bantoue}, 1945 ; Kagamé) qui identifie la philosophie africaine aux visions du monde collectives, anonymes et implicites d'un peuple, reconstruites à partir des coutumes, langues et proverbes. Critiquée par Hountondji pour son absence de discours individuel, écrit et critique.
    \item \textbf{\cle{Philosophie sagace} :} méthode d'Henry Odera Oruka consistant à rechercher, au sein des cultures africaines, des sages individuels capables d'une réflexion rationnelle, cohérente et critique sur leur propre tradition. Elle distingue la sagace populaire (sagesse collective) de la sagace philosophique (pensée critique d'un individu identifié).
    \item \textbf{\cle{Décolonialité} :} orientation actuelle de la philosophie africaine qui vise à défaire la « colonialité du pouvoir et du savoir » survivant après l'indépendance politique : décoloniser les concepts, les institutions, l'éducation et la technique pour penser l'Afrique avec ses propres catégories.
    \item \textbf{\cle{Palabre} :} institution traditionnelle africaine de parole publique et de délibération collective, visant la résolution consensuelle des conflits et la décision commune. Comme méthode de gouvernance, elle fonde la légitimité sur la discussion argumentée, l'écoute de tous et la recherche du consensus plutôt que sur la contrainte.
\end{enumerate}
\end{corrige}

\begin{corrige}[Exercice 4 — Correspondances : Auteurs / Œuvres / Thèses centrales]
Les associations correctes sont :
\begin{center}
\begin{tabularx}{\linewidth}{|X|X|X|}\hline
\rowcolor{popblueL}
\textbf{Auteur} & \textbf{Œuvre majeure} & \textbf{Thèse / Apport central} \\ \hline
Paulin Hountondji & \emph{Sur la philosophie africaine} (1976) & Critique radicale de l'ethnophilosophie ; philosophie comme discours critique, scientifique et individuel. \\ \hline
Fabien Eboussi Boulaga & \emph{La Crise du Muntu} (1977) & Analyse de l'aliénation du « Muntu » ; appel à une refondation anthropologique et historique. \\ \hline
Kwasi Wiredu & \emph{Philosophie et langue africaine} & Décolonisation conceptuelle ; distinction philosophie / pensée traditionnelle ; logique culturelle. \\ \hline
Henry Odera Oruka & \emph{Sagesse africaine : philosophie sagace} & Identification de sages individuels critiques ; distinction « sagace populaire » / « sagace philosophique ». \\ \hline
Léopold Sédar Senghor / Kwame Nkrumah & \emph{Liberté / Consciencisme} & Philosophie nationaliste / idéologique : Négritude, Socialisme africain, Consciencisme ; affirmation d'une identité culturelle comme base politique. \\ \hline
\end{tabularx}
\end{center}
\textbf{Erreurs fréquentes à éviter :} attribuer \emph{La Crise du Muntu} à Hountondji (c'est Eboussi Boulaga, philosophe camerounais) ; confondre Wiredu (décolonisation conceptuelle) et Oruka (sagace). Retenir le repère : Hountondji = critique de l'ethnophilosophie ; Oruka = sages individuels ; Wiredu = langue et concepts ; Eboussi Boulaga = aliénation du Muntu ; Senghor/Nkrumah = affirmation identitaire politique.
\end{corrige}

\begin{attention}[Points de vigilance — Exercice 4]
Le programme exige la connaissance des \emph{quatre tendances fondatrices} : 1. Ethnophilosophie (Tempels, Kagamé), 2. Philosophie nationaliste/idéologique (Senghor, Nkrumah, Nyerere), 3. Philosophie sagace (Oruka), 4. Philosophie critique / herméneutique (Hountondji, Wiredu, Eboussi Boulaga). Le tableau ci-dessus les couvre toutes.
\end{attention}

\begin{aretenir}[Les quatre tendances fondatrices de la philosophie africaine moderne]
\begin{enumerate}
    \item \textbf{Ethnophilosophie} (Tempels, Kagamé) : reconstruction de la pensée collective implicite (coutumes, proverbes) comme « philosophie » du groupe. Critique : anonymat, absence de rationalité critique.
    \item \textbf{Philosophie nationaliste / idéologique} (Senghor, Nkrumah, Nyerere, Césaire) : mise en œuvre de concepts philosophiques (Négritude, Consciencisme, Ujamaa) au service de la libération politique et de l'affirmation identitaire.
    \item \textbf{Philosophie sagace} (Oruka) : découverte de sages \emph{individus} critiques dans la tradition orale ; distinction sagace populaire / sagace philosophique.
    \item \textbf{Philosophie critique / herméneutique} (Hountondji, Wiredu, Eboussi Boulaga, Okolo) : exigence de scientificité, décolonisation conceptuelle, analyse de l'aliénation, interprétation critique de la tradition.
\end{enumerate}
\end{aretenir}

\begin{corrige}[Exercice 5 — Analyse d'argumentaire : Hountondji contre l'ethnophilosophie]
\begin{enumerate}
    \item \textbf{Argument 1 — Confusion entre philosophie et vision du monde.} Pour Hountondji, l'ethnophilosophie prend pour « philosophie » ce qui n'est qu'une \emph{Weltanschauung} collective (croyances, mythes, coutumes) ; or la philosophie est un discours explicite, argumenté et problématisant. \emph{Exemple camerounais :} présenter le proverbe beti « la case du père ne s'oublie pas » ou les croyances sur les ancêtres comme « la philosophie du peuple » revient à dissoudre la philosophie dans la sagesse populaire.
    \item \textbf{Argument 2 — Anonymat et absence de sujet pensant.} L'ethnophilosophie décrit une pensée « de tous et de personne », sans auteur qui assume et défende ses thèses ; or philosopher, c'est un \emph{je} qui argumente et répond à la critique. \emph{Exemple camerounais :} dire « chez nous, on pense que... » lors d'une palabre ou d'un débat illustre ce sujet collectif anonyme : nul n'est responsable de la thèse, nul ne peut la discuter.
    \item \textbf{Argument 3 — Caractère figé, immuable, non historique.} L'ethnophilosophie présente la pensée africaine comme un système achevé, identique à lui-même depuis toujours, sans ruptures ni controverses ; or toute philosophie véritable est \emph{historique} : elle débat, se contredit, évolue. \emph{Exemple camerounais :} invoquer « la tradition immuable de nos ancêtres » pour interdire toute discussion sur l'héritage ou le statut des femmes fige la culture et exclut le progrès critique.
\end{enumerate}
\end{corrige}

\begin{attention}[Points de vigilance — Exercice 5]
Chaque argument doit être formulé comme une \emph{thèse générale} (sur la nature de la philosophie) avant l'exemple. L'exemple camerounais doit illustrer ce que Hountondji \emph{dénonce}, non ce qu'il défend.
\end{attention}

\begin{corrige}[Exercice 6 — Application technique : L'Ubuntu dans la gestion de projet BTP]
\begin{enumerate}
    \item \textbf{Décision managériale argumentée.} Au nom de l'\cle{Ubuntu} (\emph{Umuntu ngumuntu ngabantu}, « je suis parce que nous sommes »), la personne se réalise dans la relation et la solidarité communautaire : le deuil et les rites funéraires ne sont pas un caprice individuel mais une obligation qui constitue l'identité même du sous-traitant. Je propose donc : (a) une \emph{médiation} (réunion de chantier élargie, sur le modèle de la palabre) exposant la situation au maître d'ouvrage ; (b) un \emph{aménagement} : avenant au contrat décalant le planning de 3 semaines, avec réorganisation des tâches non dépendantes pour limiter l'impact global ; (c) la \emph{suspension} des pénalités pour cette période, le retard n'étant pas imputable à une négligence mais à un événement social majeur. La rigueur contractuelle est préservée (avenant écrit, traçabilité), mais l'éthique relationnelle africaine commande de traiter le partenaire comme une personne insérée dans une communauté, non comme une simple variable de production.
    \item \textbf{Risque éthique.} Appliquer strictement la pénalité violerait le principe déontologique d'\emph{humanité et d'équité} dans l'exercice de la profession : le code de déontologie de l'ingénieur impose de concilier l'intérêt du projet avec le respect des personnes et du contexte social. Le risque est double : injustice envers un partenaire de bonne foi (sanction d'une force majeure sociale) et rupture durable de la confiance, donc de la qualité de la collaboration — contraire à l'intérêt même de l'ouvrage public.
\end{enumerate}
\end{corrige}

\begin{corrige}[Exercice 7 — Débat : Philosophie sagace vs Ethnophilosophie sur le terrain]
\begin{enumerate}
    \item \textbf{Méthode du sociologue (proche de l'ethnophilosophie).} Il cherche la \emph{vision collective} du groupe : il interroge de nombreux villageois, recueille proverbes, coutumes et récits sur la succession, puis reconstruit un « système de pensée bamileke » homogène. Les paroles sont traitées comme les expressions \emph{anonymes} d'une conscience collective : l'individu n'est qu'un témoin de la tradition, jamais un auteur.
    \item \textbf{Méthode du philosophe sagace (Oruka).} Il cherche un interlocuteur particulier : un \emph{sage} reconnu pour sa capacité de réflexion critique (souvent un ancien, parfois illettré). Il ne lui demande pas « que dit la coutume ? » mais : « \emph{Toi}, que penses-tu de la succession patrimoniale ? La coutume est-elle juste ? Pourquoi ? » Il valide le propos par l'\emph{examen critique} : il teste la cohérence de l'argumentation, confronte le sage aux objections, vérifie que la pensée est personnelle et raisonnée, et non la simple répétition du proverbe.
    \item \textbf{Conclusion.} Seule la méthode sagace permet d'identifier une \cle{philosophie} au sens strict : elle met au jour un discours \emph{rationnel, argumenté et individuel}, assumé par un sujet qui le défend. La méthode ethnophilosophique ne livre qu'une pensée collective et anonyme — une sagesse, non une philosophie (c'est précisément la critique de Hountondji, qu'Oruka dépasse en montrant que l'individu critique existe \emph{dans} la tradition orale).
\end{enumerate}
\end{corrige}

\begin{corrige}[Exercice 8 — Synthèse : La réorientation actuelle]
\begin{enumerate}
    \item \textbf{Deux ruptures majeures :}
    \begin{itemize}
        \item \emph{De l'existence à la pertinence :} les débats 1970-1990 demandaient « Y a-t-il une philosophie africaine ? » (légitimation face à Hegel et Lévy-Bruhl) ; le texte tient l'existence pour acquise et demande « Quelle philosophie \emph{pour} l'Afrique ? » — la philosophie devient un outil d'action.
        \item \emph{De la quête d'identité à l'autonomie épistémique :} au lieu de chercher la reconnaissance (validation occidentale), il s'agit de penser l'État, la technique, l'économie et l'écologie « avec nos propres catégories » : fin de l'attitude de demande de légitimation.
    \end{itemize}
    \item \textbf{Deux concepts opératoires :}
    \begin{itemize}
        \item \cle{Décolonialité} : déconstruction de la « colonialité » (pouvoir, savoir, être) qui survit à la colonisation dans les institutions et les esprits.
        \item \cle{Développement endogène} : modèle de développement fondé sur les ressources, les savoirs et les valeurs propres des communautés africaines, et non sur l'imitation de modèles importés. (\emph{Autres réponses acceptées :} \cle{epistemicide} = destruction des savoirs locaux par la colonisation ; \cle{biopolitique} = gouvernement des corps et des vies.)
    \end{itemize}
    \item \textbf{Exemple camerounais actuel :} la promotion des matériaux locaux (briques de terre comprimée, bambou) dans les politiques de construction durable illustre une « philosophie pour l'Afrique » : innovation technique pensée à partir des ressources et savoir-faire endogènes, répondant aux défis écologiques et économiques sans copier les modèles occidentaux. (\emph{Sont aussi recevables :} les tontines numérisées, l'intégration de la palabre dans la médiation administrative, les mouvements citoyens de revalorisation des langues nationales dans l'éducation.)
\end{enumerate}
\end{corrige}

\begin{corrige}[Exercice 9 — Sujet de Dissertation Type Bac Probatoire : « La philosophie africaine est-elle une philosophie de la libération ? »]
\textbf{1. Analyse du sujet.} \emph{Philosophie africaine} : ensemble des démarches critiques nées de la contestation des thèses ethnocentristes (Hegel, Lévy-Bruhl) et structurées en tendances (ethnophilosophie, sagace, critique nationaliste, herméneutique), aujourd'hui réorientées vers la pertinence sociale. \emph{Libération} : au triple sens épistémique (se libérer de la domination des savoirs coloniaux), politique (émancipation des peuples) et existentiel/technique (redevenir sujet de son histoire et maître de son développement). \emph{Tension centrale :} la philosophie africaine est née d'un geste de libération, mais la libération annoncée a-t-elle été réalisée, ou confisquée ?

\textbf{2. Problématique.} « Si la philosophie africaine naît de la volonté de libérer la raison africaine du mépris colonial, cette libération épistémique suffit-elle à émanciper réellement les peuples, ou faut-il la prolonger par une libération existentielle et technique ? »

\textbf{3. Plan dialectique détaillé.}
\begin{itemize}
    \item \textbf{Thèse : la philosophie africaine est d'abord une libération épistémique.}
    \begin{itemize}
        \item 1.1. Critique de l'ethnocentrisme : réfutation de Hegel et Lévy-Bruhl ; \emph{auteur :} Hountondji ; \emph{concept :} critique de l'ethnophilosophie ; \emph{argument :} arracher la raison africaine au préjugé d'infériorité ; \emph{exemple :} les débats universitaires à Yaoundé sur la scientificité.
        \item 1.2. Décolonisation conceptuelle : \emph{auteur :} Wiredu ; \emph{concept :} examen critique des concepts hérités ; \emph{argument :} philosopher suppose de repenser les catégories mentales colonisées ; \emph{exemple :} repenser « développement » hors du modèle importé.
    \end{itemize}
    \item \textbf{Antithèse : les limites d'une libération confisquée.}
    \begin{itemize}
        \item 2.1. La libération politique détournée : indépendances sans décolonisation réelle (néocolonialisme, dictatures) ; \emph{concept :} aliénation du Muntu (Eboussi Boulaga) ; \emph{exemple :} élites formées à l'école coloniale reproduisant la domination.
        \item 2.2. Le risque de l'instrumentalisation de la « tradition » : l'ethnophilosophie figée peut justifier l'immobilisme (« c'est notre culture ») ; \emph{auteur :} Hountondji ; \emph{exemple :} refus du débat sur certaines coutumes au nom de la tradition.
    \end{itemize}
    \item \textbf{Synthèse : une libération existentielle et technique par l'innovation endogène.}
    \begin{itemize}
        \item 3.1. Refondation du sujet : \emph{auteur :} Eboussi Boulaga ; \emph{concept :} refondation anthropologique ; \emph{exemple :} éducation valorisant langues et savoirs nationaux.
        \item 3.2. Éthique du développement : \emph{concepts :} Ubuntu, palabre, développement endogène ; \emph{argument :} la libération s'achève dans l'appropriation technique et la gouvernance responsable ; \emph{exemple :} matériaux locaux, médiation par la palabre.
    \end{itemize}
\end{itemize}

\textbf{4. Introduction rédigée (modèle).} En affirmant que l'Afrique est « sans histoire » et sans philosophie, Hegel a posé un préjugé que toute la pensée africaine moderne s'est employée à renverser. La philosophie africaine désigne l'ensemble des démarches critiques par lesquelles des penseurs africains — de Hountondji à Wiredu, d'Oruka à Eboussi Boulaga — ont revendiqué pour leur continent le droit au discours rationnel. La libération, elle, peut s'entendre au sens épistémique (libérer le savoir), politique (libérer les peuples) ou existentiel (libérer le sujet de son aliénation). Dès lors, une question s'impose : la philosophie africaine est-elle une philosophie de la libération ? Nous montrerons d'abord qu'elle naît comme libération épistémique du savoir ; nous verrons ensuite que cette libération a été en partie confisquée par la postcolonie ; nous soutiendrons enfin qu'elle ne s'accomplit pleinement que comme libération existentielle et technique, au service d'un développement endogène.

\textbf{Barème indicatif (sur 20) :} analyse des termes et problématique : 4 pts ; pertinence du plan dialectique : 6 pts ; exactitude des références (auteurs, concepts) : 4 pts ; exemples camerounais : 3 pts ; introduction rédigée et langue : 3 pts.
\end{corrige}

\begin{attention}[Points de vigilance — Exercice 9]
Ne pas réduire « libération » au seul sens politique ; citer au moins trois auteurs avec leurs thèses exactes ; la synthèse ne doit pas être un compromis mou mais un dépassement (libération existentielle et technique).
\end{attention}

\begin{corrige}[Exercice 10 — Explication de texte : Eboussi Boulaga, \emph{La Crise du Muntu}]
\textbf{1. Thèse de l'auteur.} L'homme africain (le Muntu), dépossédé par la colonisation de toutes les dimensions de son existence, ne subit pas seulement une spoliation extérieure : il est devenu \emph{étranger à lui-même}, vivant une vie inauthentique qui n'est pas la sienne.

\textbf{2. Concepts clés.} \cle{Muntu} : mot bantu signifiant « l'homme », désignant chez Eboussi Boulaga l'Africain concret, saisi dans sa situation historique singulière (et non une essence abstraite). \cle{Aliénation} : au sens ontologico-existentiel (et non seulement économique comme chez Marx), processus par lequel un être perd la possession de soi — histoire, langue, dieux, corps — jusqu'à ne plus se reconnaître dans ses propres actes, paroles et pensées : l'aliénation est « \emph{dans} l'être ».

\textbf{3. Argumentation (trois étapes).}
\begin{enumerate}
    \item \emph{Constat de la dépossession multiple :} énumération cumulative (histoire, culture, langage, dieux, terre, travail, famille, corps) montrant que la colonisation a touché la totalité de l'existence, pas seulement l'économie.
    \item \emph{Enfoncement de l'aliénation dans l'être :} la dépossession extérieure devient intérieure — le Muntu « devient étranger à lui-même » ; la logique est celle d'un passage de l'\emph{avoir} (ce qu'on m'a pris) à l'\emph{être} (ce que je suis devenu).
    \item \emph{Conséquence existentielle : la vie inauthentique :} « il vit une vie qui n'est pas la sienne » — l'aliéné ne peut plus être l'auteur de sa propre existence ; il imite, répète, survit dans des formes qui ne sont pas les siennes.
\end{enumerate}

\textbf{4. Portée critique.} L'analyse dépasse la dénonciation du colonisateur : l'aliénation peut \emph{survivre à la colonisation}. Après l'indépendance, l'État postcolonial, ses élites, son école et ses techniques importées peuvent reproduire la même dépossession : le Camerounais d'aujourd'hui peut rester « étranger à lui-même » dans des institutions qui ne sont pas pensées à partir de son expérience. La question devient postcoloniale : qui détient le pouvoir d'aliéner, sinon désormais nos propres structures ?

\textbf{5. Application contextuelle.} Exemple : le système éducatif où l'élève camerounais apprend l'histoire, la géographie et la technique presque exclusivement à travers des références européennes, dans une langue qui n'est pas celle de sa communauté — il réussit scolairement en se dépossédant de ses savoirs et de sa langue : il « vit une vie qui n'est pas la sienne ». (Autres exemples recevables : administration copiée sur le modèle colonial ; cadre qui dévalorise les savoir-faire locaux ; consommation technique sans appropriation.)

\textbf{Barème indicatif (sur 20) :} thèse : 3 pts ; définitions des concepts : 4 pts ; analyse des trois étapes : 6 pts ; portée critique : 4 pts ; exemple contextualisé : 3 pts.
\end{corrige}

\begin{attention}[Points de vigilance — Exercice 10]
Bien distinguer l'aliénation ontologique (Eboussi Boulaga) de l'aliénation économique (Marx) ; ne pas plaquer une lecture uniquement « anticoloniale » : le texte vise aussi la postcolonie ; l'exemple doit être camerounais et actuel.
\end{attention}

\begin{corrige}[Exercice 11 — Cas Pratique : Éthique de l'ingénieur et Philosophie africaine]
\textbf{Modèle de réponse argumentée (lettre au Directeur Général, copie au Contrôleur).}

\emph{Objet : Refus de toute transaction financière liée à la réception des travaux du pont.}

Monsieur le Directeur Général,

J'ai l'honneur de porter à votre connaissance qu'une somme de 5\,\% du marché m'a été demandée en échange de la signature sans délai du PV de réception. Je refuse cette transaction, et je motive mon refus.

Notre profession engage d'abord notre \cle{honneur} et notre \cle{dignité} : dans nos traditions, la parole donnée et le nom de la famille valent plus que l'argent ; céder serait déshonorer l'entreprise et ma communauté. Ensuite, l'\cle{Ubuntu} m'oblige : « je suis parce que nous sommes » — le pont appartient au « nous » ; un ouvrage réceptionné contre pot-de-vin, sans contrôle loyal, mettrait en danger la sécurité des usagers, c'est-à-dire la communauté qui me constitue.

Ce refus s'appuie sur le code de déontologie de l'ingénieur : le principe d'\emph{intégrité} interdit tout avantage indu dans l'exercice de la fonction ; le principe de primauté de l'\emph{intérêt public} et de la sécurité impose que la réception soit prononcée sur la seule qualité technique de l'ouvrage, laquelle est ici reconnue.

Je propose une issue procédurale conforme à notre éthique de la \cle{palabre} : réunir une réunion contradictoire et consignée (maître d'ouvrage, contrôle, entreprise) pour examiner le dossier ; à défaut, saisir notre hiérarchie et, s'il y a blocage, l'ARMP ou le juge administratif. La loi et la relation seront ainsi respectées.

Veuillez agréer, Monsieur le Directeur Général, l'expression de ma considération distinguée.

\textbf{Barème indicatif (sur 20) :} deux concepts philosophiques africains correctement mobilisés : 6 pts ; deux principes déontologiques explicites : 4 pts ; issue procédurale concrète et légale : 4 pts ; forme (registre, longueur 15-20 lignes, ton) : 3 pts ; cohérence de l'argumentation : 3 pts.
\end{corrige}

\begin{attention}[Points de vigilance — Exercice 11]
Distinguer « coutume » (la pratique du pot-de-vin invoquée par le contrôleur) et « tradition philosophique » (Ubuntu, palabre, honneur) : la corruption n'est pas une valeur africaine. Éviter le relativisme (« c'est la culture ») comme le moralisme sans ancrage. La procédure (ARMP, réunion contradictoire) doit être portée par l'argument éthique, non s'y substituer.
\end{attention}

\newpage

\coursec{Fiche bilan}

\courssub{Synthèse visuelle : Carte mentale de la leçon}
\begin{popfigure}
\begin{tikzpicture}[
  mindmap,
  concept color=popblue,
  level 1 concept/.append style={level distance=4.5cm, sibling angle=72},
  level 2 concept/.append style={level distance=3cm},
  every node/.style={font=\small, text width=2.8cm, align=center},
  concept/.append style={minimum size=0, inner sep=4pt},
  grow cyclic
]
\node[concept, align=center]{Philosophie Africaine\\ (1re Tech)}
  child[concept color=poporange] { node[concept, align=center]{1. Procès de la raison\\Ethnocentrisme colonial\\ (Hegel, Lévy-Bruhl)} }
  child[concept color=popblue] { node[concept, align=center]{2. Quatre tendances\\Ethnophilosophie, Sage,\\Nationaliste, Professionnelle} }
  child[concept color=popgreen] { node[concept, align=center]{3. Réorientation actuelle\\Décolonialité, Développement\\Pertinence sociale} }
  child[concept color=poppurple] { node[concept, align=center]{4. Atelier pratique (TP)\\Analyse de documents\\Outils heuristiques} }
  child[concept color=poppink] { node[concept, align=center]{5. Intégration\\Éthique technique\\Gouvernance, Citoyenneté} };
\end{tikzpicture}
\legende{Carte mentale récapitulative des 5 axes de la leçon. Chaque branche correspond à une étape de la progression officielle MINESEC.}
\end{popfigure}

\courssub{Bilan des savoirs et savoir-faire}

\textbf{1. Déclinaison des notions clés}
\begin{itemize}
  \item \cle{Ethnocentrisme} : Attitude consistant à juger les autres cultures à partir des valeurs de la sienne, en la posant comme universelle. En philosophie : déni de rationalité à l'Africain (Hegel : « Afrique sans histoire » ; Lucien Lévy-Bruhl : « mentalité prélogique »).
  \item \cle{Philosophie africaine} : Réflexion critique, rationnelle et argumentée produite par des Africains ou sur l'Afrique, portant sur l'être, le savoir, la valeur, la société.
  \item \cle{Ethnophilosophie} (Tempels, Kagamé) : Reconstruction systématique de la pensée traditionnelle implicite (ontologie bantoue, force vitale). Risque : essentialisme, « philosophie de salon ».
  \item \cle{Philosophie sagace} (Oruka, Paulin Hountondji critique) : Pensée critique individuelle des sages traditionnels (ex. Okemba, Chaungo). Distinction : \emph{folk philosophy} (sagesse populaire) vs \emph{philosophic sagacity} (réflexion personnelle argumentée).
  \item \cle{Philosophie nationaliste-idéologique} (Nkrumah, Nyerere, Senghor) : Mise en concept de l'idéologie politique (consciencisme, ujamaa, négritude). Outil de libération mais risque de dogmatisme d'État.
  \item \cle{Philosophie professionnelle/académique} (Hountondji, Wiredu, Mudimbe, Mbembe) : Exigence de rigueur universitaire, critique de l'ethnophilosophie, dialogue avec l'histoire de la philosophie universelle. \cle{Hountondji} : « Surrogate » (philosophie de substitution) à dépasser.
  \item \cle{Décolonialité} (Ndlovu-Gatsheni, Maldonado-Torres, adaptés) : Délier le savoir du pouvoir colonial, épistémologies du Sud, pertinence du développement endogène.
\end{itemize}

\textbf{2. Les quatre tendances fondatrices (Classification Oruka) — Tableau de révision}
\begin{center}
\begin{tabularx}{\linewidth}{|l|X|X|X|}
\hline
\rowcolor{popblueL}
\textbf{Tendance} & \textbf{Figures emblématiques} & \textbf{Thèse centrale} & \textbf{Apport / Limite (Bac)} \\
\hline
Ethnophilosophie & Placide Tempels (\emph{La Philosophie bantoue}, 1945), Alexis Kagamé (\emph{La philosophie bantu-rwandaise de l'être}, 1956) & La pensée africaine est une ontologie implicite : l'être = force vitale. & + Mise en valeur du patrimoine. -- Essentialisme, absence de sujet critique. \\
\hline
Philosophie sagace & Henry Odera Oruka, Sages : Paul Mbuya Akoko (Okemba), Chaungo Barasa & Existence de philosophes traditionnels individus, critiques, rationnels. & + Preuve de rationalité africaine. -- Difficulté de transcription écrite, risque de folklorisation. \\
\hline
Nationaliste-idéologique & Kwame Nkrumah (\emph{Consciencisme}), Julius Nyerere (\emph{Ujamaa}), Léopold Sédar Senghor (\emph{Négritude}) & La philosophie sert la libération nationale et la construction étatique. & + Ancrage politique concret. -- Risque de propagande, confusion philosophie/idéologie. \\
\hline
Professionnelle / Académique & Paulin Hountondji (\emph{Sur la philosophie africaine}), Kwasi Wiredu, V.Y. Mudimbe, Achille Mbembe & Exigence de scientificité, critique radicale de l'ethnophilosophie, universalisme critique. & + Rigueur épistémologique. -- Risque d'aliénation aux canons occidentaux, déconnexion du terrain. \\
\hline
\end{tabularx}
\end{center}

\textbf{3. Méthodes express (Fiches Bac)}
\begin{methode}[Dissertation type : « L'ethnophilosophie est-elle la vraie philosophie africaine ? »]
\begin{enumerate}
  \item \textbf{Analyse} : Définir ethnophilosophie (sens fort/sens faible), « vraie » (exclusive ? légitime ?), philosophie africaine (critères : rationalité, universalité, historicité).
  \item \textbf{Problématique} : L'ethnophilosophie suffit-elle à fonder une philosophie africaine scientifique ou en est-elle seulement le matériau premier ?
  \item \textbf{Plan dialectique} :
    \begin{enumerate}
      \item Thèse : L'ethnophilosophie comme fondement légitime (Tempels, Kagamé) : ancrage culturel, décolonisation des esprits.
      \item Antithèse : Les limites épistémologiques (Hountondji, Wiredu) : anonymat, absence de contradiction, essentialisme ; nécessité de la philosophie professionnelle.
      \item Synthèse : Dépassement par la \cle{décolonialité} : réappropriation critique du patrimoine (philosophie sagace) + rigueur académique + engagement développemental.
    \end{enumerate}
  \item \textbf{Exemples Cameroun} : Référence aux travaux de Fabien Eboussi Boulaga (critique théologique/philosophique dans \emph{La Crise du Muntu}), à la pensée Bamiléké/Bassa sur la personne et la communauté, aux GIC (Groupements d'Initiative Commune) comme éthique du développement endogène.
\end{enumerate}
\end{methode}

\begin{methode}[Commentaire de texte (ex. Hountondji, \emph{Sur la philosophie africaine}, 1983)]
\begin{enumerate}
  \item \textbf{Contextualisation} : Auteur, ouvrage, date (1983), débat sur l'ethnophilosophie (réponse à Tempels, Kagamé).
  \item \textbf{Structure} : Identifier la thèse (critique du « surrogate »), les arguments (absence de sujet pensant, fétichisme de la tradition, confusion tradition/philosophie), les concepts-clés (tradition, modernité, scientificité, endogénéité).
  \item \textbf{Analyse philosophique} : Expliquer \emph{en propres mots} le sens de « philosophie de substitution ». Montrer la tension universalisme/particularisme.
  \item \textbf{Discussion critique} : Accorder (nécessité de rigueur, subjectivité) / Nuancer (risque d'eurocentrisme inversé, valeur heuristique de l'ethnophilosophie comme source). Ouvrir sur la décolonialité actuelle (Mbembe, Ndlovu-Gatsheni).
\end{enumerate}
\end{methode}

\begin{aretenir}[Checklist avant le Bac — Philosophie Africaine]
\begin{enumerate}
  \item $\square$ Je définis \cle{ethnocentrisme} et cite 2 auteurs coloniaux (Hegel, Lucien Lévy-Bruhl) avec leur thèse fausse.
  \item $\square$ Je distingue net les \cle{4 tendances} d'Oruka (noms, thèse, limite) et je sais les mettre en tableau.
  \item $\square$ J'explique la critique de \cle{Hountondji} : « ethnophilosophie = surrogate » ; sujet absent ; tradition figée.
  \item $\square$ Je caractérise la \cle{philosophie sagace} : individu + rationalité critique + oralité (Oruka).
  \item $\square$ Je relie \cle{philosophie nationaliste} à l'indépendance politique (Nkrumah, Nyerere, Senghor) sans confondre avec l'ethnophilosophie.
  \item $\square$ Je définis la \cle{décolonialité} (épistémologie du Sud, déliaison savoir/pouvoir) et son lien avec le \cle{développement endogène}.
  \item $\square$ Je produis un \cle{plan dialectique} complet (Thèse/Antithèse/Synthèse) sur un sujet type Bac.
  \item $\square$ Je maîtrise la \cle{méthode de commentaire} : Contextualiser $\to$ Structurer $\to$ Analyser $\to$ Discuter.
  \item $\square$ J'intègre au moins \textbf{un exemple camerounais} concret (coutumes, GIC, auteur local comme Eboussi Boulaga, institution) dans mes copies.
  \item $\square$ J'évite l'anachronisme : pas de notions de Terminale (ex. « philosophie de la technique » approfondie, « bioéthique » hors cadre), je reste en 1re Tech.
  \item $\square$ Je soigne la rédaction : connecteurs logiques, définitions en début de développement, conclusion qui ouvre.
\end{enumerate}
\end{aretenir}

\findecours

\end{document}
